\documentclass[10pt]{article}
\usepackage{amssymb,enumerate,mathrsfs, amsmath, amsthm, graphicx, mathtools}
\usepackage{url,titletoc,enumitem}
\usepackage{tikz-cd}
\usepackage{bbold}
\usepackage{hyperref}
\hypersetup{
linkbordercolor = {cyan}}



%Yifan: I'm moving symbols definition to a different file so we can modify it easily.
\input{symbols.tex}


\linespread{1.18}
\textwidth = 6.5 in %6.2
\textheight = 8.5 in %9
\oddsidemargin = 0.0 in %.5
\evensidemargin = 0.0 in%.5
\topmargin = 0.0 in
\headheight = 0.0 in
\headsep = 0.3 in
\parskip = 0.05 in
\parindent = 0.3 in

\theoremstyle{plain}

\newtheorem{theorem}{theorem}[section]
\newtheorem{cor}[theorem]{Corollary}
\newtheorem{lemma}[theorem]{Lemma}
\newtheorem{Claim}[theorem]{Claim}
\newtheorem{sublemma}[theorem]{Sublemma}
\newtheorem{prop}[theorem]{Proposition}

\newtheorem{alg}[theorem]{Algorithm}
\newtheorem{conjecture}[theorem]{Conjecture}
\newtheorem{thm}[theorem]{Theorem}
\newtheorem{lem}[theorem]{Lemma}
\newtheorem{corollary}[theorem]{Corollary}
\newtheorem*{lemma*}{Lemma}
\newtheorem{Heuristic}[theorem]{Heuristic}
\newtheorem{Question}[theorem]{Question}

\theoremstyle{definition}

\newtheorem{definition}[theorem]{Definition}
\newtheorem{notation}[theorem]{Notation}
\newtheorem{obs}[theorem]{Observation}
\newtheorem{hypothesis}[theorem]{Hypothesis}
\newtheorem{wish}[theorem]{Wish}
\newtheorem{question}[theorem]{Question}
\newtheorem{example}[theorem]{Example}
\newtheorem{remark}[theorem]{Remark}
\newtheorem{rem}[theorem]{Remark}


\numberwithin{equation}{section}

\usepackage{biblatex}
\addbibresource{reference.bib}



\begin{document}

\title{A Chebotarev Density Theorem over Local Fields}
\author{Asvin G.\footnote{email id: gasvinseeker94@gmail.com}, Yifan Wei, John Yin\footnote{All three authors are affiliated with the University of Wisconsin-Madison}}
\maketitle

\begin{abstract}
    We compute the $p$-adic densities of points with a given splitting type along a (generically) finite map, analogous to the classical Chebotarev theorem over number fields and function fields. Under some mild hypotheses, we prove that these densities satisfy a functional equation in the size of the residue field. This functional equation is a direct reflection of Poincar\'e duality in \'etale cohomology. As a consequence, we prove a conjecture of Bhargava, Cremona, Fisher, and Gajovi\'c on factorization densities of p-adic polynomials. 
    
    The key tool is the notion of \emph{admissible pairs} associated to a group, which we use as an invariant of the inertia and decomposition action of a local field on the fibers of the finite map. We compute the splitting densities by M\"obius inverting certain p-adic integrals along the poset of admissible pairs. The conjecture on factorization densities follows immediately for tamely ramified primes from our general results. We reduce the complete conjecture (including the wild primes) to the existence of an explicit ``Tate-type" resolution of the ``resultant locus" over $\Spec \mathbb Z$ and complete the proof of the conjecture by constructing this resolution. 
\end{abstract}

\tableofcontents

\section{Introduction}

Let $h(z) = c_nz^n + c_{n-1}z^{n-1} \dots + c_0$ be a random polynomial having coefficients $ c_i \in \mathbb Z_p$. In this paper, we develop a general method to compute the density of polynomials with a given factorization type over $\mathbb Z_p$, thus answering a conjecture of Bhargava, Cremona, Fisher and Gajović \cite{bhargava2022density}. As we will see, our method is far more general than this special case and can be considered as a Chebotarev-type theorem for finite maps over local fields (as opposed to the classical versions which apply to varieties over finitely generated rings). For concreteness, we first explain the case of a random polynomial.

We parametrize polynomials $h(z)$ as above by $\mathbb Z_p^{n+1}$ with the Haar measure $\mu_{\mathrm{Haar}}$ normalized so that the total measure is $1$. If $h(z)$ is irreducible, define $K = \mathbb Q_p[z]/(h(z))$ and the \emph{factorization type} of $h$ by $\sigma(h) = \{f^e\}$ where $f$ is the size of the residue field of $K$ and $e$ its inertial degree. If $e=1$, we will often omit the superscript and simply write $\{f\}$ instead of $\{f^1\}$. In general, if $h(z)$ is squarefree with a factorization $h = g_1\dots g_r$ into irreducible polynomials over $\mathbb Q_p$, we define its factorization type to be the multiset $\sigma(h) = \{\sigma(g_1),\dots,\sigma(g_r)\}$. 

Fixing such a factorization type $\sigma$, let $U_\sigma(p) \subset \mathbb Z_p^{n+1}$ be the $p$-adic open subset of squarefree polynomials with factorization type $\sigma$ and $\rho(n,\tau;p) = \mu_{\mathrm{Haar}}(U_\tau(p))$. As a few sample examples, \cite{bhargava2022density} computes
\begin{center}
\begin{tabular}{ |c|c||c|c| } 
\hline
 $\rho(2,(11);p) = \frac{1}{2}$ & $p$ splits as $\pi_1\pi_2$ &$\rho(2,(2);p) = \frac{p^2-p+1}{2(p^2+p+1)}$& $p$ splits as $\pi^2$\\
 $\rho(3,(111);p) = \frac{p^4+2p^2+1}{6(p^4+p^3+p^2+p+1)}$ & $p$ splits as $\pi_1\pi_2\pi_3$ & $\rho(3,(12);p) = \frac{p^4+1}{2(p^4+p^3+p^2+p+1)}$& $p$ splits as $\pi_1\pi_2^2$\\
 \hline
\end{tabular}
\end{center}
% \begin{align*}
%     \rho(2,(11);p) &= \frac{1}{2} &\rho(2,(2);p) &= \frac{p^2-p+1}{2(p^2+p+1)}\\
%     \rho(3,(111);p) &= \frac{p^4+2p^2+1}{6(p^4+p^3+p^2+p+1)} &\rho(3,(12);p) &= \frac{p^4+1}{2(p^4+p^3+p^2+p+1)}.
% \end{align*}
Based on numerical evidence and some of their results, \cite[Conjecture 1.2]{bhargava2022density} states:
\begin{conjecture}\label{conj: intro bhargava}
The densities $\rho(n,\sigma;p)$ are rational functions in $p$ and satisfy the following symmetry:
\[\rho(n,\sigma;p^{-1}) = \rho(n,\sigma;p).\]
\end{conjecture}

In more detail, let $X: \mathbb Z_p^{n+1} \to \mathbb N$ be the random variable that sends a polynomial to the number of $\mathbb Q_p$-rational roots. Then, \cite[Theorem 1]{bhargava2022density} shows that all the moments of this random variable are rational functions in $p$ that are invariant under the transformation $p \to p^{-1}$ and they conjecture that in fact, the densities themselves satisfy both properties. Since the number of rational roots is far from determining the factorization type of a polynomial, the conjecture is much stronger than what they prove.

In the case where the splitting fields of the polynomials are tame extensions of $\mathbb{Q}_p$, \cite{del_corso_dvornicich_2000} (and independently, the last author in \cite{john}) prove that the densities $\alpha(n,\sigma;p)$ form a rational function in $p$. They compute a recurrence equation for the above densities in terms of certain other densities. The individual terms in this recurrence are rational but do not satisfy the above symmetry and therefore, we need a new idea in order to prove the functional equation (and the rationality at wild primes). More recently (and in fact simultaneously with the current paper), \cite{roy} proves the above conjecture including the symmetry for the factorization type $\{n^1\}$.

In this paper, we prove a generalization of this conjecture formulated for any finite extension of $\bbQ_p$. The case for tame primes, e.g., for $p > n$, follows from our main general results and the previous work on rationality cited above. However, we offer an independent proof that works uniformly for all primes $p$. This is the most technically involved part of our paper and involves constructing good resolutions of ``resultant varieties" in \emph{all} characteristics. 

The rationality in $p$ for factorization densities boils down to the fact that certain varieties $D/\mathbb Z$ related to the problem have point counts $p \to |D(\mathbb F_p)|$ that are polynomial functions of $p$ and is special to the case of polynomial factorizations. The functional equation ultimately reduces to the Weil Conjectures (more specifically, Poincar\'e duality) over finite fields. The key insight that led to this paper was decoupling these two parts of the conjecture. The rationality turns out to be very special but the functional equation, suitably generalized, applies to a much wider setting.

Before stating our general results, we briefly recall the classical Chebotarev theorem since our results can be considered as an analogue of that theorem for local fields. Let $f: X \to Y$ be an \'etale Galois map between varieties with constant Galois group $G$ over a finite field $\mathbb F_q$. Given a closed point $y \in Y$ with residue field $\kappa(y)$ and a geometric lift $x \in X(\overline{\mathbb F}_q)$, we have an induced map
\[\widehat{\mathbb Z} \cong \Gal(\overline{\mathbb F}_q/\kappa(y)) = \pi_1^{\text{\'et}}(y;x) \to G\]
by functoriality of the \'etale fundamental group. The image of the Frobenius element $\sigma_{\kappa(y)} \in \Gal(\overline{\mathbb F}_q/\kappa(y))$ is called the Frobenius of $y$ with respect to the base point $x$ and different choices of base points correspond to conjugation by $G$. Therefore, we can define a well defined conjugacy class $\frob{y} \subset G$. We then have the following classical theorem.

\begin{thm}[Chebotarev]
Let $c \subset G$ be a conjugacy class. Then,
\[\lim_{N \to \infty}\frac{|\{y \in Y : [\kappa(y):\mathbb F_q] \leq N, \frob{y} = c\}|}{|\{y \in Y : [\kappa(y):\mathbb F_q] \leq N\}|} = \frac{|c|}{|G|}.\]
\end{thm}


\subsection{Our Main Results}

We work in the context of an arbitrary local ring $\A$ over $\mathbb Z_p$ with fraction field $\F$ and residue field $\A/\mathfrak m_{\F} = \mathbb F_q$. Let $f: X \to Y$ be a generically finite, Galois map between smooth proper varieties defined over $\A$ with Galois group $G$, a constant group scheme over $\A$. By this, we mean that there is a Zariski open subset $\U{f} \subset Y$ such that $f: f^{-1}(\U{f}) \to \U{f}$ is an \'etale Galois map with Galois group $G$ and moreover, the natural map $\Aut(f) \to G$ is an isomorphism. 

Let $P \in \U{f}(\F)$ with $Q \in \U{f}(\overline{\F})$ a geometric lift of $P$. By functoriality of the \'etale fundamental group, we have a map
\[\Gal(\overline{\F}/\F) = \pi_1^{\text{\'et}}(P;Q) \to G.\]
We denote the image of the entire Galois group by $D_Q$, the image of the inertia group by $I_Q$ and the image of the canonical Frobenius coset by $\frob{Q}$. The Galois group of a local field is significantly more complicated than that of a finite field and thus, we need a new definition to capture the data of $(I_Q,\frob{Q})$. We call a pair $\tau = (H,gH)$ for a subgroup $H \subset G$ and an element $ g \in G$ an \emph{admissible pair} if $gH = Hg$ (in this general setting, $\tau$ will refer to more general admissible pairs as opposed to factorization types). The pair $(I_Q,\frob{Q})$ as above is an admissible pair and we define 
\[U_{f,\tau}(\F) = \{P \in \U{f}(\F) : (I_Q,\frob{Q}) = \tau \text{ for some } Q \text{ with } f(Q) = P \}.\]
For a smooth projective variety $Y$ as above, one can define a canonical measure $\mu_Y$, see \S\ref{section: canonical measure} for more details. We normalize the canonical measure $\mu_Y$ on $Y(\F)$ so that $\mu_Y(Y(\F)) = |Y(\mathbb F_q)|$. In analogy with the factorization densities conjecture, one might hope that $\rhotype{\tau}(\F) \coloneqq \mu_Y(U_{f,\tau}(\F))$ is a rational function in $q$ as we vary $\A$ over extensions of $\mathbb Z_p$. The simplest examples (such as the degree $2$ map from an elliptic curve $E \to \mathbb P^1$) very quickly disabuse us of this hope or any reasonable modification of it. Indeed, we will show that $\rhotype{\tau}(\F)$ is a linear combination of the point counts $|Z(\mathbb F_q)|$ weighted by $\eta_k(q)$ as $Z$ ranges over smooth proper varieties of dimension at most $\dim Y$ where $\eta_k(q) = (q^{k+1}-q)/(1-q^{k+1})$. The obstruction to rationality is precisely that the point counts $|Z(\mathbb F_q)|$ are generally not rational functions of $q$.

Despite this, the symmetry in Conjecture \ref{conj: intro bhargava} mentioned above continues to hold in this general context when formulated correctly. More precisely, we define the notion of a palindromic form of weight $k$ over $\A$ in Definition \ref{defn: palindromic forms} and prove that the densities are palindromic forms. For the introduction, we will only need that palindromic forms of weight $k$ are functions $\rho: \mathbb N \to \mathbb Q$ of a special form that allow a natural extension to a function $\rho: \mathbb Z \to \mathbb Q$ with the property that
\[\rho(-m) = q^{-km}\rho(m).\]

With this definition, our main theorem is as follows. Note that it is not even obvious a-priori that the densities $\rhotype{\tau}(\L)$ lie in $\mathbb Q$.

\begin{thm}[Theorem \ref{thm: Main theorem, density fixing inertia and frobenius coset}]
Let $f: X \to Y/\Spec\A$ be a generically finite, Galois map with Galois group $G$ between smooth, projective varieties and $\tau = (H,gH)$ an admissible pair as above. We denote its inverse by $\tau^{-1} = (H,g^{-1}H)$. Suppose that, for every $\tau' = (H',g'H')$ with $H' \subset H$ and $g' \in gH$, there exists a $g'$-equivariant resolution $\pi_{H'}: \widetilde{X_{H'}} \to X/H'$ such that the ramified locus in $\widetilde{X_{H'}}$ for the map $\widetilde{X_{H'}} \to Y$ is a simple normal crossings divisor over an unramified extension of $\A$. 

Then, for any finite extension of local rings $\A \subset \B$ with corresponding residue field extensions $\mathbb F_q \subset \mathbb F_{q^m}$, the densities $\rhotype{\tau}(\L) \in \mathbb Q$ and the function
\[ m \to \rhotype{\tau}(\L) + \rhotype{\tau^{-1}}(\L)\]
depends on $\L$ only through the size of the residue field $\mathbb F_{q^m}$ and is a palindromic form (in the variable $m$) of weight $\dim_{\F} Y_{\F}$.
\end{thm}

One can consider this as a refinement of the classical Chebotarev density theorem over finite fields as $q \to \infty$ (Corollary \ref{cor: q to infty limiting behaviour}). In this limit, the ramified densities tend to $0$ and the unramified densities tend to the classical limits. Moreover, this general theorem will imply Conjecture \ref{conj: intro bhargava} for all but finitely many primes as we explain later.

Simple examples \footnote{such as the map $\mathbb P^1 \to \mathbb P^1; t \to pt^2$} show that it's necessary to assume the existence of such a resolution. If $f: X \to Y$ is generically finite with Galois group $G$ over (an open subset) of the ring of integers of a number field, then the resolution hypothesis is satisfied for all but finitely many completions of the number ring since we can equivariantly resolve singularities over characteristic $0$ (for instance, see \cite{abramovich1996equivariant},\cite{encinas1998good} and \cite{bierstone1997canonical} among many others) and \emph{spread out} to all but finitely many primes as we show in Theorem \ref{thm: main theorem for almost all primes}.

Moreover, the densities of the $\Utype{\tau}(\L)$ individually are not palindromic forms of any weight. The key here is Lemma \ref{lemma: twisted point counts have the right weight} and Remark \ref{rmk: individual densities are not palindromic} shows that $\rhotype{\tau}$ is not palindromic by itself in general. Similarly, the densities where we require that the splitting field is a fixed extension are also not palindromic (for example, if we require the splitting field of $f$ to be $\bbQ_p(sqrt p)$, and not $\bbQ_p(\sqrt{dp}))$, where $d$ is a non-square mod $p$). As Caruso explains immediately following \cite[Theorem C]{caruso2022zeroes}, the expected number of roots in a fixed ramified quadratic extension (such as $\bbQ_p(sqrt p)$ of a random $p$-adic degree $n$ polynomial (for $n \geq 3$) is not invariant under the $p \to p^{-1}$ transformation but this symmetry is regained when we consider instead the expected number of roots in any ramified quadratic extension. 

In the context of a random polynomial over $\mathbb Z_p$, the relevant map is
\[f: X = (\mathbb P^1)^n \to (\mathbb P^1)^n/S_n = \mathbb P^n = Y; ([t_1,s_1],\dots,[t_n,s_n]) \to [a_0, \dots, a_n]\]
where $(s_1z-t_1)\dots(s_nz-t_n) = \sum a_i z^i$. For any local field $K/\mathbb Q_p$ with residue field $\mathbb F_q$, one can define an analogous notion of factorization density $\rho(n,\sigma;q)$. A priori, this might depend on $K$ but we show that it is purely a function of $q$. We prove that it is a rational function satisfying the following functional equation.

\begin{thm}[Theorem \ref{thm: rationality of wild-type densities}]
The factorization densities $\rho(n,\sigma;q)$ form a rational function in $q$ and satisfy the symmetry
\[\rho(n,\sigma;q^{-1}) = \rho(n,\sigma;q).\]
\end{thm}


\begin{remark}[A motivic analogue]\label{rmk: motivic duality}
The methods of our proof can be adapted to the setting of motivic integration with the following differences:
\begin{itemize}
    \item We have a satisfactory theory of resolutions of singularities and the main theorems are therefore hypothesis free.
    \item The absolute Galois group of $\mathbb C[\![t]\!]$ is much simpler than that of $\mathbb Q_p$ which obviates the need for the notion of \emph{admissible pairs} and the associated induction of their poset. Instead, we can simply use the poset of cyclic subgroups (under inclusion) and an induction along this poset instead.
    \item One important complication however is that motivic integration needs to be modified to take values in an \emph{equivariant} Grothendieck ring of motives in order to define a change of variables formula for generically finite maps as opposed to the standard theory dealing with birational maps. This modification is discussed in \cite[\S2.1]{carvajal2021behavior}, for example.
\end{itemize}
Given these modifications, exactly the analogous theorems of this paper should hold althought we do not pursue this direction here. In this context, the transformation $m \to -m$ should be replaced by Bittner duality \cite[Corollary 3.4]{bittner2004universal}. This duality is a manifestation of the duality in the conjectured abelian category of pure motives and is closely related to Poincar\'e duality. As such, our transformation $m \to -m$ can be seen as a $p$-adic shadow of this motivic duality. 
\end{remark}


\subsection{An overview of the proof}


The first observation in proving Theorem \ref{thm: Main theorem, density fixing inertia and frobenius coset} is that the \emph{split points}, i.e., the points corresponding to the admissible pair $(\mathrm{id},\mathrm{id})$, correspond exactly to the image of the points $X(\F)$ under $f$. This density can then be computed by an integral over the $p$-adic analytic set $X(\F)$ of an appropriate function of the form
\[\rhotype{(\mathrm{id},\mathrm{id})}(\F) = \eta_f(\F) = \int_{X(\F)}|\Jac(f)|_{\F}d\mu_{X}.\]
We have traded the measure of a possibly complicated space for the integral of a possibly complicated function over a nice space. However, we show that the function is not too bad if there exists a nice resolution of singularities and we explicitly compute the above integral in Theorem \ref{thm: twist symmetry under pullback} which proves that the integrals are palindromic forms.

In the simplest case when the ramified locus $Z_f \subset X$ can be resolved by a simple normal crossings divisor, we have the following.

\begin{corollary}[Corollary \ref{cor: pullback formula in sncd case}]
Let $f: X \to Y$ be a generically finite, Galois map and suppose that there exists a resolution of singularities $\pi: \widetilde{X} \to X$ over $\A$ with  ramification locus $\pi^{-1}(Z_f) = \bigcup_{i=1}^r D_i \subset \tilde{X}$ a simple normal crossings divisor over $\A$. Then
\[\eta_f(\F) = \sum_{J \subset \{1,\dots,r\}}\left|\left(\bigcap_{j\in J}D_j\right)(\mathbb F_{q})\right|\prod_{j\in J}\frac{q - q^{e_j+1}}{q^{e_j+1}-1}\]
for some positive integers $e_j$.
\end{corollary}

Similar formulas as in Corollary \ref{cor: pullback formula in sncd case} are well known in slightly different forms in various \emph{geometric} situations where $Z_f$ is a simple normal crossings divisor (\cite[Proposition 4.11]{yasuda2017wild}, \cite[Theorem 2.3.2]{denef1992caracteristiques}, \cite[Theorem 2.8.1]{loeser2009seattle}) and of most relevance to us \cite[Theorem 3]{denef1991functional}. Indeed, the last reference proves a similar functional equation to the ``remarkable symmetry" using similar methods to this paper. There are two major new complications that we need to deal with.

First, \cite{denef1991functional} only deals with the case of an irreducible hypersurface in projective space. We need to consider the more general case of a finite map between varieties which follows with few modifications the proof of the hypersurface case. More importantly, they assume that the irreducible divisors appearing in the complement $Z_f$ are geometrically irreducible. We \emph{need} to allow for \emph{reducible} divisors since our inductive step deals with \emph{unramified Galois twists} of our finite map. Even if the $Z_f$ were initially composed of geometrically irreducible divisors, this will almost certainly not be true after twisting. This significantly complicates the local computations and involves new ideas. The key here is Lemma \ref{lemma: local integrals}.

To compute the densities corresponding to an arbitrary admissible pair $\tau = (H,gH)$ as above, we compute the densities of ${^\tau f}({^\tau X}(\F)) \subset Y(\F)$ in terms of integrals similar to the ones discussed above. Here, ${^\tau X} \coloneqq  {^{g^{-1}} X}/H$ is the quotient of an unramified $g$-twist of $X$ by $H$  while $ {^\tau f}: {^\tau X} \to Y$ is the quotient map. We show that the sum of the densities of ${^\tau f}({^\tau X}(\L))$ and $ {^{\tau^{-1}} f}({{^{\tau^{-1}}} X}(\L))$ are palindromic forms of weight $\dim Y$ as we vary over $\A \subset \B$ in Theorem \ref{thm: twist symmetry under pullback}. This is not an immediate consequence of our integral computation because we twist \emph{after} base changing. The crucial input is Lemma \ref{lemma: twisted point counts have the right weight}.

We express the densities of these images as a linear combination of the densities $\rhotype{\tau'}$ for $\tau' \leq \tau$. This system of linear equations can be considered as an identity in the incidence algebra of the poset of admissible types which lets us express the splitting type densities in terms of the densities of the images discussed above. Said more simply, the system of linear equations is upper triangular and invertible with respect to an appropriate basis.

Finally, we consider the case of polynomial factorization densities. In this case, the relevant map is the $S_n$ quotient $(\mathbb P^1)^n \to \mathbb P^n$. In this case, one might naively hope that all the splitting densities associated to each admissible pair for $S_n$ satisfy the general conjecture in \cite{bhargava2022density} but this hope turns out to be false. Nevertheless, we relate the polynomial factorization densities directly to integrals over a certain family of "nice" quotients corresponding to "polynomial multiplication" maps
\[\mathbb P^{e_1}\times\dots\times \mathbb P^{e_r} \to \mathbb P^{e_1+\dots+e_r}\]
where we identify points of projective space with univariate polynomials of the appropriate degree so that the above map corresponds to
\[(P_1(z),\dots,P_r(z)) \to P_1(z)\dots P_r(z).\]
The ramification locus for this corresponds exactly to sets of polynomials $(P_1(z),\dots,P_r(z))$ where at least two of the polynomials of a common root. Finally, in \S\ref{sec: resolving resultant locus}, we describe an explicit resolution of this resultant locus by an Artin stack. The resolution is explicit enough that we can verify the rationality by hand, thus completing the proof of rationality for all primes, including the wild ones.


\begin{remark}
In \cite[Conjecture 9.1]{john}, the last author conjectured a further functional equation for a generating function $\rho(\sigma,e,f;p,t)$, which specializes to $\rho(n,\sigma;p)$ after setting $t=p^{-e/2}$. The methods of this paper generalize without any difficulty to prove this conjecture by considering \[\eta_f(\F,t) = \int_{X(\F)}|\Jac(f)|_{\F}t^{\nu_K (\Jac(f))}d\mu_{X}\] instead of \[\eta_f(\F) = \int_{X(\F)}|\Jac(f)|_{\F}d\mu_{X}.\]
\end{remark}

\subsection{Relation to previous work}

Our work draws from multiple strands of research as described below.

\subsubsection*{Densities of random polynomial with fixed splitting data} 

The question of splitting field densities in the archimedean case goes back at least to work of Bloch and P\"olya \cite{bloch1932roots} who proved asymptotic bounds on the number of expected roots of polynomials with coefficients in $\{-1,0,1\}$, followed by a series of papers of which a highlight is the work of Kac \cite{kac1943average}, which provides an \emph{exact formula} for the average number of roots of a random polynomial with Gaussian coefficients. For a survey and related results, \cite{dembo2002random} and \cite{nguyen2021random}. The $p$-adic version was first considered by Evans \cite{evans2006expected}, followed by a series of authors (\cite{buhler2006probability},\cite{kulkarni2021p}) all of whom were concerned with the expected number of roots of a random $p$-adic polynomial for which the methods of Kac can be extended. Beyond the mean, the aforementioned paper \cite{bhargava2022density} was among the first to consider the density of $p$-adic polynomials with a given number of $\mathbb Q_p$ roots and they proved that these densities are rational in $p$ and are symmetric under $p \to p^{-1}$. Another recent interesting paper is by Caruso \cite{caruso2022zeroes}, which is interested in the expected number of roots in an extension $\mathbb Q_p \subset F$ of a random $p$-adic polynomial.

\subsubsection*{Motivic and $p$-adic integration}

Our proof owes the most debt to the literature on $p$-adic and motivic integration. The story begins with Batyrev's paper \cite{batyrev1997birational} proving that birational Calabi-Yau varieties have the same Betti numbers using ideas from $p$-adic integration. Following this, Kontsevich in a 1995 lecture in Orsay introduced the theory of Motivic integration to prove that birational Calabi-Yau varieties have the same Hodge numbers. In both these contexts, the key tool is a change of variables formula relating ($p$-adic or motivic) integrals on the $Y_i$ to corresponding integrals on a common resolution, where the integral is much simpler to compute. There has been much work since then in this area of which we would like to highlight the literature on McKay correspondences (\cite{reid2002correspondance} for a survey).

Given a finite group $G$ with an action on a variety $V$, this correspondence seeks to relate invariants of $V/G$ to invariants of $G$. The invariants involved are ``stringy integrals" very similar to the ones appearing in our work (as in Corollary \ref{cor: pullback formula in sncd case}). Indeed, our general formula involves computing these stringy invariants for quotients $X/H$ for $H \subset G$ a subgroup with respect to a \emph{ramification divisor} on $X$. \cite[Theorem 1.1]{yasuda2017wild} proves a result relating these stringy invariants on $X/H$ to certain other stringy invariants on $X$. It is quite plausible that the methods of our final section could be replaced by using the McKay correspondence and thus generalized and made more explicit.

Another landmark result with thematic similarities to our paper is Serre's mass formula \cite{serre1978formule} and especially the first proof in that paper using $p$-adic integration. In fact, there have been many recent results along similar themes of which \cite{caruso2022zeroes}, \cite{bhargava2007mass} and \cite{yasuda2017wild} are especially relevant to our work. Some other recent highlights are \cite{groechenig2020geometric}, \cite{groechenig2020mirror}.

Finally, we would like to mention the work of \cite{denef_loeser_2000}. They associate to every formula $\varphi$ in the first order language of rings with coefficients in a field of characteristic $0$ a virtual motive $\chi_c(\varphi)$ in the Grothendieck ring of an appropriate category of motives. The formula can be interpreted as defining certain subsets of varieties as in the following example.

For $\mathbb P^1 \to \mathbb P^1; t \to t^n$, the subset of points in the image of this map corresponds to the formula $\varphi(x) = ``\exists y, x = y^n"$. More generally, for a finite Galois map $f: X \to Y$, the subset of points in $Y$ with prescribed decomposition group correspond to some such formula.

In \cite[Corollary 7.1.2]{denef_loeser_2000}. they prove that $\chi_c(\varphi)$ belong to the subring of motives obtained by localizing at certain polynomials in the Lefschetz motive. Under \'etale realization, $\chi_c(\varphi)$ specializes to the canonical measure of the corresponding subsets. To emphasize the strength of this theorem, it is not even clear a priori that these densities should be rational!

In relation to this work, our theorem provides a different perspective to the uniformity phenomenon. While Denef and Loeser use quantifier elimination and motivic integration to prove the uniformity, our proof uses more geometric methods. As such, we obtain an explicit formula that can be used to prove the required functional equation (and rationality for factorization densities).
     



\subsection{Organization of the paper}

In \S \ref{section: preliminaries}, we review some well known theorems that will be useful for us and standardize notations and normalizations. The most important novelty in this section is the notion of a \emph{palindromic form of weight k} (Definition \ref{defn: palindromic forms}). The other novelty is in \S \ref{section: admissible pairs}, where we define the notion of an admissible pair related to a group and a corresponding partial order on it. In \S \ref{section: on the palindromicity}, we prove our main Theorem \ref{thm: Main theorem, density fixing inertia and frobenius coset} and derive some corollaries. This section contains the core arguments of our paper. Finally, in \S \ref{section: factorization types density}, we adapt our general method to prove the conjecture of Bhargava, Cremona, Fisher and Gajović. The main work is involved in showing that we can construct a nice, \emph{Tate}-type resolution of singularities in this case.


\textit{Acknowledgements.} We would like to thank Jordan Ellenberg, Daniel Litt and Andrew O'Desky for general helpful discussions. Moreover, we would like to thank Dan Abramovich, Colin Crowley, Aise Johan de Jong, Connor Simpson, Jason Starr, and Botong Wang for helping with the resolution of the resultant variety in the final section.

\section{Preliminaries}\label{section: preliminaries}

\subsection{Geometry over $p$-adic rings}

Unless mentioned otherwise, we fix a prime $p$ throughout the paper and let $\A$ be a local ring over $\bbZ_p$ with maximal ideal $\mathfrak m_\F$ generated by a uniformizer $\pi_\F$, fraction field $\F$ and residue field $\A/\mathfrak m_\F= \RF$ of size $q$, a power of $p$ (and use similar notation for extensions $\F \subset \L$). We will think of $\A$ as fixed and vary over extensions $\A \subset \B$ with residue field extension of degree $m = [\RL:\RF]$. We define $|\cdot |_{\L}$ to be the non-archimedean metric on $\L$ (so that $|\pi_\L|_\L = q^{-m}$) and $\frob{\L}$ to be the Frobenius acting on the unramified closure of $\L$. We denote the Frobenius element in $\Gal(\overline{\mathbb F}_q/\mathbb F_q)$ by $\frob{q}$.

We fix $\SS= \Spec\B$ and define a variety $X \to \SS$ to be a finite type, geometrically reduced scheme that is flat over $\SS$. By $\dim X$, we mean the relative dimension over $\SS$. Given an extension of schemes $T \to \SS$, we denote the base change by $X_T$ and similarly for morphisms between varieties. For a ring $R$, we let $X(R)$ denote the $\Spec R$ valued points of $X$. 

Given a scheme $X\to\Spec \B$ and a point $x \in X(\B)$, we define $\cO_{X,x}$ to be the localization at the maximal ideal $\mathfrak m_{X,x}$ corresponding to the reduction $\overline{x} \in X(\B/\mathfrak m_\B)$ and $\widehat{\cO}_{X,x}$ to be its completion at its maximal ideal. When $X$ is smooth, a complete set of local co-ordinates $t_1,\dots,t_{\dim X}$ at $x$ determines an isomorphism $\widehat\cO_{X,x} \cong \B[\![t_1,\dots,t_{\dim X}]\!]$ and consequently, 
\[\widehat{\cO}_{X,x}(\B) = \widehat{\cO}_{X,x}(\L) = \mathfrak m_B^{\dim X}.\]
Here, $\widehat{\cO}_{X,x}(\B)$ is taken to be the space of \emph{continuous} ring homomorphisms to $\B$.




\subsection{$p$-adic integration.}\label{section: canonical measure}

We will use techniques of $p$-adic integration throughout this paper. In this section, we give a brief introduction and define our normalizations. We fix an extension $K \subset L$ with residue field $\mathbb F_{q^m}$ and $U \to \Spec \B$ a smooth variety of dimension $d$ in this section. The set $U(\L)$ carries the natural structure of a \emph{$\L$-analytic manifold} \cite[\S2.4]{igusa2000introduction}.

\begin{remark}
    In fact, we need not restrict to $U$ being a smooth variety, the theory works without modification for $U$ an algebraic space and we will need this generality in the final section. However, for simplicity we stick to the case of $U$ being a scheme in this introduction and refer to \cite[\S 2.1]{groechenig2020geometric} for the details in the case of $U$ being an algebraic space.
\end{remark}

\begin{definition}[$p$-adic integration with respect to a form]
Let $\omega \in \Omega_U^d$ be a top form. For $V \subset U(\L)$ a compact subset, one can define the integral
\[\mu_{\omega}(V) \coloneqq \int_{V}|\omega| \in \mathbb R_{\geq 0}.\]
If $\omega = f(x)dx_1\wedge\dots\wedge dx_n$ for a $\L$ analytic function $f$ and local co-ordinates $x_1,\dots,x_n$ on $V$, then
\[\int_{V}|\omega| = \int_{V'}|f(x)|d\mu_{\L^d} \in \mathbb R_{\geq 0}\]
where we identify $V$ with an open subset $V'$ of $\L^d$ using the $x_i$ and $\mu_{\L^d}$ is the canonical Haar measure on $\L^d$ normalized so that $\mu_{\L^d}(\B^d) = q^{md}$. Note that this normalization is slightly non-standard.

\end{definition}

\begin{definition}[Measures from compactifications]\label{defn: measures from scheme compactifications}

Given an open immersion $U \subset X$ over $\Spec\B$ with $X$ a smooth and proper variety, we can assign a canonical measure $\mu_X$ on the set $U(\L)$ \cite[\S 4.1]{yasuda2017wild}. We cover $X$ by Zariski opens $U_i$ such that $U_i(\B)$ covers $X(\B) = X(\L)$ and moreover trivializes $\Omega_X^d|_{U_i}$. Let $\omega_i \in \Omega_X^d(U_i)$ be gauge forms, i.e., nowhere vanishing sections. While these $\omega_i$ might not glue together on intersections, the transition functions are are nowhere vanishing on $(U_i\cap U_j)(\B)$ and hence have $p$-adic norm $1$. Therefore, the measures $\mu_{\omega_i}$ on the $p$-adic analytic sets $U_i(\B)$ glue together to give a unique measure $\mu_X$.
\end{definition}

We note that our normalization is such that $\mu_X(U(\L)) = |X(\mathbb F_{q^m})|$ which is \emph{not} the standard one found in the literature. 

The following is easy but very useful.

\begin{thm}\cite[Proposition 2.1]{groechenig2020geometric}\label{thm: change of variables}
Suppose we have a top form $\omega \in \Omega_U^d$ and $V \subset U(\L)$ a compact subset so that $\mu_\omega(V)$ is well defined. Then
\begin{enumerate}
    \item If $V = Z(\L)$ for $Z$ a Zariski-closed subscheme, $\mu_\omega(V) = 0$.
    \item Suppose we have an \'etale map $f: \widetilde U \to \widetilde U/G = U$ where $G$ is an \'etale group scheme over $\B$ acting freely on $\widetilde U$. Then
    \[\frac{1}{|G(\L)|}\int_{f^{-1}(V)}|f^*\omega| = \int_{V}|\omega|.\]
\end{enumerate}
\end{thm}
\begin{proof}
The first part is well known while the second follows from the observation that the map $f$ induces a covering space map of $\L$-analytic manifolds
\[f: \widetilde{U}(\L) \to U(\L)\]
of degree $|G(\L)|$. The theorem is then immediate.
\end{proof}
\subsection{Palindromic forms}

Let $K_0(\mathrm{var})$ be the Grothendieck ring of varieties over $\mathbb F_q$. We consider the extension given by 
\[\widetilde{K_0(\mathrm{var})} = \left(K_0(\mathrm{var})\otimes\mathbb Q\right)\left[\frac{\mathbb L^k - \mathbb L}{\mathbb L^{k}-1}\right]_{k\geq 2}\]
where $\mathbb L = [\mathbb A^1]$ is the Lefschetz motive. This ring has implicitly and explicitly appeared before in the literature on motivic integration (\cite{denef1991functional}, \cite[\S 7]{denef_loeser_2000}).

\begin{definition}\label{defn: palindromic forms}

We will be concerned with functions $\rho: \mathbb N \to \mathbb Q$ which are of the form
\[m \to |X(\mathbb F_{q^m})|\]
for $X \in \widetilde{K_0(\mathrm{var})}$. Explicitly, such functions are $\bbQ$-linear combinations of products of functions of the form
\[m \to \frac{q^{mk}-q^m}{q^{mk}-1} \text{ and } m \to \sum_{i}\alpha_{i}^m\]
where the $\alpha_i$ are Weil numbers (relative to $q$).\footnote{A Weil number relative to $q$ is an algebraic integer $\alpha$ such that there exists a $w \in \mathbb N$ so that under any complex embedding $\sigma: \overline{\mathbb Q} \to \mathbb C, |\sigma(\alpha)| = q^{w/2}$.} 
Both types of functions can be extended to have domain $\mathbb Z$ by allowing negative powers in the exponents. We say that such a function has weight $k$ if the natural extension of the function to $\mathbb Z$ given by the same formula satisfies
\[\rho(-m) = q^{-mk}\rho(m)\]
and we call such functions \emph{palindromic forms of weight $k$}.
\end{definition}

\begin{example}\label{ex: roots of untiy are palindromic forms}
Consider the class of the scheme $\cP_k = \Spec \mathbb F_{q^k} \in \widetilde{K_0(\mathrm{var})}$. The point counts $\cP_k(\mathbb F_{q^m})$ correspond to the function
\[m \to \gcd(m,k) = \sum_{a=1}^k\zeta_{k}^{am}.\]
This is a palindromic form of weight $0$.
\end{example}

\begin{remark}\label{rmk: palindromic forms closed under sums and products}
It is immediate that the sum of two palindromic forms of weight $k$ has weight $k$ while the product of palindromic forms of weights $k_1,\dots,k_r$ has weight $k_1+\dots+k_r$.
\end{remark}

\begin{lemma}\label{lemma: point counts have right weight}
Let $X/\mathbb F_q$ be a smooth, proper variety. The function $m \to \rho_X(m) = |X(\mathbb F_{q^m})|$ is a palindromic form of weight $\dim X$.
\end{lemma}
\begin{proof}
By the Grothendieck-Lefschetz trace formula, one has Weil numbers $\alpha_{ij}$ for $0 \leq i \leq 2\dim X$ and $1 \leq j \leq \dim H^i_{\text{\'et}}(X\times_{\mathbb F_q}\overline{\mathbb F}_q,\mathbb Q_\ell)$ such that
\[\rho_X(m) = \sum_{i=0}^{2\dim X}\sum_{j}\alpha_{ij}^m.\]
Since $X$ is smooth and proper, Poincar\'e duality guarantees us that for each $i,j$, there exists a $j'$ (with $j \to j'$ a permutation) such that $\alpha_{2\dim X-i,j'} = q^{\dim X}\alpha_{ij}^{-1}$. Therefore
\[\rho_X(-m) = \sum_{i=0}^{2\dim X}\sum_{j}\alpha_{ij}^{-m} = q^{-m\dim X}\sum_{i=0}^{2\dim X}\sum_{j}\alpha_{2\dim X-i,j'}^m = q^{-m\dim X}\rho_X(m)\]
as required.
\end{proof}

\begin{lemma}\label{lemma: twisted point counts have the right weight}
Let $X/\mathbb F_q$ be a smooth proper variety of dimension $k$ and let $\ell$ be a prime coprime to $q$. For $g \in \mathrm{Aut}_{\mathbb F_q}(X)$ of finite order $d$, the function
\[m \to \rho_{X,g}(m) \coloneqq \sum_{i=0}^{2\dim X}(-1)^i\tr\left((g+g^{-1})\frob{q}^m|H^i_{\text{\'et}}(X\times_{\mathbb F_q}\overline{\mathbb F}_q,\mathbb Q_\ell)\right)\]
is a palindromic form of weight $\dim X$.
\end{lemma}
\begin{proof}
Since $g$ is defined over $\mathbb F_q$, it commutes with the action of $\frob{q}$ on \'etale cohomology. Since it is of finite order, it acts semisimply and we can find a common set of eigenvectors $v_{ij} \in H^i_{\text{\'et}}(X\times_{\mathbb F_q}\overline{\mathbb F}_q,\overline{\mathbb Q}_\ell)$ with
\[\frob{q}(v_{ij}) = \alpha_{ij}v_{ij}, g(v_{ij}) = \lambda_{ij}v_{ij}\]
with the $\alpha_{ij}$ as in the previous lemma and the $\lambda_{ij}$ roots of unity of order dividing $d$. Also as in the previous lemma, Poincar\'e duality guarantees us that for each $i,j$, there exists a $j'$ with a $G$-equivariant pairing
\[\overline{\mathbb Q}v_{ij} \otimes \overline{\mathbb Q} v_{ij'} \to H^{2\dim X}(X\times_{\mathbb F_q}\overline{\mathbb F}_q,\overline{\mathbb Q}_\ell) = \overline{\mathbb Q}(-\dim X)\]
so that $\alpha_{2\dim X-i,j'} = q^{\dim X}\alpha_{ij}^{-1}$ and $\lambda_{2\dim X-i,j'} = \lambda_{ij}^{-1}$ since $G$ acts trivially on the top degree cohomology\footnote{This can be seen, for instance, through the cycle class map since the class of a point generates the top degree cohomology and any two points are algebraically equivalent, therefore the $g$-action on a point is trivial.}. Therefore,
\begin{align*}
    \rho_{X,g}(m) &= \sum_{i=0}^{2\dim X}\sum_{j}\alpha_{ij}^m\left(\lambda_{ij} + \lambda_{ij}^{-1}\right)\\
    \implies \rho_{X,g}(-m) &= \sum_{i=0}^{2\dim X}\sum_{j}\alpha_{ij}^{-m}\left(\lambda_{ij} + \lambda_{ij}^{-1}\right)\\
    &=\sum_{i=0}^{2\dim X}\sum_{j}q^{-m\dim X}\alpha_{2\dim X-i,j'}^{m}\left(\lambda_{2\dim X-i,j'}^{-1}+\lambda_{2\dim X-i,j'}\right)\\ &= q^{-m\dim X}\rho_{X, g}(m).
\end{align*}
\end{proof}

\begin{remark}\label{rmk: individual densities are not palindromic}
We note that the function
\[m \to \sum_{i=0}^{2\dim X}(-1)^i\tr\left(g\frob{q}^m|H^i_{\text{\'et}}(X\times_{\mathbb F_q}\overline{\mathbb F}_q,\mathbb Q_\ell)\right)\]
is, perhaps surprisingly, not a palindromic form in general. Let $E$ be the elliptic curve defined by the equation $y^2 = x^3 + x$ over $\mathbb Z_p$ with $p \equiv 1 \pmod{4}$ and a fixed  $i = \sqrt{-1} \in \mathbb Z_p$. We have an automorphism $g: E \to E$ defined over $\mathbb Z_p$ given by $g(x,y) = (-x,iy)$. Since $g$ has eigenvalues $\pm i$ on $H^1_{\text{\'et}}(\overline{E},\mathbb Q_\ell)$, there exist $a,b \in \mathbb Z$ such that $a^2 + b^2 = p$ and\footnote{For instance, if $p = 5$ with $i \equiv 2 \pmod{5}$, we have $a = 1, b = 2$.}
\begin{align*}
    \nu(m) \coloneqq \sum_{i=0}^{2\dim X}(-1)^i&\tr\left(g\frob{p}^m|H^i_{\text{\'et}}(X\times_{\mathbb F_p}\overline{\mathbb F}_p,\mathbb Q_\ell)\right) = 1 + (i(a+ib)^m -i(a-ib)^m) + p^m.
\end{align*}
Therefore,
\begin{align*}
\nu(-m) =& 1 + (i(a+ib)^{-m} -i(a-ib)^{-m}) + p^{-m}\\
= &p^{-m}\left(p^m + (i(a-ib)^{m} -i(a+ib)^{m}) + 1 \right) \\\neq& p^{-m}\nu(m)
\end{align*}
which proves that $\nu$ is not a palindromic form.
\end{remark}


\begin{example}\label{defn: rho, eta}
For any extension $\A\subset \B$ where $\B$ has residue field $q^m$, we define
\[\rhol_k(m;q) \coloneqq \int_{\mathfrak m_\L}|t^k|_\L dt = \frac{q^m-1}{q^{m(k+1)}-1}\]
and $\eta_k(m;q) = \rhol_k(m;q) - 1$. The function $m \to \eta_k(m;q)$ is a palindromic form of weight $1$:
\[\eta_k(-m;q) = \frac{q^{-m} - q^{-m(k+1)}}{q^{-m(k+1)}-1} = q^{-m}\frac{q^{m(k+1)}-q^m}{1-q^{m(k+1)}} = q^{-m}\eta_k(m;q).\]
Note that $\rhol_k,\eta_k$ depend on $\B$ only through the size of its residue field $q^m$. Indeed, $\eta_k(m;q)$ is a function purely of $q^m$ but we write it this way to emphasize that we will think of $q$ as fixed and $m$ as varying.

\begin{definition}\label{defn: palindromic forms as a function of rings}
We say that a function $\rho: \{\text{extensions } \F \subset \L\} \to \mathbb Q$ is a palindromic form of weight $k$ if
\begin{enumerate}
    \item The function depends on $\L$ only through the degree of residue field extensions $m$.
    \item The function $m \to \rho(\L)$ is a palindromic form of weight $k$.
\end{enumerate}
\end{definition}
\end{example}

\subsection{The poset of admissible pairs}\label{section: admissible pairs}

We discuss here a purely group theoretic notion that will become crucial in the proof of Theorem \ref{thm: Main theorem, density fixing inertia and frobenius coset}. Let $G$ be a finite group.

We say that a tuple $\tau = (H_\tau,g_\tau H_\tau)$ consisting of a subgroup and a left coset is an \emph{admissible pair} if $g_\tau$ normalizes $H_\tau$, i.e., $g_\tau H_\tau = H_\tau g_\tau$. The data of an admissible pair $\tau$ is equivalent to the choice of a pair of subgroups $H_{\tau} \subset H_{1,\tau} \subset G$ such that $H_\tau$ is normal in $H_{1,\tau}$ and $C_\tau \coloneqq H_{1,\tau}/H_{\tau}$ is cyclic along with the choice of a distinguished generator $g_\tau$ for this cyclic group. Here, $H_{1,\tau}$ corresponds to the subgroup generated by $g_{\tau}$ and $H_{\tau}$.

We say that $\tau' \leq \tau$ for two admissible pairs if $H_{\tau'} \subset H_{\tau}$ and $g_{\tau'} \in g_{\tau}H_{\tau}$. This is easily checked to be a partial order. There is one maximal element $(G,G)$ while the minimal elements correspond exactly to $(e,g)$ for $e$ the identity subgroup and $g \in G$ any element.

For $\lambda \in G$, we define the conjugate of $\tau = (H,gH)$ by $\lambda$ by
\[\lambda \tau \lambda^{-1} \coloneqq (\lambda H \lambda^{-1}, \lambda gH \lambda^{-1}).\]
We also define the inverse of an admissible pair as
\[(H,gH)^{-1} \coloneqq (H,g^{-1}H).\]
Both conjugation and inversion are easily seen to be automorphisms of the poset of admissible pairs. 

% Finally, for an admissible pair $\tau$, we define the centralizer 
% \[C_\tau = \{h \in H_\tau : g_{\tau}h = hg_{\tau}\}\]
% with size $c_\tau = |C_\tau|$.

We will use the language of an incidence algebra on a poset. We briefly recall this notion for the convenience of the reader. Given a poset $\cP$, an interval $[a,b]$ in it is the set of elements $c$ such that $a \leq c \leq b$ (where we tacitly assume that $a\leq b$). Its incidence algebra $\cI_{\cP}$ (over any ring $R$) is the ring of functions from the set of intervals to $R$. Addition in $\cI_{\cP}$ is pointwise while multiplication is defined by convolution:
\[\text{For all } \alpha,\beta \in \cI_{\cP}, \hspace{10 mm} (\alpha\ast\beta)([a,b]) = \sum_{a \leq c \leq b}\alpha([a,c])\beta([c,d]).\]
This multiplication need not be commutative. An element $\alpha \in \cI$ is invertible precisely when $\alpha([a,a]) \in R^\times$ for all $a \in \cP$. We will suppress the ring $R$ from the notation, it will usually be taken to be $\mathbb Q$ or $\mathbb R$.

\begin{definition}\label{defn: alpha, beta}
In particular, we will make use of the following elements in $\cI_{\cP}$ for $\cP$ the poset of admissible pairs:
\[\alpha([\tau',\tau]) \coloneqq |\{\lambda \in H_\tau\backslash G : \lambda \tau' \lambda^{-1} \leq \tau\}|,\]
\[\beta([\tau',\tau]) \coloneqq |\{\tau'' \in \cP : \tau'' = \lambda \tau' \lambda^{-1} \leq \tau \text{ for some } \lambda \in G.\}|.\]
and $\gamma([\tau',\tau]) \coloneqq \alpha([\tau',\tau])/\beta([\tau',\tau])$. If $\tau' \leq \tau$, then $\alpha([\tau',\tau]),\beta([\tau',\tau])$ are both at least $1$ since $\tau'$ is conjugate to itself. Therefore, $\gamma([\tau',\tau]) > 0$ when $\tau' \leq \tau$.

Note that the condition in the definition of $\alpha$ indeed only depends on the coset $H_\tau\lambda$. It is also clear that $\alpha$ and $\beta$ are invariant under conjugating the poset by an element of $G$.
\end{definition}


\subsection{Generically finite, Galois maps}

Let $X,Y$ be two geometrically connected, smooth and proper varieties relative to $\Spec \A$.

\begin{definition}\label{defn: generically fintie}[generically finite]
A map $f: X \to Y$ is said to be generically finite if there is some open $U \subset Y$ such that $f: f^{-1}(U) \to U$ is \'etale and finite. There is in fact a maximal such open which we will denote by $\U{f}$ and the preimage of the complement $\Z{f} = f^{-1}(Y - \U{f})$ will be called the exceptional locus. This is always codimension $1$ in $X$ by the Jacobian criterion (see \ref{defn: Jacobian}). We furthermore require that $\Z{f}$ has no components supported purely over the special fiber of $\Spec \A$ or equivalently, that there is some closed point on $Y$ over $\Spec \A/\pi \A$ over which the map $f$ is \'etale.

\end{definition}

\begin{definition}\label{defn: generically finite, galois}[generically finite, Galois]
Let $G$ be a constant group scheme over $\Spec \A$. A generically finite map $f: X \to Y$ is said to be Galois with Galois group $G$ if the restriction of the map to $\U{f}$ is an \'etale Galois cover with Galois group $G$ and every such automorphism extends to an automorphism of $f$. That is to say,
\[\mathrm{Aut}_\A(f:X\to Y) = \mathrm{Aut}_\A(f: f^{-1}(\U{f}) \to \U{f}) = G.\] 

\end{definition}

\begin{example}
An interesting family of generically finite, Galois maps are the following. Let $C$ be a curve of genus $g \geq 2$. Let $D$ be a degree $g$ divisor on $C$. Then, the morphism
\begin{align*}
    C^g &\to \Jac(C);\\
    (P_1,\dots,P_g) &\to [P_1] + \dots + [P_g] - D
\end{align*}
is an example of a generically finite map. It would be interesting to compute splitting densities for these maps although we don't pursue it in this paper.
\end{example}

This paper will be mainly concerned with such maps and to that end, we state a few preliminary definitions and properties of such maps here.

\begin{definition}\label{defn: decomposition group}[Decomposition group]
Let $P \in \U{f}(\L)$ be a rational point away from the ramified locus. Then, $f^{-1}(P)$ is an \'etale $\L$ algebra and each geometric point $Q \in f^{-1}(P)$ over $\overline{\L}$ defines a homomorphism $\Gal(\overline{\L}/\L) \to G$. We denote the image (known as the decomposition group of $Q$ over $P$) by $D_Q \subset G$. Correspondingly, the image of the inertia group $\Gal(\overline{L}/L^{\mathrm{ur}})$ is denoted by $I_Q \subset G$. The image of the Frobenius coset corresponds to a generating coset $\frob{Q} \in D_Q/I_Q$. We note that the pair $\tau_Q = (I_Q,\sigma_Q)$ is an \emph{admissible pair} as in \S \ref{section: admissible pairs}.

If $P \in \U{f}(\B)$, then $f^{-1}(P)$ is an unramified extension of $\B$ and a choice of $Q$ thus determines a well defined Frobenius element $\frob{Q} \in G$.
\end{definition}

\begin{remark}\label{rmk: conjugacy of stuff}
Different choices of $Q$ over $P \in U(\L)$ correspond to conjugating $D_Q$ by elements in $G$. More precisely, let $(Q_1,\dots,Q_n)$ be the geometric points of $X$ mapping to $P$. This set has a natural action of $G$ and is a $G$-torsor under this action.

If $Q_i = \lambda(Q_{j})$ for some $\lambda \in G$, then $I_{Q_i} = \lambda I_{Q_j}\lambda^{-1}$ and similarly for the decomposition groups and Frobenius cosets. Therefore, we can define $\tau_P$ as the conjugacy class of $\tau_{Q_i}$ by picking any lift $Q_i$ of $P$ to $X$ and similarly,
\[(I_P,D_P) = \{(\lambda I_{Q_i}\lambda^{-1},\lambda D_{Q_i}\lambda^{-1}) : \lambda \in G\}.\]
\end{remark}


The following definition stratifies $P \in \U{f}(\L)$ based on $\tau_P$, i.e., the data of the (conjugacy class) of the inertia group at $P$ along with a Frobenius coset. 

\begin{definition}\label{defn: rho_f,admissible pair}
For $\tau$ an admissible pair for the group $G$, we define
\[\Utype{\tau}(\L) = \{P \in \U{f}(\L) : \tau \in \tau_P\}\]
with measure $\rhotype{\tau}(\L) = \mu_Y(\Utype{\tau}(\L))$. This only depends on the conjugacy class of $\tau$. The $\Utype{\tau}(\L)$ are in fact $\L$-analytic open sets which proves that the measures $\rhotype{\tau}(\L)$ are well defined as the next lemma shows.
\end{definition}

\begin{lemma}
For any admissible pair $\tau$, the $\Utype{\tau}(\L)$ are $\L$-analytic open sets of $X(\L)$.
\end{lemma}
\begin{proof}
By Krasner's lemma \cite[Proposition 3.5.74]{poonen2017rational}, the isomorphism type of the \'etale scheme $f^{-1}(y)$ is a locally constant function as $y$ varies over $Y(\L)$ in the analytic topology. Therefore, both the Galois group of the fiber and its action on the geometric points are locally constant which proves that the corresponding admissible pair is also locally constant.
\end{proof}


We obtain a coarser stratification of $P \in \U{f}(\L)$ by remembering only the (conjugacy class of the) inertia group and decomposition group associated to $P$, i.e., we forget the precise generating coset corresponding to the Frobenius coset.

\begin{definition}\label{defn: rho_f,decomp,inertia}
For $\tau$ an admissible pair for the group $G$, we define
\[\Uid{\tau}(\L) = \{P \in \U{f}(\L) : (H_\tau,H_{1,\tau}) \in (I_P,D_P)\}\]
with measure $\rhoid{\tau}(\L) = \mu_Y(\Uid{\tau}(\L))$. As before, this only depends on the conjugacy class of $\tau$.
\end{definition}

We also have a change of variables formula for generically finite maps. We begin by defining the Jacobian.

\begin{definition}\label{defn: Jacobian}(Jacobian)
Let $f:X \to Y$ be a generically finite map over $\Spec \A$. Let $V \subset X$ and $U \subset Y$ be affine opens such that $f(V) \subset U$ and there exist gauge forms $\omega_V,\omega_U$ on $V,U$ respectively (well defined as usual up to a $p$-adic unit). We define the Jacobian of $f$ on $V$ by
\[\Jac(f) = \frac{f^*(\omega_U)}{\omega_V}.\]
It is independent of the choice of local gauge forms up to a unit in $\A$ and in particular, the norm $|\Jac(f)|$ is a well defined function on $X(\F)$.

Alternatively, one could view $\Jac(f)$ as the canonical map between the canonical bundles $f^*\mathscr O(K_Y) \to \mathscr O(K_X)$ or equivalently as a section of $f^*\mathscr O(K_Y)^{\vee}\otimes\mathscr O(K_X)$. Specifying local basis $f^*(\omega_U),\omega_V$ for the two line bundles then determines the function $\Jac(f)$ as defined above. It is clear from this description that the vanishing locus of $\Jac(f)$ is precisely the exceptional locus $Z_f \subset X$.

\end{definition}

\subsection{Galois twists}

We will use Galois twists of generically finite, Galois maps at multiple points in this paper. For convenience, we recall this definition and state some basic properties. We fix $f: X \to Y$ to be a generically finite, Galois map with Galois group $G$. 

\begin{definition}[Galois twists]
We say that $\widetilde{f}: \widetilde{X} \to Y$ defined over $\Spec \A$ is a twist of $f$ if there exists an unramified extension $\F \subset \F'$ such that the map $\widetilde{f}_{\mathcal O_{\F'}}: \widetilde{X}_{\mathcal O_{\F'}} \to Y_{\mathcal O_{\F'}}$ is isomorphic to $f_{\mathcal O_{\F'}}: X_{\mathcal O_{\F'}} \to Y_{\mathcal O_{\F'}}$.  
\end{definition}

Given an element $g \in G$ of order $r$, we can define the twist of $f$ with respect to $g$ as follows. 

\begin{definition}\label{defn: unramified twist by an element}[$g$-twist]
Let $\mathcal O_{\F'}$ be the unique degree $r$ unramified extension of $\A$ with $\Gal(\mathcal O_{\F'}/\A) = \mathbb Z/r\mathbb Z$ generated by $\frob{\F}$. We define $^gX = X\times_{\Spec \A}\Spec \mathcal O_{\F'}/\langle(g,\frob{\F})\rangle$ as a geometric quotient. As a scheme, it is defined over $\Spec \A$ and admits a map $^gf: {^gX} \to Y$. It is a standard fact that this is a twist of $f$. One sees from the definition that $^gX$ `\emph{is}' $X$ with the Frobenius acting by ${^g\frob{\F}} \coloneqq g\circ\frob{K}$ instead of $\frob{K}$. The cohomology of the twist with its Frobenius action can be identified with
\[(H^i_{\text{\'et}}({^gX}\times_{\mathbb F_q}\overline{\mathbb F}_q,\mathbb Q_\ell) , {^g\frob{\F}}) \cong (H^i_{\text{\'et}}({X}\times_{\mathbb F_q}\overline{\mathbb F}_q,\mathbb Q_\ell) , {g^{-1}\frob{\F}}). \]

\end{definition}

\begin{remark}
We note that the map $^gf: {^gX} \to Y$ does not have a constant automorphism group $G$, instead it has the \'etale automorphism group scheme $^gG$ which ``\emph{is}" $G$ with the Frobenius acting by $h \to ghg^{-1}$.
\end{remark}



\section{On the palindromicity of various natural densities}\label{section: on the palindromicity}

This section contains the main theorems of this paper. We prove that various natural densities are palindromic. Before beginning the proofs, we establish some useful definitions and notations. 


\begin{definition}\label{defn: sncd}[Simple normal crossings divisor]
Let $X/\Spec \B$ be a variety. We say that a divisor $D \subset X$ is a simple normal crossings divisor (sncd) relative to $\Spec \B$ if $X$ is smooth over $\Spec \B$ and
\begin{enumerate}
    \item $D = \bigcup_{i=1}^{r} D_i$ is a decomposition into geometrically irreducible components such that for any subset $I \subset \{1,\dots,r\}$, $D_I \coloneqq \bigcap_{i \in I}D_i$ is smooth over $\Spec \B$. We use the convention that $D_{\emptyset} = X$.
    \item For every $x \in D_I(\B)$, there exists a regular system of parameters $t_i, i\in I$ over $\B$ such that $D_i$ is cut out by $t_i$ in $\cO_{X,x}$.\footnote{i.e., $t_{j}$ is a non zero divisor in $\cO_{X,x}/(t_{i_1},\dots,t_{i_r})$ for $j \not\in \{i_1,\dots,i_r\} \subset I$ and $\cO_{X,x}/(t_{i_1},\dots,t_{i_r})$ is smooth over $\B$.}
\end{enumerate}
\end{definition}

\begin{definition}\label{defn: geometric sncd}[Geometric sncd]
Let $X/\Spec \A$ be a variety. We say that a divisor $D\subset X$ is a geometric sncd if there exists an unramified extension $\A \subset \mathcal O_{K'}$ and $D_{\mathcal O_{K'}} \subset X_{\mathcal O_{K'}}/\Spec \mathcal O_{K'}$ is a sncd relative to $\Spec \mathcal O_{K'}$. In this case, we write $D = \bigcup_{i=1}^{r}D_i$ with the $D_i$ irreducible over $\Spec \A$ (but possibly not geometrically irreducible). For each $i \leq r$, we define $\F_i$ to be the smallest unramified extension over which $D_{i,\mathcal O_{K_i}} = \bigcup_{j \in J_i}D_{ij.\mathcal O_{K_i}}$ with the $D_{ij}$ geometrically irreducible (and hence smooth). In this case, the Frobenius automorphism $\frob{\F}$ of $\F_i/\F$ will transitively permute the $D_{ij}$ and we define $\deg(\mathcal O_{K_i}/\A) = k_i$. We will consider the set $J_i = \{{1},\dots,{k_i}\}$ to be equipped with the action of the Frobenius $\frob{\F}$ (corresponding to the action on the divisors $D_{ij}$) and denote by $J_i/\frob{\F}^\ell$ the set of orbits for the action of $\frob{\F}^\ell$. At each closed point $x \in X$, we let $t_{i,x} \in \cO_{X,x}$ be a local function cutting out $D_i$ and similarly, we define $t_{ij,x} \in \cO_{X,x}\otimes_{\A}\mathcal O_{K'}$ cutting out $D_{ij}$ such that $\frob{\F}(t_{ij}) = t_{i\frob{\F}(j)}$.

\end{definition}


We note that $D \subset X$ being a geometric sncd is invariant under taking unramified twists as in Definition \ref{defn: unramified twist by an element} since this condition can be checked over the unramified closure.

\begin{hypothesis}
\label{hypothesis:res of sing}
A tuple $(g,X\to\Spec\A, Z \subset X)$ where $g \in \Aut(X)$ and $Z \subset X$ is a $g$-equivariant closed subscheme has an equivariant resolution if there exists a smooth variety $\widetilde{X} \to \Spec \A$ with an automorphism $g: \widetilde{X} \to \widetilde{X}$ and a proper, $g$-equivariant birational map $\pi: \widetilde{X} \to X$ such that $\pi^{-1}(Z)$ is a geometric sncd. 

\end{hypothesis}

% \yif{don't know where to put the following lemma... also not sure if we really need this, I want to make sure every quotient is smooth, and possibly have rationality result from this. Let me think...}
% \as{I think we need something like the following for the proof of the twist symmetry. Found a reference! https://mathoverflow.net/a/332698/58001}


% \begin{lemma}[Replacement Lemma]\label{lem: replacement lemma}
% Let $X$ be an irreducible projective variety over a field $k$ with char $= 0$. Suppose $G$ is a constant finite group acting on $X$, let $Y = X/G$ be the quotient variety. Then there exists smooth birational models $X' \to X$ and $Y' \to Y$ and a morphism $X' \to Y'$ that is generically Galois with Galois group $G$. In other words, the following diagram commutes, with horizontal maps being birational projective (proper) and $G$-equivariant, and vertical maps being generically Galois with Galois group G. \yif{diagram commented out to speed up compilation.}
% % \[
% % \begin{tikzcd}
% % X' \arrow[r]\arrow[d] & X \arrow[d] \\
% % Y' \arrow[r] & Y
% % \end{tikzcd}
% % \]
% \end{lemma}
% \begin{proof}
% Let $Y' \to Y$ be a resolution of singularity of $Y$, consider the pullback $Y' \times_Y X$ along the quotient map $X \to Y$. Suppose $U \subset X$ is a $G$-invariant open subscheme where $U \to U/G \subset Y$ is \'etale. Then $U$ naturally embeds into $Y' \times_Y X$, denote the Zariski closure of $U$ by $Z$, then $G$ acts on $Z$. Take $X' \to Z$ be a $G$-equivariant resolution of singularity, and we have the desired model.
% \end{proof}


This hypothesis need not be satisfied even if $g = \mathrm{id}, X = \mathbb P^1$ as can be seen by considering $D = V(x^2 - p^2y^2) \subset \mathbb P^1$. The minimal resolution of singularities is given in this example by blowing up $X$ at the closed point $V(t, p)$ and $\pi^{-1}(D)$ has a component purely supported in characteristic $p$ and hence is not a relative normal crossings divisor. Nevertheless, it is generally not a restrictive hypothesis.

\begin{remark}\label{rmk: char 0 resolutions}
If $f: X \to Y$ is generically finite with Galois group $G$ over (an open subset) of the ring of integers of a number field, then the above hypothesis for all the tuples $(g,X,Z_f)$ where $g \in G$ is simultaneously satisfied for all but finitely many $p$ since we can equivariantly resolve singularities over characteristic $0$ (for instance, see \cite{abramovich1996equivariant},\cite{encinas1998good} or \cite{bierstone1997canonical}) and \emph{spread out} to all but finitely many primes.
\end{remark}


Before proving the main theorems of this section, we first do a local computation. For any predicate $\mathscr P$ (such as ``$\gcd(k,m) = \ell$"), we define $\cI(\mathscr P)$ to be the indicator function that is $1$ when the predicate is satisfied and zero otherwise.

Since our divisors are not geometrically irreducible, our formulas depend on the divisibility properties of the size of the residue field of $\L$ and we use the indicator functions to concisely express this.

\begin{lemma}\label{lemma: local integrals}
Let $D \subset X$ be a geometric sncd as in Definition \ref{defn: geometric sncd} (with the same notation).  For any $e_1,\dots,e_r \in \mathbb Z_{\geq 0}$ and $M_1 \subset \{1,\dots,k_1\},\dots, M_r \subset \{1,\dots,k_r\}$, define $D_{M_i} = \bigcap_{j \in M_i}D_{ij}$ and let $x \in D_{M_1,\dots,M_r} \coloneqq \bigcap_{i=1}^rD_{M_i}$ be a point defined over an extension $\B \supset \A$.  We have the identity
\begin{align*}
    \rhol_{x,J_1,\dots,J_r}(\L) &\coloneqq \int_{\widehat\cO_{X,x}(\B)}|t_{1,x}|_{\L}^{e_1}\dots |t_{r,x}|_{\L}^{e_r}d\mu_{X_{\SS}}\\
    &= \sum_{\ell_1|k_1,\dots,\ell_r|k_r}\prod_{i=1}^r\cI(\gcd(k_i,m)=\ell_i)\left(\prod_{i=1}^r\rhol_{e_i}^{|M_i/\frob{\F}^{\ell_i}|}(m;q^{k_i/\ell_i})\right),
\end{align*}
where $m$ is the degree of the residue field of extension of $L/K$ and $\rhol_k(m;q)$ is as in Definition \ref{defn: rho, eta}.
\end{lemma}

We make the preliminary observation that $\frob{\L}$ action preserves the subsets $M_1,\dots,M_r$ since $x \in \bigcap_{i=1}^rD_{M_i}$.


\begin{proof}
Recall Definition \ref{defn: geometric sncd}, i.e., we have minimal unramified extensions $\A\subset \mathcal O_{K_i}$ of degree $k_i$ over which the $D_i = \bigcup_{j\in J_i}D_{ij}$ decompose into smooth components for $J_i = \{1,\dots,k_i\}$. We let $\L_i$ be the compositum of $\L$ with ${K_i}$ so that $\L \subset \L_i$ has degree $\ell_i = \gcd(m,k_i)$ and take  $\L'$ to be the compositum of all the $\L_i$. As before, we define local parameters $t_{ij}$ corresponding to the $D_{ij}$ with $\frob{\L}(t_{ij}) = t_{i\frob{\L}(j)}$ so that in particular, we have $t_i \coloneqq t_{i,x} = \prod_{j \in J_i}t_{ij}$. 

Consider the following map of $\B$ modules with a $\frob{\L}$ action \[\mathfrak m_{X,x} \to \mathfrak m_{X_{\mathcal O_{L'}},x}.\]
The preimage of the span of the $t_{ij}$  is a $\B$-submodule $\cL \subset \mathfrak m_{X,x}$ with the property that its generic rank is equal to its rank modulo the uniformizer $\pi_\L$ precisely because the $t_{ij}$ are a regular system of parameters and $\B \subset \mathcal O_{L'}$ is a faithfully flat extension. Therefore, we can extend $\cL$ to a spanning lattice $\cL'$ by adding in generators $s_1,\dots,s_v$.

Consider now the Frobenius equivariant isomorphism of free $\B$ modules 
\begin{equation}\label{eqn: co-ord isomorphism}
    \widehat\cO_{X_{\mathcal O_{L'}},x}(\mathcal O_{L'}) \xrightarrow{\sim} \prod_{i=1}^{r}\prod_{j \in J_i}\mathfrak m_{\L'}\times (\mathfrak m_{\L'})^v; P \to (t_{11}(P),t_{12}(P),\dots,t_{rk_r}(P),s_1(P),\dots,s_v(P)),
\end{equation} 
where the Frobenius $\frob{\L}$ acts on the right by
\[(x_{11}, x_{12},\dots, x_{rk_r}, y_1, \dots, y_v) \to (\frob{\L}(x_{1\frob{\L}^{-1}(1)}),\dots, \frob{\L}(x_{r\frob{\L}^{-1}(k_r)}), \frob{\L}(y_1), \dots, \frob{\L}(y_v)).\]
On taking $\frob{\L}$ invariants on both sides, we obtain
\[ \widehat\cO_{X,x}(\B) \cong \prod_{i=1}^r\prod_{\lambda \in J_i/\frob{\L}}\mathfrak m_{\L_i}\times (\mathfrak m_{\L})^v.\]
This isomorphism preserves the additive structure and is therefore measure preserving since the induced measure is the Haar measure with the correct normalization. All together, our integral thus reduces to
\begin{align*}
    \int_{\widehat{\cO}_{X,x}(\B)}|t_{1,x}|_{\L}^{e_1}\dots |t_{r,x}|_{\L}^{e_r}d\mu_{X_{\SS}} &= \int_{{\widehat{\cO}}_{X,x}(\B)}\prod_{i=1}^r\prod_{j=1}^{k_j}|t_{ij}|_{\L}^{e_i}d\mu_{X_{\SS}}\\
    &= \prod_{i=1}^r\prod_{\lambda \in M_i/\frob{\L}}\int_{\mathfrak m_{\L_i}}\left(\prod_{j \in \lambda}|t_{ij}|_{\L_i}^{e_i}\right)^{\ell_i/k_i}d\mu_{\mathfrak m_{\L_i}},
\end{align*}

where in the second equality, we use the above identification of our region of integration and the fact that $|t_{ij}| = 1$ over our region of integration unless $j \in M_i$. The functions $|t_{ij}|$ depend only on the $\frob{\L}$ orbit $\lambda$ of $j$ since $|\frob{\L}t_{ij}(P)| = |t_{ij}(P)|$ and we denote this function by $|t_{i\lambda}|$. Moreover, the size of the orbit $\lambda$ is exactly $k_i/\ell_i$ and orbits of $\frob{\L}$ are exactly the same same as the orbits of $\frob{\F}^{\ell_i}$ on the set $J_i$.

Therefore, the inner integrand simplifies to $|t_{i\lambda}|^{e_i}$ and the above integral is equal to
\[\prod_{i=1}^r\prod_{\lambda \in M_i/\frob{\F}^{\ell_i}}\int_{\mathfrak m_{\L_i}}|t_{i\lambda}|_{\L_i}^{e_i}d\mu_{\mathfrak m_{\L_i}} = \prod_{i=1}^{r}\rhol_{e_i}^{|M_i/\frob{\F}^{\ell_i}|}(m;q^{k_i/\ell_i})\]
since the residue field of $\mathcal O_{L_i}$ has size $q^{k_i/\ell_i}$.

Since we assumed that $\gcd(k_i,m) = \ell_i$, the proof is completed by summing over all the possibilities for the $\ell_i$ the corresponding indicator functions weighted by the above product.

\end{proof}

\begin{thm}\label{thm: symmetry for pullback}
Let $X,Y$ be \emph{smooth proper} varieties over $\Spec \A$ with $f: X \to Y$ a generically finite map with $(\mathrm{id},X,Z_f)$ satisfying Hypothesis \ref{hypothesis:res of sing}. The function
\[\L \to \eta_f(\L) = \int_{X(\B)} f_{\SS}^{*}(d\mu_{X_{\SS}})\]
is a palindromic form of weight $\dim Y$. In fact, if $m$ is the degree of the residue field extension of $L$ and we pick a resolution of singularities of $\widetilde\pi : \widetilde{X} \to X$ as in Hypothesis \ref{hypothesis:res of sing} with $\pi^{-1}(Z_f)$ playing the role of $D$ in Lemma \ref{lemma: local integrals} with the same notation as before, we have the identity
\[\eta_f(\L) = \sum_{\ell_1|k_1,\dots,\ell_r|k_r}\prod_{i=1}^r\cI(\gcd(m,k_i) = \ell_i)\left(\sum_{\cL }\prod_{i=1}^{r}\eta_{e_i}^{|\Lambda_i|}(m;q^{k_i/\ell_i})\left|D_{\cL}(\mathbb F_{q^m})\right|\right)\]
where the $e_i$ are positive integers and $\cL$ ranges over tuples of the form $(\Lambda_1,\dots,\Lambda_r)$ with $\Lambda_i \subset J_i/\frob{\F}^{\ell_i}$ and $D_\cL \coloneqq \bigcap_{i=1}^r\bigcap_{\lambda \in \Lambda_i}\bigcap_{j\in\lambda}D_{ij}$.
\end{thm}
\begin{proof}

By Hypothesis \ref{hypothesis:res of sing} and the change of variables formula for birational maps, we can replace $X$ by a resolution of singularities $\widetilde{X}$ and hence suppose that the exceptional locus $Z_f \subset X$ is a normal crossings divisor relative to $\Spec \A$. Let $x \in X$ be a closed point. Since $X/\Spec \A$ is smooth, $\widehat\cO_{X,x} \cong \A[\![x_1,\dots,x_n]\!]$ for $n = \dim X$ which is a unique factorization domain by the Weierstrass preparation theorem.

Note that $\Jac(f)$ (as in Definition \ref{defn: Jacobian}) is a local function that precisely cuts out $Z_f$ geometrically, well defined up to a function that is invertible when evaluated on ${\mathcal O}_{\overline\F}$ points. If $Z_f = \bigcup_{i=1}^rD_i$ as in the statement of Lemma \ref{lemma: local integrals}, we have $\Jac(f) = u\prod_{i=1}^rt_{i,x}^{e_i} \in \widehat\cO_{X,x}$ for some $u,e_i \geq 0$ with the $u$ coprime to the $t_{i,x}$. In fact, the $e_i$ are exactly given by the order of vanishing of $\Jac(f)$ at the $D_i$ (computed on any local chart). This is well defined precisely because $\Jac(f)$ is well defined up to a unit on that chart. Due to this, we see that
$u = \Jac(f) \prod_i t_{i,x}^{-e_i}$ does not vanish away from $Z_f$ and also does not vanish at the generic point of each $D_i$. Therefore, it is non-vanishing on an open with complement of codimension at least $2$ and hence a unit in $\widehat{\cO}_{X,x}$.

Finally, recall the notation from before that $D_i = \bigcup_{j \in J_i}D_{ij}$ over an unramified extension $\mathcal O_{L_i}/\A$ with the $D_{ij}$ smooth (with $\mathcal O_{L_i}$ the minimal such extension). Given $\cM = (M_1,\dots,M_r)$ with the $M_i \subset J_i$, we define by $D_{\cM}^\circ$ the locally closed strata of $D_{\cM} = \bigcap_{j \in J_i} D_{ij}$ not contained in $D_{\cM'}$ for any other $\cM'$ corresponding to a smaller strata. We carry over other notation from Lemma \ref{lemma: local integrals} with $Z_f$ playing the role of $D$.

In order to compute the integral, we pick a set of points $\cS \subset X(\B)$ in bijection with the points of $X(\mathbb F_{q^m})$ under the reduction map. This gives rise to a decomposition of the $p$-adic analytic set $X(\B)$ into discs around each point $s \in \cS$ isomorphic to $\widehat{\cO}_{X,s}$, giving rise to the equality 
\begin{align*}
    \int_{X(\B)}f_{\SS}^{*}d\mu_{X_\SS} &= \sum_{\cM}\sum_{s \in \cS\cap D_{\cM}^\circ}\int_{\widehat\cO_{X,s}(\B)}|t_1^{e_1}\dots t_r^{e_r}|_{\L}d\mu_{X_\SS}.
\end{align*}
By Lemma \ref{lemma: local integrals}, the above is equal to
\[\sum_{\ell_1|k_1,\dots,\ell_r|k_r}\prod_{i=1}^r\cI(\gcd(m,k_i) = \ell_i)\left(\sum_{\cM}\sum_{s \in \cS \cap D_{\cM}^\circ}\prod_{i=1}^{r}(\eta_{e_i}(m;q^{k_i/\ell_i})+1)^{|M_i/\frob{\F}^{\ell_i}|}\right),\]
where we recall that $\delta_e(m;q) = \eta_{e}(m;q)+1$. Since $s \in D_\cM(\B)$, the $M_i$ are necessarily $\frob{L}$, and hence, $\frob{\F}^{\ell_i}$ stable. Next, we pass to the orbits by defining $\cL = (\Lambda_1,\dots,\Lambda_r)$ with $ \Lambda_i = M_i/\frob{K}^{\ell_i} \subset J_i/\frob{\F}^{\ell_i}$. Summing over such $\cL$ and noting that $|\cS \cap D_\cL^\circ| = D_\cL^\circ(\mathbb F_{q^m})$, we obtain
\[\sum_{\ell_1|k_1,\dots,\ell_r|k_r}\prod_{i=1}^r\cI(\gcd(m,k_i) = \ell_i)\left(\sum_{\cL \text{ as above}}\left|D^\circ_{\cL}(\mathbb F_{q^m})\right|\prod_{i=1}^{r}\prod_{\lambda \in \Lambda_i}\left(\eta_{e_i}(m;q^{k_i/\ell_i})+1\right) \right).\]
We now expand the inner product as
\[\prod_{\lambda \in \Lambda_i}\left(\eta_{e_i}(m;q^{k_i/\ell_i})+1\right) = \sum_{\Lambda_i' \subset \Lambda_i}\eta_{e_i}^{|\Lambda_i'|}(m;q^{k_i/\ell_i}).\]
We say that $\cL' = (\Lambda_1',\dots,\Lambda_r') \leq \cL$ when $\Lambda_i' \subset \Lambda_i$ for all $i$. Then, the above expansion lets us rewrite the integral as
\[\sum_{\ell_1|k_1,\dots,\ell_r|k_r}\prod_{i=1}^r\cI(\gcd(m,k_i) = \ell_i)\left(\sum_{\cL \text{ as above}}\sum_{\cL' \leq \cL}\prod_{i=1}^{r}\eta_{e_i}^{|\Lambda_i'|}(m;q^{k_i/\ell_i})|D^\circ_{\cL}(\mathbb F_{q^m})|\right).\]
Switching the order of summation, we have
\[\sum_{\ell_1|k_1,\dots,\ell_r|k_r}\prod_{i=1}^r\cI(\gcd(m,k_i) = \ell_i)\left(\sum_{\cL'}\left(\prod_{i=1}^{r}\left(\eta_{e_i}^{|\Lambda_i'|}(m;q^{k_i/\ell_i})\right) \sum_{\cL' \leq \cL}|D^\circ_{\cL}(\mathbb F_{q^m})|\right)\right).\]
The claim now follows from the observation that $\sum_{\cL' \leq \cL}|D^\circ_{\cL}(\mathbb F_{q^m})| = |D_{\cL}(\mathbb F_{q^m})|$ and is thus a palindromic form of weight $\dim D_{\cL}$ while the $\eta_e(m;q^d)$ factors are palindromic forms of weight $d$. Together, 
\[\prod_{i=1}^{r}\eta_{e_i}^{|\Lambda_i|}(m;q^{k_i/\ell_i})|D_{\cL}(\mathbb F_q)|\]
is a palindromic form of weight \[\sum_{i=1}^r|\Lambda_i|\frac{k_i}{\ell_i} + \dim D_\cL = \dim X\]
since the first term above is precisely the number of smooth components of $Z_f$ cutting out $D_\cL$ because each $\lambda \in \Lambda_i$ corresponds to an orbit of divisors of size $k_i/\ell_i$. 

The indicator functions $\cI(\ell|\gcd(m,k))$ (for $\ell|k$) can be written as a sum of roots of unity $\zeta_k$ of order $k$ as 
\[\cI(\ell|\gcd(m,k)) = \frac{1}{k}\sum_{a=1}^{k}\zeta_{k}^{amk/\ell}.\]
Therefore, they are palindromic forms of weight $0$. Since
\begin{equation}\label{eqn: open stratification formula}
    \cI(\gcd(m,k) = \ell) = \cI(\ell|\gcd(m,k)) - \left(1-\prod_{1 < d, d\ell|k}\left(1-\cI(d\ell|\gcd(m,k))\right)\right),
\end{equation}
$\cI(\gcd(m,k) = \ell)$ is also a palindromic form of weight $0$. All together, this shows that our integral is a palindromic form of weight $\dim Y = \dim X$.
\end{proof}

Despite appearances, the function $\eta_f$ does not depend on the resolution $\pi: \widetilde X \to X$ by the birational change of variables formula. We record the following simplification when all components of the exceptional divisor are geometrically irreducible over $\Spec \A$.

\begin{corollary}\label{cor: pullback formula in sncd case}
Let $m$ be the degree of the residue field extension of $L/K$. Let $f: X \to Y$ be as in the previous theorem and suppose that there exists a resolution of singularities $\pi: \widetilde{X} \to X$ over $\A$ with $\pi^{-1}(Z_f) = \bigcup_{i=1}^r D_i$ a sncd over $\A$. Then
\[\eta_f(\L) = \sum_{J \subset \{1,\dots,r\}}\left|D_{J}(\mathbb F_{q^m})\right|\prod_{j\in J}\eta_{e_j}(m;q)\]
for some positive integers $e_j$.
\end{corollary}

We prove a modification of the above theorem taking into consideration unramified twists, as in Definition \ref{defn: unramified twist by an element}.

Let $\tau = (H,gH)$ be an admissible pair, i.e., $gH=Hg$. As above, we take $f: X \to Y$ to be a generically finite, Galois map with Galois group $G$ over $\Spec \A$. Define $f_{H}: X \to X/H$ to be the natural quotient map. Since $gH = Hg$, $g$ descends to an automorphism of $X/H$.

We suppose that $(g,X/H,f_H(Z_f))$ satisfies Hypothesis \ref{hypothesis:res of sing} and define $\widetilde{X/H}$ to be a $g$-equivariant resolution
\[\pi_H: \widetilde{X/H} \to X/H.\]
Since all the maps involved are $g$-equivariant and twisting by $g$ is functorial, we obtain maps
\[\pi_\tau: {^g\left(\widetilde{X/H}_{\SS}\right)} \to {^g{(X/H)_{\SS}}}\]
where we base change to $\SS$ \emph{first} and then twist by $g$. We denote ${^g\left(\widetilde{X/H}_{\SS}\right)}$ by ${^\tau X_{\SS}}$ and the resulting map by ${^\tau f_{\SS}}: {^\tau X_{\SS}} \to Y$. By construction, ${^\tau f_{\SS}}$ is an \'etale map over the unramified open locus $\U{f}$.

\begin{thm}\label{thm: twist symmetry under pullback}
Let $f: X \to Y$ be a generically finite, Galois map with group $G$ between \emph{smooth, proper} varieties with $(g,X/H,f_H(Z_f))$ and $(g^{-1},X/H,f_H(Z_f))$ satisfying Hypothesis \ref{hypothesis:res of sing} as above and define
\[\L \to \eta_{\tau}(\L)\coloneqq \int_{{^{\tau^{-1}} X_{\SS}}(\B)}\Jac({^{\tau^{-1}} f_\B})d\mu_{{^{\tau^{-1}} X_{\SS}}}.\]
Then, $\eta_\tau(\L), \eta_{\tau^{-1}}(\L) \in \mathbb Q$ and $\eta_{\tau}(\L) + \eta_{\tau^{-1}}(\L)$ is a palindromic form of weight $\dim Y$.
\end{thm}

\begin{proof}

We maintain the notation from the beginning of the proof of Theorem \ref{thm: symmetry for pullback} so that the exceptional locus for ${^\tau f}_{\SS}: {^\tau X}_{\SS} \to Y_{\SS}$ is $\bigcup_{i=1}^r {^gD}_{i,\SS}$ with the $D_i$ defined over $\Spec \A$ and irreducible, $D_i = \bigcup_{j \in J_i}D_{ij}$ and so on. By Theorem \ref{thm: symmetry for pullback}, 
\[\eta_{\tau^{-1}}(\L) = \sum_{\ell_1|k_1,\dots,\ell_r|k_r}\prod_{i=1}^r\cI(\gcd(m,k_i) = \ell_i)\left(\sum_{\cL }\prod_{i=1}^{r}\eta_{e_i}^{|\Lambda_i|}(m;q^{k_i/\ell_i})\left(|{^gD}_{\cL,\SS}(\mathbb F_{q^m})|\right)\right) \in \mathbb Q,\]
while $\eta_{\tau}(\L) + \eta_{\tau^{-1}}(\L)$ is equal to
\[\sum_{\ell_1|k_1,\dots,\ell_r|k_r}\prod_{i=1}^r\cI(\gcd(m,k_i) = \ell_i)\left(\sum_{\cL }\prod_{i=1}^{r}\eta_{e_i}^{|\Lambda_i|}(m;q^{k_i/\ell_i})\left(|{^gD}_{\cL,\SS}(\mathbb F_{q^m})| + |{^{g^{-1}}D}_{\cL,\SS}(\mathbb F_{q^m})|\right)\right).\]
It is possible (and indeed likely!) that even though $D_{\cL,\SS}$ has $\mathbb F_{q^m}$ points, the twist ${^gD}_{\cL,\SS}$ is not defined over $\B$ and consequently has no $\mathbb F_{q^m}$ points. We emphasize that the twists ${^gD}_{\cL,\SS}$ are obtained after base changing \emph{first}. Therefore, it is not obvious a-priori that the above expression is a palindromic form of weight $\dim Y$. Nevertheless, we will prove that it is so using the Lefschetz trace formula. 

To that end, fix an auxiliary prime $\ell$ and a smooth, projective variety $D/\mathbb F_q$. As is well known, the action of $\frob{\L}$ on $H_{\text{\'et}}^*({^{g^{-1}}}D\otimes_{\mathbb F_q}\overline{\mathbb F}_q,\mathbb Q_\ell)$ is equivalent to the action of $g\frob{\L}$ on $H_{\text{\'et}}^*(D\otimes_{\mathbb F_q}\overline{\mathbb F}_q,\mathbb Q_\ell)$. Thus, in conjunction with Grothendieck-Lefschetz trace formula, we obtain 
\[|{^{g^{-1}}D}_{\cL,\SS}(\mathbb F_{q^m})| = \sum_{i=0}^{2\dim D}(-1)^i\tr\left(g\frob{q}^m|H^i_{\text{\'et}}(D_{\cL}\times_{\mathbb F_q}\overline{\mathbb F}_q,\mathbb Q_\ell)\right),\]
and consequently,
\begin{align*}
    |{^gD}_{\cL,\SS}(\mathbb F_{q^m})|+|{^{g^{-1}}D}_{\cL,\SS}(\mathbb F_{q^m})| &= \sum_{i=0}^{2\dim D}(-1)^i\tr\left((g+g^{-1})\frob{q}^m|H^i_{\text{\'et}}(D_{\cL}\times_{\mathbb F_q}\overline{\mathbb F}_q,\mathbb Q_\ell)\right)\\
    &= \rho_{D_\cL,g}(m).
\end{align*} Thus, our integral above reduces to
\[\eta_{\tau}(\L) + \eta_{\tau^{-1}}(\L) = \sum_{\ell_1|k_1,\dots,\ell_r|k_r}\prod_{i=1}^r\cI(\gcd(m,k_i) = \ell_i)\left(\sum_{\cL }\prod_{i=1}^{r}\eta_{e_i}^{|\Lambda_i|}(m;q^{k_i/\ell_i})\rho_{D_\cL,g}(m)\right)\]
where $\rho_{D_\cL,g}$ is as in Lemma \ref{lemma: twisted point counts have the right weight}. By lemma, $\rho_{X,g}$ is a palindromic form of weight $\dim D_\cL$ and the proof is completed exactly as in Theorem \ref{thm: symmetry for pullback}.
\end{proof}

\begin{thm}\label{thm: incidence identity}
Let $f:X\to Y$ be a generically finite, Galois map with Galois group $G$. For an admissible pair $\tau = (H,gH)$, suppose that the tuple $(g^{-1},X/H,f_H(Z_f))$ satisfies Hypothesis \ref{hypothesis:res of sing}. Then, we have
\begin{equation}\label{eqn: incidence identity}
     {\eta_{\tau}(\L)} = \sum_{\tau' \leq \tau}\rhotype{\tau'}(\L)\gamma([\tau',\tau]),
\end{equation}
where $\rhotype{\tau}$ is as in Definition \ref{defn: rho_f,admissible pair} and $\gamma([\tau',\tau])$ are as in \S\ref{section: admissible pairs}.
\end{thm}
\begin{proof}
Suppose $\tau = (H,gH)$ and for ease of notation, we let $k = g^{-1}$ with $\nu = \tau^{-1}$. Let $\widetilde{U} = {^k f_{\SS}}^{-1}(U_f), \widetilde{U}_\nu = {^k\left( f_{\SS}^{-1}(U_f)/H\right)}$ with quotient map $f_\nu: \widetilde{U}_\nu \to \U{f}$.

Suppose $P \in \U{f}(\L)$ with geometric pre-images $Q_1,\dots,Q_n$ in $X(\overline{\L}) = {^kX}(\overline{\L})$ (where the two sets are identified using the canonical identification with the only difference being the Frobenius action). For such a $Q_i$, we denote the corresponding inertia group and Frobenius coset in $G$ by $I_i$ and $\frob{i}$, respectively. We similarly define the group schemes ${^kI}_i$ and ${^k\frob{i}}$ in $^k G$. In fact, ${^kI}_i = I_i$ under the canonical identification of ${^kG}(\B^{\mathrm{ur}}) \cong G(\B^{\mathrm{ur}}) = G$ over the unramified closure of $\B$ since such a base change does not affect inertia groups. Finally, we let $\tau_{Q_i} = (I_i,\frob{i})$ be the corresponding admissible pair. 

Let $\overline{Q}_i$ be the image of $Q_i$ in $\widetilde{U}_\nu$. It lands in $\widetilde{U}_\nu(\L)$ precisely when both the inertia group and Frobenius fix it. Therefore
\begin{align*}
    \overline{Q}_i \in \widetilde{U}_\nu(\L) &\iff I_i\overline{Q}_i = \overline{Q}_i \text{ and } {^k\frob{i}}(\overline{Q}_i) = \overline{Q}_i\\ 
    &\iff I_i \subset H \text{ and } k\frob{i} \in H
\end{align*}
where the final equivalence follows from the faithfulness of the action of $G$ on $Q_1\dots,Q_n$. Altogether, we see that the image of $Q_i$ in ${^\nu X_{\SS}}$ is $\B$-rational precisely when $\tau_{Q_i} \leq \tau$, i.e., $I_i \subset H$ and $\frob{i} \in gH$. In other words, we have shown that $f_\nu$ is a covering space map with image 
\[f_\nu(\widetilde{U}_\nu(\L)) = \bigsqcup_{\substack{\tau' \leq \tau\\\text{ up to conjugacy }}}\Utype{\tau'}(\L).\]

For every such $P$ with some $i$ such that $\tau_i \leq \tau$, the $j$ such that $\tau_j \leq \tau$ correspond to $\lambda \in G$ such that $\lambda({^kI}_i,{^k\frob{i}})\lambda^{-1} \leq \tau$. The number of distinct images of these $Q_j$ in ${^\nu X_{\SS}}$ is precisely equal to $\alpha([({^kI}_i,{^k\frob{i}}),\tau])$ since the ones in the same $H$ orbit get identified (recall the definitions of $\alpha,\beta$ in Definition \ref{defn: alpha, beta}).

Therefore, the degree of $f_\nu$ over $\Utype{\tau'}(\L)$ is precisely $\alpha(\tau',\tau)$. If we pick a compact open $V \subset \Utype{\tau'}(\L)$ and a differential form $\omega$ on it, we have
\[\alpha(\tau',\tau)\int_{V}|\omega| = \int_{f_{\nu}^{-1}(V)}|{f_{\nu}}^*\omega|\] 
by Theorem \ref{thm: change of variables}. We pick $\omega$ to be a gauge form with respect to the measure $\mu_Y$ and sum over a disjoint union cover $V_i$ of $\bigsqcup_{\tau'\leq \tau}\Utype{\tau'}(\L)$ to obtain the required identity
\[\sum_{\tau'\leq \tau}\rhotype{\tau'}(\L)\frac{\alpha(\tau',\tau)}{\beta(\tau',\tau)} = \int_{\widetilde{U}_\nu(\L)}|\Jac(f_{\nu})|d\mu_{X_{\SS}}.\]
We divide by $\beta(\tau',\tau)$ in the first term because that is precisely the amount we over-count by when we sum over all types $\tau' \leq \tau$ instead of up to conjugacy. 
\end{proof}


\begin{thm}\label{thm: Main theorem, density fixing inertia and frobenius coset}
Let $f:X\to Y$ be a generically finite, Galois map with Galois group $G$ between \emph{smooth, proper} varieties and $\tau$ an admissible pair (with respect to $G$). For every admissible pair $\tau' \leq \tau$ and $\tau' \leq \tau^{-1}$, suppose that the tuples $(g_{\tau'},X/H_{\tau'},f_{H_{\tau'}}(Z_f))$ satisfy Hypothesis \ref{hypothesis:res of sing}. Then, the densities $\rhotype{\tau}(\L) \in \mathbb Q$ and the function
\[\L \to \rhotype{\tau}(\L) + \rhotype{\tau^{-1}}(\L)\]
is a palindromic form of weight $\dim Y$, where $\rhotype{\tau}$ is as in Definition \ref{defn: rho_f,admissible pair}. We note that the above function is the measure of the set $\Utype{\tau}(\L) \cup \Utype{\tau^{-1}}(\L)$ since these are disjoint subsets of $Y(\L)$.
\end{thm}
\begin{proof}

We fix $\L$ and elements $\widetilde{\eta},\widetilde{\rho}$ in the incidence algebra for the poset of types satisfying (for all $g \in G$ and types $\tau' \leq \tau$) 
\[\widetilde{\rho}([(e,g),\tau]) = \rhotype{\tau}(\L),\]
\[\widetilde{\eta}([\tau',\tau]) = \widetilde\rho\ast\gamma([\tau',\tau]).\]
By Equation \ref{eqn: incidence identity}, we see that $\widetilde{\eta}([(e,g),\tau]) = \eta_\tau(\L)$. Since $\gamma([\tau,\tau]) > 0$, $\gamma$ is an invertible element of the incidence algebra so that
\[\widetilde\rho = \widetilde\eta\ast\gamma^{-1}\]
and evaluating at $[(e,g),\tau]$ for some type $\tau \geq (e,g)$, we obtain
\[\rhotype{\tau}(\L) = \sum_{\tau'\leq \tau}{\eta_{\tau'}(\L)}\gamma^{-1}([\tau',\tau]).\]
This is a rational number since $\eta_{\tau'}(\L) \in \mathbb Q$ (Theorem \ref{thm: twist symmetry under pullback}) and $\gamma([\tau',\tau]) \in \mathbb Q$ (and hence also its inverse in the incidence algebra). By definition, we have $\alpha([\tau',\tau]) = \alpha([\tau'^{-1},\tau^{-1}])$ and similarly for $\beta$ and hence also for $\gamma$ and $\gamma^{-1}$. Therefore, if we sum the above identity with the corresponding one for $\tau^{-1}$, we obtain
\begin{align*}
    \rhotype{\tau}(\L) + \rhotype{\tau^{-1}}(\L) &= \sum_{\tau'\leq \tau}{\eta_{\tau'}(\L)} \gamma^{-1}([\tau',\tau]) + \sum_{\tau'^{-1}\leq \tau^{-1}}{\eta_{\tau'^{-1}}(\L)}\gamma^{-1}([\tau'^{-1},\tau^{-1}])\\
    &=\sum_{\tau'\leq \tau}\left(\eta_{\tau'}(\L) + \eta_{\tau'^{-1}}(\L)\right){\gamma^{-1}([\tau',\tau])}.
\end{align*}
By Lemma \ref{lemma: twisted point counts have the right weight}, we see that $\eta_{\tau'}(\L) + \eta_{\tau'^{-1}}(\L)$ is a palindromic form of weight $\dim Y$ and since the poset of admissible pairs and $\gamma^{-1}$ depend only on the group $G$ and not on $\L$, the above identity shows that $\rhotype{\tau}(\L) + \rhotype{\tau^{-1}}(\L)$ is a palindromic form of weight $\dim Y$.
\end{proof}

As an immediate corollary, we have that the measures in Definition \ref{defn: rho_f,decomp,inertia} take rational values and are palindromic forms of the correct weight.

\begin{corollary}\label{cor: symmetry for fixed inertia, decomp}
Let $f: X \to Y$ be a generically finite, Galois map between \emph{smooth, proper} varieties. For every admissible pair $\tau' \leq \tau$ and $\tau' \leq \tau^{-1}$, suppose that the tuples $(g_{\tau'},X/H_{\tau'},f_{H_{\tau'}}(Z_f))$ satisfy Hypothesis \ref{hypothesis:res of sing}. 

Then the measures $\rhoid{\tau}(\L)$ are palindromic forms of weight $\dim Y$.
\end{corollary}
\begin{proof}

For a fixed normal inclusion $H \subset H_1$ with the quotient cyclic, we obtain an admissible pair $\tau_{\theta} = (H, \theta H)$ for any generator $\theta \in H_1/H$. The corresponding $\Utype{\tau_\theta}(\L)$ are disjoint and moreover, 
\[\Uid{\tau}(\L) = \bigsqcup_\theta \Utype{\tau_\theta}(\L) \implies \rhoid{\tau}(\L) = \sum_{\theta \text{ a generator}} \rhotype{\tau_\theta}(\L).\]
so that $2\rhoid{\tau}(\L) = \sum_{\theta \text{ a generator}} \left(\rhotype{\tau_\theta}(\L) + \rhotype{\tau^{-1}_\theta}(\L)\right)$ is a palindromic form of weight $\dim Y$ by the previous theorem and hence, so is $\rhoid{\tau}(\L)$.
\end{proof}

As the next theorem shows, the hypothesis on the resolution of singularities above are automatically satisfied at almost all primes.

\begin{thm}\label{thm: main theorem for almost all primes}
Let $M$ be a number field and $T = \Spec \mathcal O_M$. Let $f:X\to Y$ over $T$ be a generically finite, Galois map with $Y_M$ a smooth proper variety over $M$. 

Then, for almost all primes $\mathfrak p$ and every admissible pair $\tau$, the function
\[\L \to \rhotype{\tau}(\L) + \rhotype{\tau^{-1}}(\L)\]
is a palindromic form of weight $\dim Y$, as $\L$ ranges over extensions of the completion $\widehat{\mathcal O}_{M,\mathfrak p}$.
\end{thm}
\begin{proof}
We pick a characteristic zero $G$-equivariant resolution of varieties $\tilde{X} \to X$ (see \cite{abramovich1996equivariant},\cite{encinas1998good} or \cite{bierstone1997canonical}). As in Remark \ref{rmk: char 0 resolutions}, we can spread the resolution out to an open subset of $T$. The assumptions of Theorem \ref{thm: Main theorem, density fixing inertia and frobenius coset} are thus satisfied for all but finitely many primes $\mathfrak p \in T$ and the theorem follows.
\end{proof}

We also note that as the size of the residue field tends to infinity, the densities $\rhotype{\tau}(L)$ simplify drastically.

\begin{corollary}\label{cor: q to infty limiting behaviour}
We maintain the hypothesis of the above theorem. As $N(\mathfrak p) \to \infty$, we have the following limiting behaviour
    \[  \lim_{N\mathfrak p \to \infty}\frac{\rhotype{\tau}(\widehat{\mathcal O}_{M,\mathfrak p})}{(N\mathfrak p)^{\dim X}} = \begin{cases} 0 & \text{ if } |H_\tau| > 1\\
    \frac{|[g_\tau]|}{|G|} & \text{ otherwise}\end{cases}  .\]
where $G$ is the Galois group of $f$ and $[g_\tau]$ denotes the $G$-conjugacy class of $g_\tau$.
\end{corollary}
\begin{proof}
The points $P \in Y(\widehat{\mathcal O}_{M,\mathfrak p})$ where the fiber $f^{-1}(P)$ has a non trivial action of the inertia group are all contained in the locus of points whose reductions modulo $\mathfrak p$ lie in the image of the ramification locus $f(Z_f)(\mathbb F_{\mathfrak p})$. Therefore, their density is upper bounded by
\[\frac{|f(Z_f)(\mathbb F_{\mathfrak p})|}{Y(\mathbb F_{\mathfrak p})} \to 0 \text{ as } N\mathfrak p \to \infty\]
by the Lang-Weil estimates. This proves the first part of the corollary.

Similarly, for the second part, we can assume that the map $f: X \to Y$ is \'etale and Galois by excising the ramified locus since it has measure $0$ in the limit $N\mathfrak p \to \infty$. In this case, the problem reduces to the finite field version by Hensel lifting and the second claim follows from the Chebotarev density theorem over finite fields (or indeed, also from our proof of Theorem \ref{thm: incidence identity} which extends the proof of the classical Chebotarev density theorem).
\end{proof}
 
\section{A conjecture on the density of polynomials with fixed factorization type}\label{section: factorization types density}

In this section, we will prove \cite[Conjecture 1.2]{bhargava2022density} as an application of the ideas in this paper. We recall some notation from their paper first.

\subsection{Factorization types of polynomials}

A \emph{factorization type of degree $n$} is a multiset $\{f_1^{e_1}f_2^{e_2}\dots f_r^{e_r}\}$ where the $f_i,e_i$ are positive integers satisfying $\sum_i f_ie_i = n$. We allow repeats in the list of symbols $f_i^{e_i}$ but the order in which they appear does not matter. We will often omit exponents $e_i$ if $e_i = 1$. 

For an \'etale extension $L/K$ of degree $n$ over a local field $K$, we define its factorization type to be $(f_1^{e_1},\dots,f_r^{e_r})$ if a uniformizer of $K$ factors in $L$ as $P_1^{e_1}\dots P_r^{e_r}$ where the $P_i$ are primes of $L$ having residue degree $f_i$. 

We consider non-zero polynomials $h(z)$ of degree $n$ in $\A[z]$ as elements in $\A^{n+1}$ through their coordinates with the associated Haar measure $\mu_{\mathrm{Haar}}$ normalized so that the total measure is $1$. 

\begin{conjecture}[Conjecture 1.2, \cite{bhargava2022density}, generalized to a local field]\label{conj: bhargava}
Let $\sigma$ be any factorization type of degree $n$ and $K$ a local field with size of residue field $q$. Set
\begin{itemize}
    \item $\rho(n,\sigma;q) \coloneqq$ the density of polynomials $h \in K[z]$ of degree $n$ such that $L = K[z]/h(z)$ is \'etale over $K$ with factorization type $\sigma$,
    \item $\alpha(n,\sigma;q) \coloneqq$ the density of monic polynomials $h \in \A[z]$ of degree $n$ such that $L = K[z]/h(z)$ is \'etale over $K$ with factorization type $\sigma$,
    \item $\beta(n,\sigma;q) \coloneqq$ the density of polynomials $h(z) \equiv z^n \pmod{\mathfrak m_K}$ of degree $n$ such that $L = K[z]/h(z)$ is \'etale over $K$ with factorization type $\sigma$.
\end{itemize}
Then $\rho(n,\sigma;q), \alpha(n,\sigma;q)$ and $\beta(n,\sigma;q)$ are rational functions of $q$ and satisfy
\begin{equation}\label{eqn: symmetry for rho}
    \rho(n,\sigma;q^{-1}) = \rho(n,\sigma;q);
\end{equation}
\begin{equation}\label{eqn: symmetry for alpha}
    \alpha(n,\sigma;q^{-1}) = \beta(n,\sigma;q).
\end{equation}
Implicitly, we have made the claim that the densities depend on $K$ only through the size of its residue field $\mathbb F_q$.
\end{conjecture}

We will prove Equation \ref{eqn: symmetry for rho} for all factorization types $\sigma$ and all primes $q$ in the rest of this paper. It should be possible to also prove Equation \ref{eqn: symmetry for alpha} by extending our methods although we do not do so here.

Let us first relate the above conjecture to the densities appearing in the rest of this paper. To every squarefree polynomial $h(z)$ with coefficients in a local field $K$ associated to a point in $\mathbb P^n(K)$ through its coefficients, we can associate a conjugacy class of an admissible pair $[\tau_h]$ (for $S_n$) using the map 
\[f: (\mathbb P^1)^n \to (\mathbb P^1)^n/S_n = \mathbb P^n.\]
If we further fix an ordering of the roots of $h(z)$, that corresponds to fixing an element in the conjugacy class $[\tau_h]$. For an admissible pair $\tau$, recall that
\[\Utype{\tau}(K) = \{h(z) \in \mathbb P^n(K) \text{ squarefree with } \tau \in [\tau_h]\}\]
with density $\rhotype{\tau}(K) = \mu_{\mathbb P^n}(\Utype{\tau})(K)$. The Haar measure on $\A^{n+1}$ relates to the canonical measure on $\mathbb P^n(\A)$ as follows.

\begin{lemma}\label{lemma: bhargava densities to our densities}
Consider the canonical quotient map
\[\pi: \A^{n+1} \setminus \{0\} \to \mathbb P^n(\A).\]
For a measurable set $S \subset \mathbb P^n(\A)$, we have the equality of measures
\[\mu_{\mathrm{Haar}}(\pi^{-1}(S)) = \frac{\mu_{\mathbb P^n}(S)}{|\mathbb P^n(\mathbb F_q)|}\]
where $q$ is the size of the residue field of $K$.
\end{lemma}
\begin{proof}
In brief, the canonical top form on $\A^{n+1}$ integrated over a $\mathbb G_m$ orbit (acting diagonally) gives a gauge form on $\mathbb P^n(\A)$. Therefore, the measures of $S$ and $\pi^{-1}(S)$ are equal to a global normalization. More explicitly:

Consider local coordinates where $v_p(a_0) \leq v_p(a_i)$ for all $0 < i \leq n$. On this chart, we can realize $\pi$ as a projection 
\[\pi: \A^{n+1} \setminus \{0\} \to \A^n; (a_0,\dots,a_n) \to (a_1/a_0,\dots,a_n/a_0)\]
onto the final $n$ co-ordinates. From this description, one computes that the preimage of a set $S$ has measure exactly
\[\mu_{\mathrm{Haar}}(S) = \frac{\mu_{\mathbb P^n}(S)}{q^{n}}\left(\frac{q-1}{q} + \frac{q-1}{q^{n+2}} + \dots\right) = \frac{\mu_{\mathbb P^n}(S)(q-1)}{q^{n+1}-1}\]
as required.
\end{proof}

In the other direction, to each admissible pair $\tau$ as above (with respect to $S_n$), we can associate a factorization type $s(\tau)$ such that if $\tau \in [\tau_h]$ for some squarefree polynomial $h(z)$, then $s(\tau)$ is the factorization type for $h(z)$:

\begin{definition}\label{defn: factorization types from admissible pairs}[Factorization types from admissible pairs]
Recall that an admissible pair $\tau$ has associated to it a normal inclusion of subgroups of $S_n$: $H_\tau \subset H_{1,\tau}$ where the smaller group corresponds to the inertia group and the bigger group to the total Galois group. The action of $H_{1,\tau}$ on the set $[n] = \{1,\dots,n\}$ will partition $[n]$ into orbits $\pi_1,\dots,\pi_r$ and we define $H_{1,\tau}^{(i)}$ to be the projection of $H_{1,\tau}$ in the symmetric group $S_{\pi_i} = \operatorname{Sym}(\pi_i)$ and similarly for $H_\tau^{(i)}$. Furthermore, the action of $H_{\tau}^{(i)}$ will partition $\pi_i$ into further orbits $\pi_{i,j}$. Since $H_\tau$ is normal in $H_{1,\tau}$, the $\pi_{i,j}$ all have the same size that we denote by $e_i$ and the number of such orbits will be denoted by $f_i$. Given this data, the factorization type $s(\tau)$ is defined to be the multiset with $r$ elements $f_i^{e_i}$ corresponding to the parts  $\pi_1,\dots,\pi_r$.
    
\end{definition}


Note that $s(\tau)$ only depends on the conjugacy class of $\tau$. We also note in passing that the coset $g_\tau$ cyclically permutes the partitions $\pi_{i,j}$ for $j=1,\dots,f_i$. The next lemma relates the factorization type densities to the splitting type densities from earlier in the paper and along the way, shows that $s([\tau_h]) = \sigma_h$ as claimed above.

\begin{lemma}\label{lemma: relation of bhargava conjecture to our theorem}
Let $\sigma = \{f_1^{e_1},\dots,f_r^{e_r}\}$ be a factorization type. Then, we have the identity
\[\rho(n,\sigma;q) = \frac{1}{|\mathbb P^n(\mathbb F_q)|}\sum_{\substack{\tau \text{ up to } S_n \text{ conjugacy}\\ \text{satisfying } s(\tau) = \sigma }}\rhotype{\tau}(K).\]
\end{lemma}
\begin{proof}\label{pf: bhargava conjecture}

We will show that the polynomials contributing to the density calculation in the right hand side correspond precisely to the polynomials contributing to the density calculation on the left hand side. This will follow if we show that a polynomial $h(z)$ with admissible pair $[\tau_h]$ satisfies $s([\tau_h]) = \sigma_h$. Since we are taking admissible pairs up to conjugacy on the right, this will prove exactly the bijective correspondence we need. The normalizing factor $|\mathbb P^n(\mathbb F_q)|^{-1}$ follows from Lemma \ref{lemma: bhargava densities to our densities}.

To this end, let $h(z) \in K(z)$ be a squarefree degree $n$ polynomial with factorization type $\sigma$. Upon fixing an ordering of its roots, we obtain an admissible pair $\tau_h$. The corresponding decomposition of $[n]$ into parts $\pi_1,\dots,\pi_r$ as in Definition \ref{defn: factorization types from admissible pairs} above corresponds exactly to the orbits of $\Gal(\overline{K}/K)$ on the set of roots and hence to the decomposition of $h(z) = h_1(z)\dots h_r(z)$ into irreducible polynomials over $K$. Moreover, $H_{1,\tau}^{(i)}, H_{\tau}^{(i)}$ correspond respectively to the Galois group, inertia group of the splitting field $L_i/K$ of $h_i(z)$ and it follows that $f_i,e_i$  are respectively the residue degree, inertia degree of the extension $K \subset K(z)/h_i(z)$.

\end{proof}

\begin{thm}\label{thm: symmetry follows from rationality for factorization densities}
    For a factorization type $\sigma$, let $\mathcal L_\sigma$ be the set of local fields for which the characteristic of the residue field is co-prime to the exponents $e_i$ occurring in $\sigma$. 
    
    For this set of local fields, $\rho(n,\sigma;q)$ is equal to a rational function $r_\sigma(t)$ evaluated at $t = q$, the size of the residue field. This rational function takes values in $\mathbb Q$ and satisfies the symmetry
    \[r_\sigma(t^{-1}) = r_\sigma(t).\]
\end{thm}
\begin{proof}
    By \cite{john}, \cite{del_corso_dvornicich_2000}, $\rho(n,\sigma;q)$ is known to be a rational function for the set of local fields considered in the theorem. It thus suffices to prove the above symmetry after evaluating $t$ at infinitely many values $q$ coming from $\mathcal L_\sigma$  since rational functions are determined by their values at any large enough set of points. 
    
    In particular, we can fix an arbitrarily large $p$ and consider all the extensions of $\mathbb Z_p$. For this set of extensions, we know by Theorem \ref{thm: main theorem for almost all primes} and Lemma \ref{lemma: relation of bhargava conjecture to our theorem} that $\rho(n,\sigma;p)$ is a palindromic form (over $\mathbb Q$) of weight $1$ since $\rhotype{\tau}(\F)$ and $|\mathbb P^n(\mathbb F_q)|$ are palindromic forms of weight $n$ so that their ratio is a palindromic form of weight $1$. In other words, this shows that
    \[r_\sigma(q^{-1}) = r_\sigma(q)\]
    as $q$ ranges over the powers of the prime $p$, thus completing the proof.
\end{proof}

The above theorem very quickly deals with ``tame" primes as a combination of earlier results and the general theorems in this paper. Unfortunately, this is the limit of this direct application of our general theorems in this context since the splitting densities $\rhotype{\tau}(K)$ fail to be rational functions in general - they are only rational along arithmetic progressions. 

In the final part of this section, we relate the factorization densities more directly to certain integrals over a particularly nice class of quotients. 

\subsection{Reducing the computation of factorization densities to certain integrals}\label{sec: relating factorization densities to integrals}

For a factorization type $\sigma = \{f_1^{e_1},\dots,f_r^{e_r}\}$ of degree $n$, we define some associated notions. We pick a partition $\mathcal P_\sigma$ of $[n] = \{1,\dots,n\}$ with distinct parts $\pi^{\sigma}_{i,j}$ of size $e_i$ for $1\leq i\leq r$ and $1 \leq j \leq f_i$. If required, we will express the $f_i,e_i,r$ by $f_i(\sigma),e_i(\sigma),r(\sigma)$ to make the dependence on $\sigma$ clear. Note that this partition is well defined up-to $S_n$ conjugacy and indeed, all the constructions here will be well defined precisely up to this relabelling.

We next define an associated admissible pair $\tau_\sigma = (H_\sigma,g_\sigma H_\sigma)$ where $H_\sigma$ is the maximal subgroup of $S_n$ preserving the partition $\mathcal P_\sigma$, i.e., $H_\sigma \cong \prod_{i=1}^r(S_{e_i}^{f_i})$ while the coset $g_\sigma H_\sigma$ cyclically permutes the cosets in the order $\pi^\sigma_{i,1},\dots,\pi^{\sigma}_{i,f_i}$ for each $i$.

Since everything is only well defined up to conjugacy, we say that $A \lesssim B$ if $A$ is less than some $S_n$-conjugate of $B$. The mapping $\tau \to s(\tau)$ is adjoint to $\sigma \to \tau_\sigma$ in the following sense. 

\begin{lemma}\label{lemma: adjointness between pairs and types}
For any admissible pair $\tau$ for $S_n$ and factorization type $\sigma$, $\tau \lesssim \tau_\sigma$ if and only if $\tau_{s(\tau)} \lesssim \tau_{\sigma}$.
\end{lemma}
\begin{proof}
The admissible pair $\tau$ induces a partition $\mathcal P_\tau$ of $[n]$ corresponding to the orbits of $H_\tau$ on $[n]$ as in Definition \ref{defn: factorization types from admissible pairs}. The inequality $\tau \lesssim \tau_\sigma$ is equivalent to there being a coarsening of $\mathcal P_\tau$ to the partition $\mathcal P_\sigma$ such that the coset $g_\tau H_\tau$ induces the same permutation as $g_\sigma H_\sigma$. Since $\tau$ and $\tau_{s(\tau)}$ induce the same partition of $[n]$ and the same induced cyclic permutation on the set of cosets, the lemma follows. 
\end{proof}

We define a partial order on the set of factorization types by 
\[\sigma' \leq \sigma \iff \tau_{s(\sigma')} \lesssim \tau_{s(\sigma)}.\]
We can then rephrase the above lemma as
\[\tau \lesssim \tau_\sigma \iff s(\tau) \leq \sigma.\]

We need one more lemma before the main theorems of this section.

\begin{lemma}\label{lemma: alpha invariant under adjointness}
The invariant $\alpha$ from definition \ref{defn: alpha, beta} satisfy the following identity
\[\alpha(\tau,\tau_\sigma) = \alpha(\tau_{s(\tau)},\tau_\sigma).\]
\end{lemma}
\begin{proof}
By definition, $\alpha(\tau,\tau_\sigma)$ is the number of elements in $H_\sigma\backslash S_n$ which induce a permutation of the partition $\mathcal P_\tau$ while preserving the properties that $\mathcal P_\tau$ is a refinement of $\mathcal P_\sigma$ and the induced cyclic permutation of $g_\tau$ on $\mathcal P_\sigma$ equals the induced cyclic permutation of $g_\sigma$ on $\mathcal P_\sigma$. Since $\tau$ and $\tau_{s(\tau)}$ induce the same partition of $[n]$ and the same induced cyclic permutation, the lemma follows.
\end{proof}

For each factorization type $\sigma$, we denote a resolution of the quotient map by $f_\sigma: X_\sigma \to {^{g_\sigma}}\left((\mathbb P^1)^n/H_{\sigma}\right) \to \mathbb P^n$ where $X_\sigma$ is a nice resolution of singularities that we assume exists (as in the rest of this paper).

\begin{thm}
For any factorization type $\sigma$ of degree $n$, we have the identity
\[\eta_{f_\sigma}(K) = \sum_{\substack{\tau : s(\tau) \leq \sigma\\\text{up to } S_n \text{ conjugacy}}}\alpha(\tau,\tau_\sigma)\rho_{f_\sigma,\tau}(K) = \sum_{\sigma'\leq \sigma}\alpha(\tau_{\sigma'},\tau_{\sigma})|\mathbb P^n(\mathbb F_p)|\rho(n,\sigma;p).\]
\end{thm}
\begin{proof}
The first identity follows immediately from Theorem \ref{thm: incidence identity} and Lemma \ref{lemma: adjointness between pairs and types}. The second identity follows from Lemma \ref{lemma: alpha invariant under adjointness} and Lemma \ref{lemma: relation of bhargava conjecture to our theorem}.
\end{proof}

Finally, we note as before that since $\alpha(\tau_{\sigma'},\tau_{\sigma}) > 0$ for all $\sigma' \leq \sigma$, we can M\"obius invert the above system of identities with respect to the poset of factorization types to obtain
\begin{equation}\label{eqn: factorization densities as integrals}
    \rho(n,\sigma;q) = \frac{1}{|\mathbb P^n(\mathbb F_q)|}\sum_{\sigma'\leq \sigma}\alpha^{-1}(\tau_{\sigma'},\tau_{\sigma})\eta_{f_{\sigma'}}(K)
\end{equation}
where $\alpha^{-1}$ is the inverse of $\alpha$ in the incidence poset of factorization types. We have thus shown that the conjecture on factorization densities is equivalent to showing that the integrals $\eta_{f_\sigma}(K)/|\mathbb P^n(\mathbb F_q)|$ are rational functions in $q$ invariant under the transformation $q \to q^{-1}$. 


We consider $\mathbb P^n$ as the parameter space for degree at most $n$ univariate polynomials as usual so that $\prod_{i=1}^r\mathbb P^{n_i}$ parametrizes tuples $(h_1,\dots,h_r)$ of polynomials. We have reduced the proof of rationality to finding a resolution for the complement of the \emph{resultant locus} $\mathcal R \subset \mathbb \prod_{i=1}^rP^{n_i}$ corresponding to the locus where two of polynomials $h_i,h_j$ share a common root. Crucially, we want such a compactification over $\Spec \mathbb Z$ (or at least over an open cover).

This problem is closely analogous to the problem of finding a resolution of the discriminant locus $\mathcal D \subset \mathbb P^n$ corresponding to polynomials $h$ having repeated roots and the resolution we describe in the next section also applies to the discriminant locus.

\section{Constructing a resolution of the resultant locus}\label{sec: resolving resultant locus}

In this section, we describe a procedure to resolve the resultant locus and more generally, to resolve \emph{skew-conically} stratified varieties which will include as a special case the discriminant locus. Our construction will be a modification of the resolution described in \cite{macpherson1998making} and as such, we will explain their paper in brief first.

\subsection{Making conical stratifications wonderful}

Their paper deals with smooth varieties $X$ with a conical stratification. They work over $\mathbb C$ with analytic neighbourhoods but in fact, their methods can be easily adapted to work in arbitrary generality with (formally) completed algebraic neighbourhoods. Throughout this subsection, we let $(X,S_{\alpha} : \alpha \in I)$ be a smooth variety with the $S_\alpha \subset X$ locally closed strata, i.e., the $S_\alpha$ are disjoint  with $\bigcup_\alpha S_\alpha = X$ and the following relation defines a partial order on the set $I$: 
$$\alpha \leq \beta \iff \overline{S}_\alpha \subset \overline{S}_\beta.$$ 

\begin{definition}[cone]
    Let $(D,S_\alpha)$ be a formal analytic neighbourhood of a smooth, stratified space $(X,S_\alpha)$. It is said to be a cone if the origin is a closed strata and moreover, the strata $S_\alpha$ are stable under the standard diagonal action of the torus $\mathbb G_m$ on the formal disc $D$. This is equivalent to the data of a stratification on the projective space defined by the quotient of $D-0$ by $\mathbb G_m$.
\end{definition}

\begin{definition}[conical stratification]
    $(X,S_\alpha)$ as above is said to be conically stratified if for every point $x \in X$, the formal analytic neighbourhood $D_x$ around $x$ can be decomposed as $D_x \cong D_T \times D_N$ where the tangent disc $D_T$ is the formal completion \emph{inside} the strata $S_\alpha \ni x$ and has the trivial stratification (i.e., the origin is the only proper strata) and moreover, the normal disc $D_N$ is stratified as a cone.
\end{definition}

% and finally,

\begin{definition}[wonderful compactification]
    $(X,S_\alpha)$ as above is said to be a wonderful compactification if the $S_\alpha$ arise as intersections of (geometric) snc divisors.
\end{definition}

We note that constructing a wonderful compactification is exactly our goal in finding a good resolution for our applications earlier in the paper.

We now explain the methods of \cite{macpherson1998making} where they show that any conically stratified variety $(X,S_\alpha)$ has a resolution $(\tilde{X},\tilde{S}_\alpha) \to (X,S_\alpha)$ where $(\tilde{X},\tilde{S}_\alpha)$ is a wonderful compactification. They identify the following special \emph{minimal irreducible} strata:

\begin{definition}[minimal irreducible strata]
    Let $(X,S_\alpha)$ be as above. We say a strata $S_\alpha$ is reducible if for every point $x \in S_\alpha$, the local decomposition $D_x \cong D_T \times D_N$ has the property that the cone $(D_N,S_\alpha \cap D_N)$ can be further factored into stratified cones as 
    \[ D_N \cong D_{N_1} \times \dots \times D_{N_r} \hspace{10 mm} (r \geq 2);\]
    where the isomorphisms preserve the stratifications. If such a non-trivial decomposition is not possible, $S_\alpha$ is said to be irreducible. Moreover, if every strata $S_\beta$ with $S_\beta \subset \overline{S}_\alpha$ is reducible, then $S_\alpha$ is said to be a minimal irreducible strata.
\end{definition}

The strategy for producing a resolution is now simple: We blow up the closures of the minimal irreducible strata of codimension greater than $1$ (in any order) until there are none left, at which point we have a wonderful compactification. Three main things have to be checked for this to work:
\begin{lst}\label{item: claims to check for resolution}
\begin{enumerate}
    \item The closures of the minimal irreducible strata are smooth.
    \item The blowup at each minimal irreducible strata has an induced conical stratification.
    \item After blowing up, the number of minimal irreducible strata of codimension $\geq 2$ decreases.
\end{enumerate}
\end{lst}

The first claim is checked in \cite[\S 2.3]{macpherson1998making} and relies on formal properties of stratification. The proof goes through just as well with formal analytic neighbourhoods instead of complex balls. 

The argument \cite[\S 2.2]{macpherson1998making} for the second claim reduces to showing that the blowup of a conically stratified formal neighbourhood is also conical. This relies on locally realizing the blowup as the tautological bundle over projective space where this base projective space inherits the induced stratification from the quotient of the cone by $\mathbb G_m$.

The third claim follows from \cite[\S 2.2, Lemma]{macpherson1998making} - the irreducible strata of the blow-up can be explicitly characterized.

There are a few more things to be checked (such as the fact that a conical stratification is wonderful if and only if all the minimal irreducible strata are of codimension $1$) but these easily adapt to our setting.

\subsection{Skew-conical compactifications}

One key example of a conical compactification comes from configuration space on a smooth variety $X$, such as $((\mathbb P^1)^n,S_\pi)$ where the strata are indexed by the partition $\pi$ corresponding to the indices of equal co-ordinates; the largest open strata is configuration space. 

This is very close to the space we need to resolve as described in the previous section - we instead need to consider the quotient of $(\mathbb P^1)^n$ by $\prod_{i=1}^r S_{n_i}$ and unfortunately, the quotient does not have a conical stratification - only a \emph{skew-conical stratification} as defined below.

In this section, we will need to work with Artin stacks for technical reasons and throughout, will work over the base ring $\mathbb Z$. The Artin stacks that appear will be global quotients of the form $\mathscr Y = Y/T$ for $Y$ a smooth scheme and $T$ a torus acting with finite stabilizers and we will denote such stacks by a \emph{space} in this subsection. A stratification of such a space will be a $T$ equivariant stratification $S_\alpha$ of $Y$.

The only technical point is to define the appropriate notion of a formal neighbourhood on an Artin stack. Here, we take the easy way out since our Artin stacks $\mathscr Y$ are in fact global quotient stacks. Given a point $y \in \mathscr Y$, we consider a neighbourhood of it to be the tubular neighbourhood of the preimage of $y$ in $Y$. Since the stratification on $Y$ is $T$ equivariant, the decompositions into tangential and normal formal discs can be carried out as before.


\begin{definition}[skew-cones]
    We say that a formal neighbourhood $(D,S_\alpha)$ of a stratified space $\mathscr Y$ is a skew-cone if $0$ is a strata and moreover, the strata $S_\alpha$ are stable for some nontrivial $\mathbb G_m$ action on $D$: in other words, we can find co-ordinates $x_1,\dots,x_n$ with $\mathbb G_m$ acting by $t \cdot x_i = t^{n_i}x_i$ with all the $n_i > 0$.
\end{definition}

\begin{definition}[skew-conical stratification]
    A space $(\mathcal Y,S_\alpha)$ is said to be skew-conically stratified if around every point $x \in \mathcal Y$, we have a decomposition of the formal neighbourhood $D_x \cong D_T \times D_N$ as before where we now require the normal direction $D_N$ to be a \emph{skew} cone.
\end{definition}

We will explain how to resolve such skew-conical stratifications by modifying the methods of \cite{macpherson1998making} in the remainder of this section. We will not provide as many details in this section since they consist of modifications to the arguments of \cite{macpherson1998making} and checking that these modifications do work. We explain the modifications needed and leave the checking to the reader since the proofs are almost unchanged.

First, we discuss our prototypical example:

\begin{example}
    The central example for us will be $X = \prod_{i=1}^r\mathbb P^{n_i}$ with strata $S_\alpha$ where $\alpha$ is a partition of $\sum_i n_i$, as in the configuration space example. Indeed, these strata are the quotients of partition strata $S_\pi$ appearing in the configuration space stratification on $\mathbb P^{\sum n_i}$ described above. 
\end{example}

In order to see that this is a skew-conical compactification, consider the quotient map
\[f: Y = \mathbb P^{\sum n_i} \to X = \prod_{i=1}^r\mathbb P^{n_i}.\]
Let $x \in S_\alpha \subset X$ be a point with local formal neighbourhood $D_x$. Let $y \in T_\pi = f^{-1}(S_\alpha)$ map to $x \in S_\alpha$. For $G \subset \prod_{i=1}^rS_{n_i}$ the isotropy group of $y$ (or equivalently of $T_\pi$), we pick a decomposition $D_y \cong \tilde D_T \times \tilde D_N$ such that $G$ acts on both $\tilde D_T$ and $\tilde D_N$ with the isomorphism equivariant for this action. Since $G$ fixes $T_\pi$ pointwise, it acts trivially on $\tilde D_T$. Moreover, the co-ordinates of $\tilde D_N$ are linear combinations, and in fact differences, of the ``standard" co-ordinates on $Y$ so that $G$ acts linearly on these co-ordinates.

This induces an isomorphism $D_x \cong \tilde{D}_T/G\times \tilde D_N/G$. Since $G$ acts trivially on $\tilde D_T$, $\tilde D_T/G = \tilde D_T$. On the other hand, since $G$ acts linearly on the co-ordinates of $\tilde D_N$, one can find a system of generators for functions on $\tilde D_N/G$ that are monomials in the co-ordinates on $\tilde D_N$. In particular, this implies that the diagonal $\mathbb G_m$ action on $\tilde D_N$ descends to a skew $\mathbb G_m$ action on $\tilde D_N/G$ (by acting on these monomial co-ordinates).

\begin{remark}\label{rmk: resolving discriminant locus}
    In the case where $r=1$ so that $X= \mathbb P^n$, the largest closed strata corresponds to the discriminant locus corresponding to degree $n$ polynomials with multiple roots. The remainder of this section will in particular, show a method to resolve this discriminant locus. In general, one obtains a resolution of the ``resultant locus."
\end{remark}

For technical reasons, the resolution passes through Artin stacks in the intermediate steps and we obtain an algebraic space in the final step. The earlier results of the paper continue to work in the settings of algebraic spaces since $p$-adic integration is perfectly well defined in this context (see \cite[\S2.1]{groechenig2020geometric} ). In more detail, the problem with modifying the argument in \cite{macpherson1998making} is in step two of \ref{item: claims to check for resolution}. Since the $\mathbb G_m$ action is skew, the standard blow-up will not inherit this skew action. The modification is simple - we use a weighted blow-up instead, where the weights are determined by the characters of the $\mathbb G_m$ action. As before, a weighted blow-up is a torsor over weighted projective space which is sufficient to show that the action can be extended to the blow up.

Unfortunately, we have the technical complication that weighted blow-ups are never smooth for non-trivial weights. In order to fix this, we move to the world of stacks where indeed, the weighted blow up stack is always smooth. The rest of the argument in \cite{macpherson1998making} goes through verbatim to produce a smooth Artin stack $(\tilde{\mathscr X},\tilde S_\alpha)$ and a resolution $\tilde{X} \to X$ respecting the stratifications.


Finally, we \emph{destackify} the above resolution to produce a honest resolution by smooth algebraic spaces using the recent paper \cite{bergh2019functorial}.

\begin{thm}[resolving skew-conically stratified spaces]
    Let $(X,S_\alpha)$ be a skew-conically stratified space. Then there exists a resolution $(\tilde{X},\tilde{S}_\alpha) \to (X,S_\alpha)$ by an algebraic space where $\tilde S_\alpha$ is a wonderful compactification.
\end{thm}
\begin{proof}
    After the discussion above, we have a found an Artin stack $\tilde{\mathscr X}$ with a wonderful stratification $\tilde{S_\alpha}$ and a resolution $\tilde{\mathscr X} \to X$. Note that since $\tilde{X}$ is obtained from taking weighted blowups, the stabilizers are all subgroups of $\mathbb G_m^s$ for some $s$ and hence the Artin stack $\tilde X$ is tame.

    Now, \cite[Theorem B]{bergh2019functorial} can be used to define an algebraic space $\tilde{X}$ obtained as the coarse space of a sequence of stacky modifications of $\tilde{\mathscr X}$ which provides a resolution $\tilde X \to X$ (so that the preimage of the stratification $S_\alpha$ is wonderful).
\end{proof}

This completes the proof of resolutions of the resultant stratification. The final thing left to show is that the strata appearing in the resolution and their twists have polynomial point counts.

\begin{thm}\label{thm: rationality of wild-type densities}
    For every factorization type $\sigma$ of degree $n$, there exists a rational function $r_\sigma(t)$ such that for \emph{every} prime power $q$, we have
        \[\rho(n,\sigma;q) = r_\sigma(q).\]
    Moreover, the rational function $r_\sigma$ satisfies the functional equation
        \[r_\sigma(t^{-1}) = r_\sigma(t).\]
\end{thm}
\begin{proof}
    By Equation \ref{eqn: factorization densities as integrals}, it suffices to show that the integrals $\eta_{f,\sigma}(q)$ form a rational function as $q$ ranges over prime powers for all factorization types $\sigma$. As in \S \ref{sec: relating factorization densities to integrals}, let $\tau_\sigma$ be the admissible type associated to the factorization type $\sigma$ and consider the the twist quotient map
        \[ f_\sigma: {^{g_\sigma}(\mathbb P^1)^n} / H_\sigma. \]
    As discussed in this section so far, we can find a resolution $X \to (\mathbb P^1)^n/H_\sigma$ by iterated stacky blowups and modifications. In fact, one can check that this resolution sequence is equivariant for the $g_\sigma$ action on $(\mathbb P^1)^n/H_\sigma$ so that we can equivariantly twist the entire construction to obtain a resolution
        \[ \pi_\sigma: X_\sigma \to {^{g_\sigma}(\mathbb P^1)^n} / H_\sigma.\]
    By the proof of Theorem \ref{thm: symmetry for pullback}, it suffices to show that the strata of $X_\sigma$ have polynomial point counts for \emph{all} prime powers $q$ (where we twist \emph{after} base changing to $\mathbb F_q$ as usual). By inclusion-exclusion, we may further reduce to proving that the closed strata have polynomial point counts.
    
    First, we have the well known lemma that if $Z \subset X$ are such that both $Z$ and $X$ have polynomial point counts and $Z,X$ are smooth, then so does the blow-up of $X$ along $Z$. This can be seen in multiple ways but perhaps the easiest is that in the Grothendieck ring of varieties (or stacks), we have the equation
        \[ [Bl_Z(X)] = [X]- [Z] + [Z]\times[\mathbb P^{\dim X-\dim Z-1}].\]

    Thus, the above lemma in combination with the fact that our compactification is obtained from a sequence of blowups shows that it is enough to prove that the closed strata of $^{g_{\sigma}}(\bbP^1)^n/H_{\sigma}$ all have polynomial point counts for all prime powers $q$. 
    
    The proof will be by an induction with respect to the resolution process. By the Grothendieck-Lefschetz trace formula, it suffices to show that the trace of the Frobenius on the cohomologies of the strata are all of the form $q^i$ for $i \in \mathbb N$. The closed strata of $(\mathbb P^1)^n / H_\sigma$ are quotients of the closed strata of $(\mathbb P^1)^n$ which are all of the form $(\mathbb P^1)^r$ for $r \leq n$ and hence the Frobenius acts by a scalar on the appropriate cohomology groups. Since the cohomology of a quotient injects into the cohomology of the cover, this shows that the closed strata of $(\mathbb P^1)^n / H_\sigma$ have polynomial point counts for all prime powers $q$.

    Now we have to show the same for the twist by $g_{\sigma}$. The action of the Frobenius on the twist $^{g_\sigma}[(\mathbb P^1)^n] / H_\sigma$ is equivalent to the action of $g_\sigma$ times the Frobenius on $[(\mathbb P^1)^n] / H_\sigma$. Since the Frobenius acts on the cohomology by a scalar and $g_\sigma$ by a permutation matrix defined independently of $q$, the trace of the Frobenius on the cohomology of the twist is of the form $c_iq^i$ for $c_i,i$ independent of $q$. This proves that the strata of the twist also have polynomial point counts for all $q$.

\end{proof}

\begin{remark}
    We can alternatively prove the above theorem by leveraging the already known tame case: Since our main theorems express the factorization densities as a rational function of certain motives which are smooth and proper over $\operatorname{Spec}\mathbb Z$ (due to the existence of an integral resolution as above), it suffices to prove that the point counts of this rational function in the motives is polynomial. We can complete the proof at this point by using a comparison theorem to Hodge theory in order to show that the motives are of Hodge-Tate type, exactly as in \cite[\S Appendix]{hausel2008mixed}. 
\end{remark}


% \subsection{Tate type resolutions}\label{section: tate type resolutions}

% As remarked above, we need to show that the integrals $\eta_{f_\sigma}(L)/|\mathbb P^n(\mathbb F_q)|$ are rational functions in $q$ invariant under the transformation $q \to q^{-1}$. By Theorem \ref{thm: symmetry for pullback}, it suffices to show this by constructing the resolutions $f_\sigma: X_\sigma \to \mathbb P^n$ alluded to above such that the ramification divisor is geometrically snc and the intersections of the components of the exceptional divisor are all \emph{Tate type}.

% \begin{definition}\label{defn: Tate type}
% We say that a variety $X \to \mathbb Z$ is of \emph{Tate type} if the point counts $|X(\mathbb F_q)|$ are a polynomial as $q$ ranges over the powers of primes.
% \end{definition}

% \begin{remark}\label{rmk: galois action of tate type varieties}
%     If $X$ is smooth, projective and of \emph{Tate type}, then the Frobenius $\mathrm{Frob}_p$ acts on the (semi-simplification of the) cohomology $H^{2i}_{\text{\'et}}(\overline{X}_{\mathbb F_p},\mathbb Z_\ell)$ as multiplication by $p^i$. Briefly, the idea is that we can use the Lefschetz trace formula to identify the action of the Frobenius on the entire cohomology ring through the point counts. Then, purity of the eigenvalues (by the Weil conjectures) identifies the action of the Frobenius on the individual cohomology groups. As before (Lemma \ref{lemma: point counts have right weight}), the symmetry $q \to q^{-1}$ is a consequence of Poincar\'e duality.
% \end{remark}

% Explicitly, we are trying to resolve the ramification locus of the \emph{multiplication map}
% \[\prod_{i=1}^r\mathbb P^{n_i} \to \mathbb P^N: (f_1,\dots,f_r) \to f_1\dots f_r\]
% where the $f_i$ are univariate polynomials (in the variable $z$) of degree at most $n_i$ and we identify projective space $\mathbb P^{n_i}$ as the moduli space of such polynomials. As is well known, the map is ramified precisely at those points where some of the $f_i$ have common roots, i.e., we are interested in finding a good resolution of the \emph{resultant locus}.

% In order to prove the conjecture, we would like to find such a resolution of singularities over all the local rings $\mathbb Z_p$. In fact, we have a good candidate for such a resolution and the only difficulty is in showing that this candidate works for all primes. Since we have a universal candidate, it suffices to show that this candidate is good (i.e., is a sncd) over $\mathbb F_p$ for every prime $p$.

% We will prove this using the methods of ``conical stratification'' \cite{macpherson1998making}, adapted to arbitrary characteristic. While that paper works in the complex analytic setting, the proofs in fact extend to arbitrary characteristic replacing "complex analytic neighbourhoods" by "analytic neighbourhoods" as we explain below. For technical reasons, we will also need to work with compactifications by fairly mild stacks instead of smooth projective schemes. Nevertheless, the $\A$ points of this compactification $\tilde{X}$ still have the structure of a $\F$-analytic manifold so that our main results using p-adic integration still go through. Because of this technical complexity and reference on background results, we will suffice ourselves with sketching a proof by explaining the modifications that need to be done to prior results in the literature. We emphasize again that the conjecture on factorization densities has already been prove for all primes $p$ co-prime to the exponents appearing in the factorization type $\sigma$ in the prior sections. 

% \yif{Weighted conical stratification is necessary to get (analytic) sncd on the smooth locus. }

% \input{asvin_stacky_blowup}
% % \input{yifan_log_pair_integraion}

% \subsubsection{Weighted conical stratification}
% By a stratification $\{S_\alpha\}_{\alpha \in I}$ of a space $X$, we mean that the $S_\alpha$ are Zariski locally closed subsets of $X$ such that the closure of any strata $S_\alpha$ is a union of other strata $S_\beta$. This gives the strata $S_\alpha$ (and $I$ itself) a partial order structure defined by
% \[ \alpha \geq \beta \iff S_\alpha \geq S_\beta \iff \overline{S}_\alpha \supset S_\beta.\]
% For the rest of this sub-section, we fix such a space $X$ and a stratification as above.

% \begin{example}[Wonderful stratification]\label{example: divisor stratification}
%     Given a geometric sncd $D = \bigcup_{i\in I} D_i \subset X$, we can define a stratification where the strata are $S_J = \bigcap_{j \in J}D_i$ as $J$ ranges over subsets of $I$. A stratification is wonderful if it is isomorphic to (a coarsening of) this one.
% \end{example}

% \begin{example}\label{example: pullback stratification}
% Given a map $X\to Y$ and a stratification $\{S_\alpha\}_{\alpha \in I}$ of $Y$, the preimages $\{\pi^{-1}(S_\alpha)\}_{\alpha \in I}$ form a stratification of $X$ and induce the same partial order on $I$ as $\{S_\alpha\}_{\alpha \in I}$.
% \end{example}

% \begin{example}\label{example: product stratification}
%     Given stratifications $\{S_\alpha\}_{\alpha \in I}, \{T_\beta\}_{\beta \in J}$ of $X,Y$, the product $X\times Y$ can be stratified by $\{S_\alpha\times T_\beta\}_{\alpha,\beta \in I,J}$. In this case, $(\alpha,\beta) \leq (\alpha',\beta') \iff \alpha \leq \alpha'$ and $\beta \leq \beta'$.
    
% \end{example}

% The key notion is that of a \emph{weighted conical stratification}. We define an analytic $n$-dimensional disk $D^n(\vec{d})$ as the stacky completed local ring of weighted projective space $\mathbb P(\vec{d})$ as above. Explicitly, they are of the form $[\Spec k[\![x_1,\dots,x_n]\!]/\mu_d]$ where $\mu_d$ acts by the $t \to t^{d_i}$ character in the $i$-th co-ordinate. When the $d_i$ are all $1$, this is simply isomorphic to $\Spec k[\![x_1,\dots,x_n]\!]$ and we denote it by $D^n$. 

% \begin{definition}[Stratification of a disk as a weighted cone]
%     Let $\{S_\alpha\}_{\alpha \in I}$ be a stratification of $D_{n}(1)$ induced from a stratification of $\mathbb A^n$ where we think of $D_n(1)$ as the analytification of the origin. We say that this stratifies $D_n(1)$ as a weighted cone if there exist some positive integers $(d_1,\dots,d_n)$ (with gcd $1$) such that each strata $S_\alpha$ is preserved under the $\mathbb G_m$ action
%     \[t(x_1,\dots,x_n) \xrightarrow{} (t^{d_1}x_1,\dots,t^{d_n}x_n).\]
%     The tuple $\vec{d} = (d_1,\dots,d_n)$ is called the weight of the stratification.
% \end{definition}

% \begin{definition}[A conical stratification]
%     We say that a stratification $\{S_\alpha\}_{\alpha\in I}$ of $X$ is conical if for every $x \in S_\alpha \subset X$,
%     \begin{itemize}
%         \item there is a product decomposition of the completion $\widehat{\mathcal O}_{X,x} \cong D_T \times D_N$ mapping the origin to the origin (where $D_T$ is isomorphic to a product of $ D_{k}(d)$ for some $k \geq 1,\vec{d}$ and $D_N \cong D_{\dim X - dim S_\alpha}(1)$),
%         \item and the restriction of the stratification $S_\alpha$ to $\widehat{\mathcal O}_{X,x}$ is induced by the trivial stratification on $D_T$\footnote{the strata are the origin and the complement of the origin } and a stratification of $D_N$ as a weighted cone.
%     \end{itemize} 
% \end{definition}

% One should think of $D_T$ as an analytic disc tangent to the strata containing $x$ and $D_N$ as an analytic disc normal to this strata. 

% \begin{remark}
%  Giving a conical stratification of $D_n(1)$ as a cone with weight $(d_1,\dots,d_n)$ is equivalent to a conical stratification of the weighted projective stack $\mathbb P(d_1,\dots,d_n)$: We can realize $\mathbb A^n\setminus\{ 0\}$ as a bundle over $\mathbb P(d_1,\dots,d_n)$ where the fibers are precisely the $\mathbb G_m$ orbits. Therefore, a neighbourhood decomposition $D_T\times D_N$ of any point $\tilde{x}$ in a fiber over $x \in \mathbb P(d_1,\dots,d_n)$ descends to a neighbourhood decomposition $x$ given by $D_T/\mu_d\times D_n$ and vice-versa. Here $\mu_d$ is the stabilizer of $x$ for the $\mathbb G_m$ action. Moreover, this decomposition is independent of the choice of $\tilde{x}$.
% \end{remark}

% The key definitions in \cite{macpherson1998making} are those of \emph{irreducibility} and \emph{minimal irreducibility}.

% \begin{definition}[Irreducible strata]
% Let $X$ be a conically stratified variety with strata $\{S_\alpha\}_{\alpha \in I}$. We say that a strata $S_\alpha$ is \emph{reducible} if for every point $x \in X$, there exists an analytic neighbourhood decomposition $ \Spec \widehat{\mathcal O}_{X,x} \cong D_T\times D_N$ where the $D_N$ can be further factored into positive dimensional factors $D_{N,1}\times D_{N,2}$ (as cones). If this is not possible, we say that $X$ is \emph{irreducible}.
% \end{definition}

% \begin{definition}[Minimal irreducible strata]
%     We say that a strata $S_\alpha \subset X$ is minimal irreducible if it is irreducible and every strata $S_\beta < S_\alpha$ is reducible.
% \end{definition}

% The weighted blow-up stack of a conically stratified space $X$ along the closure of a strata carries a natural conical stratification as we now explain. 

% \begin{lemma}\label{lemma: blow-ups have a conical stratification}\cite[\S2.2]{macpherson1998making}
%     Consider $\mathbb A^n$ with $\mathbb G_m$ acting by 
%     \[t(x_1,\dots,x_n) \xrightarrow{} (t^{d_1}x_1,\dots,t^{d_n}x_n)\]
%     for positive integers $d_1,\dots,d_n$ whose common gcd is $1$. The weighted blow-up stack of $\mathbb A^n$ along the origin denoted by $\Bl_{0,\vec{d}}(\mathbb A^n)$ has a natural conical stratification.

%     The irreducible strata in this stratification are the irreducible strata of $\mathbb A^n\setminus\{0\}$ together with the generic stratum of $\mathbb P(d_1,\dots,d_n)$ (which is therefore a \emph{minimal} irreducible stratum).
% \end{lemma}
% \begin{proof}

% We then note that $\Bl_{0,\vec{d}}(D)$ can be realized as the total space of of the tautological bundle with the exceptional bundle embedded as the $0$ section. For every point on the exceptional bundle, we can pick a local trivialization and hence a local product structure as $D\times\mathbb A^1$ for $D \subset \mathbb P^n$. The claims then follow immediately.

% \end{proof}

% We now briefly explain the methods of \cite{macpherson1998making} and how they adapt to this setting.

% \begin{thm}\label{thm: resolution of conically stratified varieties}\cite[\S2.4]{macpherson1998making}
%     If $X$ is a conically stratified variety, then one can blow-up along the minimal irreducible strata in any order to obtain a resolution $\tilde{X}$ with an induced \emph{wonderful} conical stratification (Example \ref{example: divisor stratification}).
% \end{thm}
% \begin{proof}[Proof sketch]
%     The key difficulty in proving any type of resolution of singularities statement is to show that the situation gets \emph{better} as we perform modifications (usually blow-ups along nice centers). The proof in this case proceeds in the same way.

%     The first observation is that a conical stratification is wonderful precisely when the minimal irreducible strata are all of codimension $1$. The key property is that the closure of a minimal irreducible strata is smooth (\cite[Proposition, \S2.2]{macpherson1998making}). The proof adapts verbatim to the our setting since smoothness is reflected faithfully across analytifications. This part only uses that we have analytic neighbourhood decomposition $D_T\times D_N$ where both $D_T$ and $D_N$ are smooth and of the right dimensions.

%     This suggests the strategy to construct the resolution - we blow up the minimal irreducible strata of codimension strictly greater than $1$ in some order. As \cite[\S2.4]{macpherson1998making} shows, this always works independent of the order in which we blow up, as long as we blow up at each step along the closure of a minimal irreducible strata:
    
%     Weighted blow-ups of $X$ along a minimal irreducible strata $S$ of codimension at least $2$ preserves all the irreducible strata of $X$ except for replacing $\overline{S}$ with the exceptional divisor (which will be minimal irreducible and of codimension $1$). This shows that we can reduce the number of minimal irreducible strata of codimension at least $1$ at each step, thus eventually ending up in a situation where all the minimal irreducible strata are of codimension $1$.
% \end{proof}

% In the above setting of resolving the \emph{resultant locus}, we will consider the following conical stratification. For a tuple of polynomials $(f_1,\dots,f_r)$ corresponding to a point $F \in X = \prod_{i=1}^r\mathbb P^{n_i}$, we define its base locus $B_F$ to be the subset of $\mathbb P^1$ where at least one of the polynomials $f_i$ vanishes. We define the strata of $X$ by the size of its base locus. The generic strata corresponds to $|B_F| = \sum_i n_i$, the codimension one strata corresponds to $|B_F| = \sum_i n_i - 1$ and so on. The resulting partial order on the indexing set is $\{\sum_i n_i > \dots > 1\}$.

% \begin{lemma}\label{lemma: resultant variety is conically stratified}
%     The above stratification on $X$ is weighted conical.
% \end{lemma}
% \begin{proof}
%     Let us pick a point $F \in X$ with corresponding base locus $B_F = \{\gamma_1,\dots,\gamma_s\}$ as above. For simplicity of notation, we assume that the $\gamma_i$ are all distinct from $\infty$. Let $m_{i,j}$ be the number of times $z-\gamma_j$ divides $f_i(z)$. Moreover, for each $j$, we fix some $i_0$ with $m_{i_0,j} \geq 1$ and define
%     \[\tilde{m}_{i,j} = \begin{cases} m_{i,j} &\text{ if } i \neq i_0\\ m_{i,j}-1 &\text{ if } i = i_0 \end{cases}.\]
%     For each such element $\gamma_j$, we  define a set of variables $u_j$ and $\vec{v} = \{v_{i,k}, i \leq r, k \leq \tilde{m}_{i,j}\}$ and a map
%     \begin{align*}
%         \mathbb A^1 \times \prod_{i=1}^r\mathbb A^{\tilde{m}_{i,j}} &\xrightarrow{} X\\
%         (u_j,\vec{v}) &\xrightarrow{} ((z-\gamma_j + u_j)^{\delta_{i,i_0}}\prod_{k=1}^{\tilde{m}_{i,j}}(z-\gamma_j + v_{i,k}))_{i=1}^r.
%     \end{align*}
%     Taking the product over all $j$ gives us a map
%     \[\mathbb A^s \times \prod_{j=1}^s\prod_{i=1}^r\mathbb A^{\tilde{m}_{i,j}} \to X\]
%     that maps the origin to $F$. Passing to the completed local ring at the origin gives rise to a product decomposition
%     \[D_s(1) \times \prod_{i=1}^{r}D_{\tilde{m}_{i,j}}/S_{\tilde{m}_{i,j}} \cong \widehat{\mathcal O}_{X,x}.\]
%     The first factor corresponds to the tangent disk $D_T$ while the second factor is the normal disk $D_N$. The normal disk is stratified by the points where either two of the co-ordinates in $\vec{v}$ are equal or $v_{i_0,k} = 0$. Before passing to the quotient by the symmetric groups $S_{\tilde{m}_{i,j}}$, these strata are stable for the diagonal $\mathbb G_m$ action. Once we pass to the symmetric group quotient, the new co-ordinates correspond to the elementary symmetric polynomials $e_{i,k}(v_{i,1},\dots,v_{i,\tilde{m}_{i,j}})$ of degree $k$. These coordinates are stable under the $\mathbb G_m$ action that scales $e_{i,k}$ by $t^k$ thus showing that this decomposition is indeed as a weighted cone, thus completing the proof.
% \end{proof}
% % Consider $(\mathbb P^1)^N$ with strata corresponding to the various diagonals indexed by $\mathscr P_N$, the set of partitions of $\{1,\dots,N\}$, i.e., given a partition $\alpha \in \mathscr P_N$, we define
% % \[S_\alpha = \{ (x_1,\dots,x_n) : x_i = x_j \iff i,j \text{ are in the same part of }\alpha \}.\]
% % Next, taking $H \subset S_N$ to be of the form $S_{n_1}\times\dots \times S_{n_r}$ with $\sum n_i = N$, we consider the quotient $(\mathbb P^1)^n/H = \prod_{i=1}^r\mathbb P^{n_i}$. We will consider the factors $\mathbb P^{n_i}$ to parametrize univariate polynomials $f_i$ of degree $n_i$.

% % We consider this quotient space with strata $S_\alpha/H$ - as can be checked, we either have $S_\alpha/H = S_\beta/H$ or they are disjoint. Indeed, one can interpret the strata $S_\alpha/H$ as corresponding to tuples of polynomials $(f_1,\dots,f_r)$ where the $f_i$ have specified numbers of roots with multiplicity as well as a specified number of common roots with multiplicities for every pair $f_i,f_j$. 

% % Moreover, for such a partition $\alpha$, let $\delta_1,\dots,\delta_s$ be the parts of $\alpha$ and for $i \leq r$, define $\delta_{j,i} = \delta_j \cap [\sum_{k < i}n_k + 1,\dots,\sum_{i\leq j}n_k]$. If we let $\alpha$ correspond to a polynomial $F$ of degree $N$ with labelled roots (coming from the partition of the labels corresponding to equal roots), then the $\delta_i$ are precisely in bijection with the \emph{distinct} roots $\gamma_1,\dots,\gamma_s$ of $F$. It will in fact, suffice to consider just those $\alpha$ for which at least one of the further partitions $\delta_{j,i}$ of $\delta_j$ is non-trivial but we will ignore this distinction since considering the finer stratification will still give the desired resolution.

% % \begin{lemma}\label{lemma: our stratification is conical}
% %     The above stratification $(S_\alpha/H)_{\alpha \in \mathscr P_N}$ on $\prod_{i=1}^r\mathbb P^{n_i}$ is a conical stratification.
% % \end{lemma}
% % \begin{proof}
% % Let $x = (f_1,\dots,f_r) \in S_\alpha/H \subset \prod_{i=1}^r\mathbb P^{n_i}$ for some $\alpha \in \mathscr P_N$, where the $f_i$ are the corresponding degree $n_i$ polynomials in the variable $z$.  to be the part of $\delta_j$ contained in the interval corresponding to the polynomial $f_i$. 

% % For simplicity of notation, let us assume that the roots of all the $f_i$ are contained in $\A$. By definition of $S_\alpha$ then, we can find distinct $\gamma_1,\dots,\gamma_s \in \A$ such that $f_i = \prod_{j=1}^s(z-\gamma_j)^{|\delta_{j,i}|}$.  We use these $\gamma_j$ to construct the local conical chart around $x$:

% % All of $D_T, D_N$ and the local analytic isomorphism will have a natural product decomposition corresponding to the parts $\delta_j$ (or equivalently the roots $\gamma_j$) so we fix such a $j$ and define $D_{T,j} = \A$ with co-ordinate $u_j$ and $D_{N,j} = \prod_{i=1}^r\A^{|\delta_{i,j}|-1}$ with co-ordinates $v_{i,1},\dots,v_{i,|\delta_{i,j}|-1}$. We also define $r$ polynomials $F_{1,j},\dots,F_{r,j}$ in $z$ by
% % \[ F_{i,j}(z) = \begin{cases}
% %     (z-\gamma_j+u)\prod_{k=1}^{|\delta_{i,j}|-1}(z-\gamma_j+u+v_k) &\text{ if } |\delta_{i,j}| \geq 1\\
% %     1 &\text{ otherwise.}
% % \end{cases}\]
% % Finally, we define the local analytic isomorphism by
% % \begin{align*}
% %     D_T\times D_N \coloneqq \left(\prod_{i=1}^r D_{T,i}\right)\times\left(\prod_{i=1}^r D_{N,i}\right) &\xrightarrow{} X;(u_1,\dots,u_s)\times(v_{1,1},\dots,v_{r,|\delta_{r,s}|-1}) &\xrightarrow{} (\prod_{j=1}^sF_{i,j}(s))_{i=1}^r.
% % \end{align*}

% % \end{proof}

% By Theorem \ref{thm: resolution of conically stratified varieties} then, we can find a resolution $\pi: \tilde{X} \to X$ by taking weighted blow-ups along minimal irreducible strata such that it can be refined to be a wonderful stratification. It is easy to see that the irreducible components of the exceptional divisor are in fact sncd. Moreover, the intersections of any subset of the irreducible components are of Tate type.

% \begin{lemma}\label{lemma: our stratification is tate type}
%     The closed strata (or equivalently the locally closed strata) of $\tilde{X}$ as constructed above are of Tate type, i.e., the function $q \to |\overline{S}_\alpha(\mathbb F_q)|$ is polynomial in $q$. 
% \end{lemma}
% \begin{proof}
%     The strata of $X$ defined using the length of the base locus have a natural refinement into Tate type strata: Consider the natural quotient map
%     \[\prod_{i=1}^r(\mathbb P^1)^{n_i} \to \prod_{i=1}^r\mathbb P^{n_i}.\]
%     The refinement comes from considering the image of the stratification on $\prod_{i=1}^r(\mathbb P^1)^{n_i}$ corresponding to the locus where some of the co-ordinates are equal (i.e., the configuration space stratification). Each strata in this case is clearly seen to isomorphic to a product of projective space excluding the diagonals and this is well known to be a Tate type variety.

%     Moreover, the quotient of a Tate type variety by a finite group action is once again Tate type by the Lefschetz fixed point formula and the fact that the cohomology of a quotient are given by taking invariants of the original cohomology ring.

%     Finally, the lemma follows from the observation that if we have a stratification of a variety by Tate type strata, the blow-up of the variety along smooth closed strata also has a stratification by Tate type strata: The exceptional divisor will be isomorphic to a projective bundle over the center and we stratify this bundle by stratifying each fiber as a projective space.
% \end{proof}

% In fact, our construction of $\tilde{X}$ is equivariant for the $g$ action (where $g \in S_{N}$ is any element such that $gH = Hg$ for $N = \sum_{i}n_i$) so that we can define an automorphism $g: \tilde X \to \tilde X$ commuting with the action on $\prod_{i=1}^r\mathbb P^{n_i}$. In particular, we can define a resolution of the twist ${^g}\left(\prod_{i=1}^r\mathbb P^{n_i}\right)$ by simply twisting $\tilde X$ to get ${^g}\tilde{X} \to {^g}\left(\prod_{i=1}^r\mathbb P^{n_i}\right)$.

% \begin{lemma}\label{lemma: tate type for twists of factorizations}
%     The induced stratification on ${^g}\tilde X$ also has the property that the closed strata are of Tate type.
% \end{lemma}
% \begin{proof}
%     In fact, the twist of any Tate type variety is also of Tate type: Since the Frobenius acts by a scalar on the untwisted cohomology, the trace of the action of $\mathrm{Frob}_q$ on the twisted cohomology $H^{2i}_{\text{\'et}}({{^g}\overline{S}}_{\alpha,\mathbb F_q}, \mathbb Z_\ell)$ is given by $q^{i}\tr(g|H^{2i}_{\text{\'et}}(\overline{S}_{\alpha,\mathbb F_q}, \mathbb Z_\ell))$ (while the odd degree cohomology vanishes). This is clearly a polynomial in $q$ which proves that ${^g}\tilde{X}$ is of Tate type.\footnote{Indeed, one can check that $g$ acts by a permutation matrix so that the coefficients of the polynomial are in $\mathbb Z$. Alternatively, one could use standard functorial comparison results with singular cohomology with integer coefficients to also see that the trace of $g$ on the untwisted cohomology is integral.}
% \end{proof}

% % Finally, we note that Theorem \ref{thm: symmetry for pullback} also applies to the maps ${^g}\tilde{X} \to \mathbb P^N$ as defined above even though ${^g}\tilde{X}$ is not a scheme. This is because for a space $X$ produced through a sequence of weighted blow-ups along smooth centers from a scheme $Y/\A$, the set $Y(\A)$ can be given the structure of a $\F$-analytic manifold and this is all the input that Theorem \ref{thm: symmetry for pullback} requires. To check this, it suffices to check that if $Y$ is a smooth space with a smooth center $Z \subset Y$ so that $Y(\A)$ has the structure of a $\F$-analytic manifold, then the weighted blow-up $\Bl_{Z,\vec{d}}(Y)(\A)$ also has the structure of a $\F$-analytic manifold. 

% % This quickly reduces to showing that $\mathbb P(d_1,\dots,d_n)(\A)$ has the structure of a $\F$-analytic manifold. To verify this, we note that
% % \[\mathbb P(d_1,\dots,d_n) \cong \mathbb A^{n}\setminus\{0\}/\mathbb G_m\]
% % as a stack (where the $\mathbb G_m$ acts in a weighted way). Therefore, the $\A$ points of projective space correspond to a line bundle $\mathscr L$ over $\Spec A$ and a $\mathbb G_m$ equivariant map to the total space $\mathbb A^{n}\setminus\{0\}$. However, since line bundles over local rings are trivial, this set is simply $\mathbb A^n\setminus\{0\}(\A)/\mathbb G_m(\A) = (\A^\times)^n/(\A^\times)$ which has the structure of a $\F$-analytic manifold.

% % The techniques of $p$-adic integration we have used in the previous sections carry over without difficulty to this generality. The only non trivial point is to check that one can define a measure on $X(\A)$ by picking a local gauge form on $X$. Since this is a local problem, it suffices to assume that $X = \Bl_{0\vec{d}}(\mathbb A^n)$. In this case, on a local chart $D(x_i)$ on $\mathbb A^n$, we can pick a $\mathbb G_m$-invariant differential form of the form $dx_0\wedge dx_1 \wedge \dots dx_n/x_i^{D}$ for $D = \sum d_i$. This descends to a local gauge form on $X$ as required. The results of \S\ref{section: on the palindromicity} go through then with minor modifications.

% \begin{corollary}\label{cor: factorization densities are polynomial in q}
% The factorization densities $\rho(n,\sigma;q)$ form a rational function in $q$ and satisfy the symmetry
% \[\rho(n,\sigma;q^{-1}) = \rho(n,\sigma;q).\]
% \end{corollary}
% \begin{proof}
% As remarked above, the proof of Theorem \ref{thm: symmetry for pullback} along with Lemma \ref{lemma: tate type for twists of factorizations} is sufficient to conclude that the integrals $\eta_{f_\sigma}(L)$ are rational functions in $q$. The symmetry in $q \to q^{-1}$ is then a formal consequence of Theorem \ref{thm: main theorem for almost all primes}. This completes the proof of Equation \ref{conj: bhargava}.    
% \end{proof}



% % We will construct the $X_\sigma$ by modifying the construction of the Fulton-MacPherson compactification in \cite{fulton1994compactification}. Indeed, we will take $X_\sigma$ to be a twist-quotient of a blow-down $Y_\sigma$ of $\X{n}$:

% % For a subset $S \subset [n]$, we define $\Delta^{S} \subset (\mathbb P^1)^n$ to be the diagonal cut out by the equations $x_i = x_j$ for $i,j \in S$ and $\Bl_{\Delta^S}((\mathbb P^1)^n)$ to be the blow up along this diagonal. For each part $\pi_{i,j}(\sigma)$, we pick and fix a point $p_{i,j} \in \pi_{i,j}(\sigma)$. Subsets $S \subset P = \{p_{1,1},\dots,p_{r,f_r}\}$ will be called admissible subsets of $[n]$.

% % We define $\tilde{Y}_\sigma$ as the closure of the configuration space
% % \[F(\mathbb P^1,n) \subset (\mathbb P^1)^n \times \prod_{|S| \text{ admissible }}\Bl_{\Delta^S}((\mathbb P^1)^S)\]
% % and $Y_\sigma$ as the quotient of $\tilde{Y}_\sigma$ under the action of $H_\sigma$. This action is well defined since the admissible subsets are permuted by the action of $H_\sigma$. Finally, we define $X_\sigma$ to be the twist of $Y_\sigma$ by $g_\sigma$. We note that $\tilde{Y}_\sigma = \X{P}\times (\mathbb P^1)^{[n]-P}$ where $\X{P}$ is the Fulton-MacPherson compactification of $F(\mathbb P^1, |P|)$. 

% % By \cite[\S 3]{fulton1994compactification}, one can alternatively construct $\tilde{Y}_\sigma$ as an iterated blowup along the (proper transforms of the) diagonals $\Delta^S$ for $S$ admissible and of size at least $2$. We denote the divisors corresponding to the proper transforms of the exceptional divisors along these blowups by $D(S)$. By \cite[Theorem 3]{fulton1994compactification}, any collection $D(S_1),\dots,D(S_k)$ for admissible subsets $S_1,\dots,S_k$ intersect transversely. For any collection of admissible subsets $\mathcal S = \{S_1,\dots,S_k\}$, we define $X(\S) = \bigcap_{i=1}^k D(S_i)$. By \cite[Corollary 7a]{fulton1994compactification}, this is a Tate-symmetric variety over $\Spec \mathbb Z$.

% % Finally, we denote by $\overline{D(S)}$ the image of $D(S)$ under the quotient map $\tilde{Y}_\sigma \to Y_\sigma$ and by ${^\sigma}\overline{D(S)} \subset X_\sigma$ the twist of $\overline{D(S)}$ and similarly for the intersections $\overline{X(\mathcal S)}, {^\sigma}\overline{X(\mathcal S)}$ for a collection $\mathcal S = \{S_1,\dot,S_k\}$ of admissible subsets.

% % \begin{thm}\label{thm: nice resolutions for resultants}
% % The ramification divisor of the quotient map $Y_\sigma \to \mathbb P^n$ is precisely the union of the divisors $\overline{D(S)}$ as $S$ ranges over all admissible subsets. These divisors are sncd in the following sense: 
% % \begin{enumerate}
% %     \item The divisors $\overline{D(S)}$ are irreducible and any set of the divisors $\overline{D(S)}$ (or ${^\sigma}\overline{D(S)}$) intersect transversely in the analytic sense. 
% %     \item For any collection $\mathcal S$ of admissible subsets, the intersections $\overline{X(\mathcal S)}$ and ${^\sigma}\overline{X(\mathcal S)}$ are of Tate type over $\Spec \mathbb Z$.
% % \end{enumerate}
% % \end{thm}
% % \begin{proof}
% % The quotient map $Y_\sigma \to \mathbb P^n$ factors through the birational map $Y_\sigma \to \prod_{i=1}^r(\mathbb P^{e_i})^{f_i}$ and therefore we first analyze the quotient $\prod_{i=1}^r(\mathbb P^{e_i})^{f_i} \to \mathbb P^n$. If we interpret the projective space $\mathbb P^d$ as parametrizing degree $d$ univariate polynomials, then this quotient corresponds to the multiplication 
% % \[(Q_{1,1}(t),\dots,{Q_{r,f_r}}(t)) \to \prod_{i,j}Q_{i,j}(t)\]
% % where the $Q_{i,j}(t)$ are polynomials of degree $e_i$. We will show that the ramification locus of this map corresponds to the locus where two polynomials $Q_{i,j}(t),Q_{i',j'}(t)$ have a common root:

% % We factor the quotient map through a series of maps multiplying two polynomials at a time so that we are reduced to the case $\mathbb P^d\times\mathbb P^e \to \mathbb P^{d+e}$. The Jacobian for this map around a point $(Q_1(t),Q_2(t)) \in \mathbb P^d\times\mathbb P^e$ is easily computed to be the resultant $\operatorname{Res}(P,Q)$ and therefore the claim is proven.

% % Altogether, this shows that the ramification map for $Y_\sigma \to \mathbb P^n$ is the union of the exceptional locus for $Y_\sigma \to \prod_{i=1}^r(\mathbb P^{e_i})^{f_i}$ and the preimage of the ramification divisor for $\prod_{i=1}^r(\mathbb P^{e_i})^{f_i} \to \mathbb P^n$. This is precisely the union of the divisors $\overline{D(S)}$ as claimed. We next prove the two claims.

% % \begin{enumerate}
% %     \item Since $D(S)$ is irreducible, so is $\overline{D(S)}$. Similarly, since the $D(S)$ intersect transversely and the map $\tilde{Y}_\sigma \to Y_\sigma$ is a finite group quotient, the $\overline{D(S)}$ intersect transversely in the completion of the local ring.
% %     \item An alternate presentation of $\overline{X(\mathcal S)}$ is given as the quotient of the variety $\bigcup_{h \in H_\sigma}h(X(\mathcal S)) \subset \tilde{Y}_\sigma$ by $H_\sigma$. The explicit cellular stratification in \cite[\S 2]{fulton1994compactification} then presents the cohomology groups of $\bigcup_{h \in H_\sigma}h(X(\mathcal S))$ as a signed sum of the cohomologies of $X(\mathcal T)$ for strata $\mathcal T \leq h(\mathcal S)$ as $h$ ranges over $H_\sigma$. It is well known that the cohomology of the quotient of a variety $Z$ by a finite group $G$ is given by the $G$-invariants of the cohomology of the $Z$. The above representation of the cohomology of $\bigcup_{h \in H_\sigma}h(X(\mathcal S))$ in terms of smaller strata is compatible with the $H_\sigma$ action and upon taking $H_\sigma$ invariants, the Grothendieck-Lefshetz trace formula shows that $\overline{X(\mathcal S)}$ is of Tate type since the strata $X(\mathcal T)$ are of Tate type.
    
% %     The same claim for the twist ${^\sigma}\overline{X(\mathcal S)}$ follows from the Grothendieck Lefschetz trace formula and the fact that $g_\sigma$ acts by a permutation on the cohomology groups of $\overline{X(\mathcal S)}$ (since it does so on the cohomology groups of the strata $X(\mathcal T)$ by \cite[Corollary 7a]{fulton1994compactification}.
% % \end{enumerate}
% % \end{proof}

% % \begin{corollary}\label{cor: factorization densities are polynomial in q}[Proof of Equation \ref{eqn: symmetry for rho}]
% % The factorization densities $\rho(n,\sigma;q)$ form a polynomial function in $q$ and satisfy the symmetry
% % \[\rho(n,\sigma;q^{-1}) = \rho(n,\sigma;q).\]
% % \end{corollary}
% % \begin{proof}
% % By the above theorem, the generically finite maps $f_{\sigma}: X_\sigma \to \mathbb P^n$ satisfy the assumptions of Theorem \ref{thm: symmetry for pullback}. By the conclusion of that theorem, $\eta_{f_\sigma}(K)$ are a linear combination of the point counts ${^\sigma}\overline{X(\mathcal S)}(\mathbb F_q)$ and rational functions for any local field $K$ with residue field $\mathbb F_q$. Since the above theorem proves that ${^\sigma}\overline{X(\mathcal S)}(\mathbb F_q)$ are polynomials in $q$, this completes the proof of the theorem.

% % Now, we fix a large prime $p \gg 0$. By Theorem \ref{thm: main theorem for almost all primes}, the measure $\eta_{f_\sigma}(K)$ form a palindromic function of weight $n$. Therefore, the densities $\eta_{f_\sigma}(K)/|\mathbb P^n(\mathbb F_q)|$ form a palindromic function of weight $0$. By Equation \ref{eqn: factorization densities as integrals}, this shows that $\rho(n,\sigma;q^{-1}) = \rho(n,\sigma;q)$ as $q$ varies over the powers of $p$. If a rational function satisfies such an identity over infinitely many values, then it satisfies the identity for all values thus completing the proof.

% % \end{proof}


\printbibliography

\end{document}

\DeclareMathOperator{\Reg}{Reg}
\DeclareMathOperator{\inter}{int}
\DeclareMathOperator{\tr}{tr}
\DeclareMathOperator{\op}{OS}
\DeclareMathOperator{\Gal}{Gal}
\DeclareMathOperator{\Tr}{Tr}
\DeclareMathOperator{\Aut}{Aut}
\DeclareMathOperator{\Supp}{Supp}
\DeclareMathOperator{\Res}{Res}
\DeclareMathOperator{\denom}{denom}
\DeclareMathOperator{\num}{num}
\DeclareMathOperator{\perm}{perm}
\DeclareMathOperator{\Spec}{Spec}
\DeclareMathOperator{\can}{can}
\DeclareMathOperator{\Jac}{Jac}
\DeclareMathOperator{\Hom}{Hom}
\DeclareMathOperator{\Conf}{Conf}
\DeclareMathOperator{\Bl}{Bl}
\DeclareMathOperator{\Stab}{Stab}

\DeclareMathOperator{\Proj}{Proj}

% \DeclareMathOperator{\frob}{\sigma}

\newcommand{\frob}[1]{\sigma_#1}
\newcommand{\U}[1]{U_#1}
\newcommand{\Z}[1]{Z_#1}

%Notation for paper

\newcommand{\as}[1]{\textcolor{orange}{[[Asvin: #1]]}}
\newcommand{\jo}[1]{\textcolor{blue}{[[John: #1]]}}
\newcommand{\yif}[1]{\textcolor{teal}{[[Yifan: #1]]}}

\def\A{{\cO_{K}}}
\def\B{{\cO_{L}}}
\def\F{K}
\def\L{{L}}
\def\RF{{\mathbb F_K}}
\def\RL{{\mathbb F_{L}}}
\def\SS{S}

% spaces
\def\Gm{\mathbb{G}_m}

\def\conj{\sim}

\def\rhol{\delta}
\def\rhof{\rho_f}
\def\rhou{\rho_{f,\mathrm{ur}}}
\def\Uu{U_{f,\mathrm{ur}}}
\newcommand{\rhoid}[1]{\rho_{f,H_#1, H_{1,#1}}} 
\newcommand{\rhotype}[1]{\rho_{f,#1}}
\newcommand{\Uid}[1]{U_{f,H_#1, H_{1,#1}}}
\newcommand{\Utype}[1]{U_{f,#1}}
\newcommand{\X}[1]{\mathbb P^1[#1]}

\def\bbR{ {\mathbb R}}
\def\bbN{ {\mathbb N}}
\def\bbS{ {\mathbb S}}
\def\bbF{ {\mathbb F}}
\def\bbQ{ {\mathbb Q}}
\def\bbZ{ {\mathbb Z}}
\def\bbP{ {\mathbb P}}
\def\bbC{ {\mathbb C}}
\def\bbA{ {\mathbb A}}
\def\bbG{ {\mathbb G}}
\def\bbL{ {\mathbb L}}
\def\bbI{ {\mathbb I}}
\def\bb1{ {\mathbb 1}}
\def\bbE{ {\mathbb E}}


\def\fkA{ {\mathfrak A}}
\def\fkB{ {\mathfrak B}}
\def\fkC{ {\mathfrak C}}
\def\fkI{ {\mathfrak I}}
\def\fkJ{ {\mathfrak J}}
\def\fkm{ {\mathfrak m}}
\def\fkX{ {\mathfrak X}}

\def\cA{ {\mathcal A} }
\def\cB{ {\mathcal B} }
\def\cC{ {\mathcal C} }
\def\cD{ {\mathcal D} }
\def\cE{ {\mathcal K} }
\def\cF{ {\mathcal F} }
\def\cG{ {\mathcal G} }
\def\cH{ {\mathcal H} }
\def\cI{ {\mathcal I} }
\def\cJ{ {\mathcal J}}
\def\cK{ {\mathcal K} }
\def\cL{ {\mathcal L} }


\def\cM{{\mathcal M}}
\def\cN{ {\mathcal N} }
\def\cO{ {\mathcal O} }
\def\cP{{\mathcal P}}
\def\cQ{{\mathcal Q}}
\def\cR{{\mathcal R}}
\def\cS{{\mathcal S}}
\def\cT{{\mathcal T}}
\def\cV{{\mathcal V}}
\def\cW{{\mathcal W}}

\def\vR{{R}}
\def\vX{{X}}
\def\vY{{Y}}
\def\vf{{f}}
\def\vsigma{{\sigma}}
\def\vn{{n}}
\def\vd{{d}}

\newcounter{lst}
\newenvironment{lst}{%
\refstepcounter{lst}%
\begin{center}
\begin{minipage}{.9\textwidth}}{%
\end{minipage}%
\makebox[.1\textwidth][r]{(\thelst)}%
\end{center}}
% $ biblatex auxiliary file $
% $ biblatex bbl format version 3.2 $
% Do not modify the above lines!
%
% This is an auxiliary file used by the 'biblatex' package.
% This file may safely be deleted. It will be recreated by
% biber as required.
%
\begingroup
\makeatletter
\@ifundefined{ver@biblatex.sty}
  {\@latex@error
     {Missing 'biblatex' package}
     {The bibliography requires the 'biblatex' package.}
      \aftergroup\endinput}
  {}
\endgroup


\refsection{0}
  \datalist[entry]{nty/global//global/global}
    \entry{abramovich1996equivariant}{article}{}
      \name{author}{2}{}{%
        {{hash=97187eb4a6c3cbfa74d8150b2e476067}{%
           family={Abramovich},
           familyi={A\bibinitperiod},
           given={Dan},
           giveni={D\bibinitperiod}}}%
        {{hash=7b5e35e817f2f1ae70b06c851edbe304}{%
           family={Wang},
           familyi={W\bibinitperiod},
           given={Jianhua},
           giveni={J\bibinitperiod}}}%
      }
      \strng{namehash}{db48a866b4d13cbb5ab4ac7d8ca885ad}
      \strng{fullhash}{db48a866b4d13cbb5ab4ac7d8ca885ad}
      \strng{bibnamehash}{db48a866b4d13cbb5ab4ac7d8ca885ad}
      \strng{authorbibnamehash}{db48a866b4d13cbb5ab4ac7d8ca885ad}
      \strng{authornamehash}{db48a866b4d13cbb5ab4ac7d8ca885ad}
      \strng{authorfullhash}{db48a866b4d13cbb5ab4ac7d8ca885ad}
      \field{sortinit}{A}
      \field{sortinithash}{2f401846e2029bad6b3ecc16d50031e2}
      \field{labelnamesource}{author}
      \field{labeltitlesource}{title}
      \field{journaltitle}{arXiv preprint alg-geom/9609013}
      \field{title}{Equivariant resolution of singularities in characteristic 0}
      \field{year}{1996}
    \endentry
    \entry{batyrev1997birational}{article}{}
      \name{author}{1}{}{%
        {{hash=4a09fb297d8ab06bbe5fccccde522eb2}{%
           family={Batyrev},
           familyi={B\bibinitperiod},
           given={Victor\bibnamedelima V},
           giveni={V\bibinitperiod\bibinitdelim V\bibinitperiod}}}%
      }
      \strng{namehash}{4a09fb297d8ab06bbe5fccccde522eb2}
      \strng{fullhash}{4a09fb297d8ab06bbe5fccccde522eb2}
      \strng{bibnamehash}{4a09fb297d8ab06bbe5fccccde522eb2}
      \strng{authorbibnamehash}{4a09fb297d8ab06bbe5fccccde522eb2}
      \strng{authornamehash}{4a09fb297d8ab06bbe5fccccde522eb2}
      \strng{authorfullhash}{4a09fb297d8ab06bbe5fccccde522eb2}
      \field{sortinit}{B}
      \field{sortinithash}{d7095fff47cda75ca2589920aae98399}
      \field{labelnamesource}{author}
      \field{labeltitlesource}{title}
      \field{journaltitle}{arXiv preprint alg-geom/9710020}
      \field{title}{Birational Calabi--Yau n-folds have equal Betti numbers}
      \field{year}{1997}
    \endentry
    \entry{bergh2019functorial}{article}{}
      \name{author}{2}{}{%
        {{hash=ca6069a867fd76c258f16969a6e36a66}{%
           family={Bergh},
           familyi={B\bibinitperiod},
           given={Daniel},
           giveni={D\bibinitperiod}}}%
        {{hash=2f38a92286b394fe609fe3530974ea26}{%
           family={Rydh},
           familyi={R\bibinitperiod},
           given={David},
           giveni={D\bibinitperiod}}}%
      }
      \strng{namehash}{5bc77b3436d635212772c47e91f4d0fa}
      \strng{fullhash}{5bc77b3436d635212772c47e91f4d0fa}
      \strng{bibnamehash}{5bc77b3436d635212772c47e91f4d0fa}
      \strng{authorbibnamehash}{5bc77b3436d635212772c47e91f4d0fa}
      \strng{authornamehash}{5bc77b3436d635212772c47e91f4d0fa}
      \strng{authorfullhash}{5bc77b3436d635212772c47e91f4d0fa}
      \field{sortinit}{B}
      \field{sortinithash}{d7095fff47cda75ca2589920aae98399}
      \field{labelnamesource}{author}
      \field{labeltitlesource}{title}
      \field{journaltitle}{arXiv preprint arXiv:1905.00872}
      \field{title}{Functorial destackification and weak factorization of orbifolds}
      \field{year}{2019}
    \endentry
    \entry{bhargava2007mass}{article}{}
      \name{author}{1}{}{%
        {{hash=0d486a3693f1f42eaab7c0fbdd3091bd}{%
           family={Bhargava},
           familyi={B\bibinitperiod},
           given={Manjul},
           giveni={M\bibinitperiod}}}%
      }
      \list{publisher}{1}{%
        {OUP}%
      }
      \strng{namehash}{0d486a3693f1f42eaab7c0fbdd3091bd}
      \strng{fullhash}{0d486a3693f1f42eaab7c0fbdd3091bd}
      \strng{bibnamehash}{0d486a3693f1f42eaab7c0fbdd3091bd}
      \strng{authorbibnamehash}{0d486a3693f1f42eaab7c0fbdd3091bd}
      \strng{authornamehash}{0d486a3693f1f42eaab7c0fbdd3091bd}
      \strng{authorfullhash}{0d486a3693f1f42eaab7c0fbdd3091bd}
      \field{sortinit}{B}
      \field{sortinithash}{d7095fff47cda75ca2589920aae98399}
      \field{labelnamesource}{author}
      \field{labeltitlesource}{title}
      \field{journaltitle}{International Mathematics Research Notices}
      \field{number}{9}
      \field{title}{Mass formulae for extensions of local fields, and conjectures on the density of number field discriminants}
      \field{volume}{2007}
      \field{year}{2007}
      \field{pages}{rnm052\bibrangedash rnm052}
      \range{pages}{-1}
    \endentry
    \entry{bhargava2022density}{article}{}
      \name{author}{4}{}{%
        {{hash=0d486a3693f1f42eaab7c0fbdd3091bd}{%
           family={Bhargava},
           familyi={B\bibinitperiod},
           given={Manjul},
           giveni={M\bibinitperiod}}}%
        {{hash=8691022cccfaa2d7728e8af4971a95ca}{%
           family={Cremona},
           familyi={C\bibinitperiod},
           given={John},
           giveni={J\bibinitperiod}}}%
        {{hash=6dcba89d7f14c86211cb53397b41d02e}{%
           family={Fisher},
           familyi={F\bibinitperiod},
           given={Tom},
           giveni={T\bibinitperiod}}}%
        {{hash=e7c4a8c7581196c36f2629c59ff742b8}{%
           family={Gajović},
           familyi={G\bibinitperiod},
           given={Stevan},
           giveni={S\bibinitperiod}}}%
      }
      \list{publisher}{1}{%
        {Wiley Online Library}%
      }
      \strng{namehash}{3bc4c24b91b66de0b37d18ae62b0cdd8}
      \strng{fullhash}{fad6d135ca157419b5b3288ae0ac322a}
      \strng{bibnamehash}{3bc4c24b91b66de0b37d18ae62b0cdd8}
      \strng{authorbibnamehash}{3bc4c24b91b66de0b37d18ae62b0cdd8}
      \strng{authornamehash}{3bc4c24b91b66de0b37d18ae62b0cdd8}
      \strng{authorfullhash}{fad6d135ca157419b5b3288ae0ac322a}
      \field{sortinit}{B}
      \field{sortinithash}{d7095fff47cda75ca2589920aae98399}
      \field{labelnamesource}{author}
      \field{labeltitlesource}{title}
      \field{journaltitle}{Proceedings of the London Mathematical Society}
      \field{title}{The density of polynomials of degree n overZ\_p having exactly r roots in Q\_p}
      \field{year}{2022}
    \endentry
    \entry{bierstone1997canonical}{article}{}
      \name{author}{1}{}{%
        {{hash=998e91190d64e621edf0a3d255e45b2b}{%
           family={Bierstone},
           familyi={B\bibinitperiod},
           given={Edward},
           giveni={E\bibinitperiod}}}%
      }
      \list{publisher}{1}{%
        {Springer}%
      }
      \strng{namehash}{998e91190d64e621edf0a3d255e45b2b}
      \strng{fullhash}{998e91190d64e621edf0a3d255e45b2b}
      \strng{bibnamehash}{998e91190d64e621edf0a3d255e45b2b}
      \strng{authorbibnamehash}{998e91190d64e621edf0a3d255e45b2b}
      \strng{authornamehash}{998e91190d64e621edf0a3d255e45b2b}
      \strng{authorfullhash}{998e91190d64e621edf0a3d255e45b2b}
      \field{sortinit}{B}
      \field{sortinithash}{d7095fff47cda75ca2589920aae98399}
      \field{labelnamesource}{author}
      \field{labeltitlesource}{title}
      \field{journaltitle}{Inventiones mathematicae}
      \field{number}{2}
      \field{title}{Canonical desingularization in characteristic zero by blowing up the maximum strata of a local invariant}
      \field{volume}{128}
      \field{year}{1997}
      \field{pages}{207\bibrangedash 302}
      \range{pages}{96}
    \endentry
    \entry{bittner2004universal}{article}{}
      \name{author}{1}{}{%
        {{hash=44c20a2e69af380af71d25082a124755}{%
           family={Bittner},
           familyi={B\bibinitperiod},
           given={Franziska},
           giveni={F\bibinitperiod}}}%
      }
      \list{publisher}{1}{%
        {London Mathematical Society}%
      }
      \strng{namehash}{44c20a2e69af380af71d25082a124755}
      \strng{fullhash}{44c20a2e69af380af71d25082a124755}
      \strng{bibnamehash}{44c20a2e69af380af71d25082a124755}
      \strng{authorbibnamehash}{44c20a2e69af380af71d25082a124755}
      \strng{authornamehash}{44c20a2e69af380af71d25082a124755}
      \strng{authorfullhash}{44c20a2e69af380af71d25082a124755}
      \field{sortinit}{B}
      \field{sortinithash}{d7095fff47cda75ca2589920aae98399}
      \field{labelnamesource}{author}
      \field{labeltitlesource}{title}
      \field{journaltitle}{Compositio Mathematica}
      \field{number}{4}
      \field{title}{The universal Euler characteristic for varieties of characteristic zero}
      \field{volume}{140}
      \field{year}{2004}
      \field{pages}{1011\bibrangedash 1032}
      \range{pages}{22}
    \endentry
    \entry{bloch1932roots}{article}{}
      \name{author}{2}{}{%
        {{hash=ae3c9c0eb1ddd4c5a1a05d9e455a8311}{%
           family={Bloch},
           familyi={B\bibinitperiod},
           given={André},
           giveni={A\bibinitperiod}}}%
        {{hash=bb22658a64af6026a70b8eed70c2c795}{%
           family={Pólya},
           familyi={P\bibinitperiod},
           given={György},
           giveni={G\bibinitperiod}}}%
      }
      \list{publisher}{1}{%
        {Oxford Academic}%
      }
      \strng{namehash}{3389ab428d1ee03097ba740bac98ff46}
      \strng{fullhash}{3389ab428d1ee03097ba740bac98ff46}
      \strng{bibnamehash}{3389ab428d1ee03097ba740bac98ff46}
      \strng{authorbibnamehash}{3389ab428d1ee03097ba740bac98ff46}
      \strng{authornamehash}{3389ab428d1ee03097ba740bac98ff46}
      \strng{authorfullhash}{3389ab428d1ee03097ba740bac98ff46}
      \field{sortinit}{B}
      \field{sortinithash}{d7095fff47cda75ca2589920aae98399}
      \field{labelnamesource}{author}
      \field{labeltitlesource}{title}
      \field{journaltitle}{Proceedings of the London Mathematical Society}
      \field{number}{1}
      \field{title}{On the roots of certain algebraic equations}
      \field{volume}{2}
      \field{year}{1932}
      \field{pages}{102\bibrangedash 114}
      \range{pages}{13}
    \endentry
    \entry{buhler2006probability}{article}{}
      \name{author}{4}{}{%
        {{hash=15ab5f0d04623ddff60f4de0905655ec}{%
           family={Buhler},
           familyi={B\bibinitperiod},
           given={Joe},
           giveni={J\bibinitperiod}}}%
        {{hash=fa4ad70d293ee67ea091f10e548b781d}{%
           family={Goldstein},
           familyi={G\bibinitperiod},
           given={Daniel},
           giveni={D\bibinitperiod}}}%
        {{hash=9aa7c16f07c1338532092e23698cd8e2}{%
           family={Moews},
           familyi={M\bibinitperiod},
           given={David},
           giveni={D\bibinitperiod}}}%
        {{hash=a255beb1f7ac6a404908492e4b084ce3}{%
           family={Rosenberg},
           familyi={R\bibinitperiod},
           given={Joel},
           giveni={J\bibinitperiod}}}%
      }
      \list{publisher}{1}{%
        {Taylor \& Francis}%
      }
      \strng{namehash}{430d61a9657ae05952d6978132fb27bc}
      \strng{fullhash}{9d8b264bd2bf399e036a91c4ed357308}
      \strng{bibnamehash}{430d61a9657ae05952d6978132fb27bc}
      \strng{authorbibnamehash}{430d61a9657ae05952d6978132fb27bc}
      \strng{authornamehash}{430d61a9657ae05952d6978132fb27bc}
      \strng{authorfullhash}{9d8b264bd2bf399e036a91c4ed357308}
      \field{sortinit}{B}
      \field{sortinithash}{d7095fff47cda75ca2589920aae98399}
      \field{labelnamesource}{author}
      \field{labeltitlesource}{title}
      \field{journaltitle}{Experimental Mathematics}
      \field{number}{1}
      \field{title}{The probability that a random monic p-adic polynomial splits}
      \field{volume}{15}
      \field{year}{2006}
      \field{pages}{21\bibrangedash 32}
      \range{pages}{12}
    \endentry
    \entry{caruso2022zeroes}{inproceedings}{}
      \name{author}{1}{}{%
        {{hash=7267b9ba428088e425170cc559876a36}{%
           family={Caruso},
           familyi={C\bibinitperiod},
           given={Xavier},
           giveni={X\bibinitperiod}}}%
      }
      \list{organization}{1}{%
        {Cambridge University Press}%
      }
      \strng{namehash}{7267b9ba428088e425170cc559876a36}
      \strng{fullhash}{7267b9ba428088e425170cc559876a36}
      \strng{bibnamehash}{7267b9ba428088e425170cc559876a36}
      \strng{authorbibnamehash}{7267b9ba428088e425170cc559876a36}
      \strng{authornamehash}{7267b9ba428088e425170cc559876a36}
      \strng{authorfullhash}{7267b9ba428088e425170cc559876a36}
      \field{sortinit}{C}
      \field{sortinithash}{4d103a86280481745c9c897c925753c0}
      \field{labelnamesource}{author}
      \field{labeltitlesource}{title}
      \field{booktitle}{Forum of Mathematics, Sigma}
      \field{title}{Where are the zeroes of a random p-adic polynomial?}
      \field{volume}{10}
      \field{year}{2022}
    \endentry
    \entry{carvajal2021behavior}{article}{}
      \name{author}{2}{}{%
        {{hash=703fdc72f18ecf9d5c4c288064329137}{%
           family={Carvajal-Rojas},
           familyi={C\bibinithyphendelim R\bibinitperiod},
           given={Javier},
           giveni={J\bibinitperiod}}}%
        {{hash=bd4c36c7c494307e15372a25d86ef2e7}{%
           family={Yasuda},
           familyi={Y\bibinitperiod},
           given={Takehiko},
           giveni={T\bibinitperiod}}}%
      }
      \strng{namehash}{d4659776d9ad5ff3dd3ca6f91f1ac6a9}
      \strng{fullhash}{d4659776d9ad5ff3dd3ca6f91f1ac6a9}
      \strng{bibnamehash}{d4659776d9ad5ff3dd3ca6f91f1ac6a9}
      \strng{authorbibnamehash}{d4659776d9ad5ff3dd3ca6f91f1ac6a9}
      \strng{authornamehash}{d4659776d9ad5ff3dd3ca6f91f1ac6a9}
      \strng{authorfullhash}{d4659776d9ad5ff3dd3ca6f91f1ac6a9}
      \field{sortinit}{C}
      \field{sortinithash}{4d103a86280481745c9c897c925753c0}
      \field{labelnamesource}{author}
      \field{labeltitlesource}{title}
      \field{journaltitle}{arXiv preprint arXiv:2105.05214}
      \field{title}{On the behavior of stringy motives under Galois quasi-$\backslash$'etale covers}
      \field{year}{2021}
    \endentry
    \entry{del_corso_dvornicich_2000}{article}{}
      \name{author}{2}{}{%
        {{hash=a15564c1a50a7162366745b162021769}{%
           family={Del\bibnamedelima Corso},
           familyi={D\bibinitperiod\bibinitdelim C\bibinitperiod},
           given={Ilaria},
           giveni={I\bibinitperiod}}}%
        {{hash=da5d39f0510e1b50ea64398b0bc3f41d}{%
           family={Dvornicich},
           familyi={D\bibinitperiod},
           given={Roberto},
           giveni={R\bibinitperiod}}}%
      }
      \strng{namehash}{23b5bbfc6dcb9cb3ee1567dfb6e43cbe}
      \strng{fullhash}{23b5bbfc6dcb9cb3ee1567dfb6e43cbe}
      \strng{bibnamehash}{23b5bbfc6dcb9cb3ee1567dfb6e43cbe}
      \strng{authorbibnamehash}{23b5bbfc6dcb9cb3ee1567dfb6e43cbe}
      \strng{authornamehash}{23b5bbfc6dcb9cb3ee1567dfb6e43cbe}
      \strng{authorfullhash}{23b5bbfc6dcb9cb3ee1567dfb6e43cbe}
      \field{sortinit}{D}
      \field{sortinithash}{6f385f66841fb5e82009dc833c761848}
      \field{labelnamesource}{author}
      \field{labeltitlesource}{title}
      \field{journaltitle}{manuscripta mathematica}
      \field{number}{2}
      \field{title}{Uniformity over primes of tamely ramified splittings}
      \field{volume}{101}
      \field{year}{2000}
      \field{pages}{239\bibrangedash 266}
      \range{pages}{28}
      \verb{doi}
      \verb 10.1007/pl00005848
      \endverb
    \endentry
    \entry{dembo2002random}{article}{}
      \name{author}{4}{}{%
        {{hash=85c371b8e21ccd3f9b29c88cd37f65ef}{%
           family={Dembo},
           familyi={D\bibinitperiod},
           given={Amir},
           giveni={A\bibinitperiod}}}%
        {{hash=431e6f90528021b69e26ef43cce177fa}{%
           family={Poonen},
           familyi={P\bibinitperiod},
           given={Bjorn},
           giveni={B\bibinitperiod}}}%
        {{hash=8c41e26a80c688899996800062f2f6b2}{%
           family={Shao},
           familyi={S\bibinitperiod},
           given={Qi-Man},
           giveni={Q\bibinithyphendelim M\bibinitperiod}}}%
        {{hash=4ab05143eacf83b5d63e21a9c900ebe4}{%
           family={Zeitouni},
           familyi={Z\bibinitperiod},
           given={Ofer},
           giveni={O\bibinitperiod}}}%
      }
      \strng{namehash}{5ebb6528090fffdb2787b0aeda60632b}
      \strng{fullhash}{2ad61c5178678b4531467573216055d5}
      \strng{bibnamehash}{5ebb6528090fffdb2787b0aeda60632b}
      \strng{authorbibnamehash}{5ebb6528090fffdb2787b0aeda60632b}
      \strng{authornamehash}{5ebb6528090fffdb2787b0aeda60632b}
      \strng{authorfullhash}{2ad61c5178678b4531467573216055d5}
      \field{sortinit}{D}
      \field{sortinithash}{6f385f66841fb5e82009dc833c761848}
      \field{labelnamesource}{author}
      \field{labeltitlesource}{title}
      \field{journaltitle}{Journal of the American Mathematical Society}
      \field{number}{4}
      \field{title}{Random polynomials having few or no real zeros}
      \field{volume}{15}
      \field{year}{2002}
      \field{pages}{857\bibrangedash 892}
      \range{pages}{36}
    \endentry
    \entry{denef1992caracteristiques}{article}{}
      \name{author}{2}{}{%
        {{hash=08db5faaf14779bd520b24f44574c976}{%
           family={Denef},
           familyi={D\bibinitperiod},
           given={Jan},
           giveni={J\bibinitperiod}}}%
        {{hash=70630960cb67da39c4269310b4475b91}{%
           family={Loeser},
           familyi={L\bibinitperiod},
           given={François},
           giveni={F\bibinitperiod}}}%
      }
      \list{publisher}{1}{%
        {JSTOR}%
      }
      \strng{namehash}{a5fb0950071b953de9a80d348dfed1a8}
      \strng{fullhash}{a5fb0950071b953de9a80d348dfed1a8}
      \strng{bibnamehash}{a5fb0950071b953de9a80d348dfed1a8}
      \strng{authorbibnamehash}{a5fb0950071b953de9a80d348dfed1a8}
      \strng{authornamehash}{a5fb0950071b953de9a80d348dfed1a8}
      \strng{authorfullhash}{a5fb0950071b953de9a80d348dfed1a8}
      \field{extraname}{1}
      \field{sortinit}{D}
      \field{sortinithash}{6f385f66841fb5e82009dc833c761848}
      \field{labelnamesource}{author}
      \field{labeltitlesource}{title}
      \field{journaltitle}{Journal of the American Mathematical Society}
      \field{title}{Caractéristiques d'Euler-Poincaré, fonctions zêta locales et modifications analytiques}
      \field{year}{1992}
      \field{pages}{705\bibrangedash 720}
      \range{pages}{16}
    \endentry
    \entry{denef_loeser_2000}{article}{}
      \name{author}{2}{}{%
        {{hash=08db5faaf14779bd520b24f44574c976}{%
           family={Denef},
           familyi={D\bibinitperiod},
           given={Jan},
           giveni={J\bibinitperiod}}}%
        {{hash=70630960cb67da39c4269310b4475b91}{%
           family={Loeser},
           familyi={L\bibinitperiod},
           given={François},
           giveni={F\bibinitperiod}}}%
      }
      \strng{namehash}{a5fb0950071b953de9a80d348dfed1a8}
      \strng{fullhash}{a5fb0950071b953de9a80d348dfed1a8}
      \strng{bibnamehash}{a5fb0950071b953de9a80d348dfed1a8}
      \strng{authorbibnamehash}{a5fb0950071b953de9a80d348dfed1a8}
      \strng{authornamehash}{a5fb0950071b953de9a80d348dfed1a8}
      \strng{authorfullhash}{a5fb0950071b953de9a80d348dfed1a8}
      \field{extraname}{2}
      \field{sortinit}{D}
      \field{sortinithash}{6f385f66841fb5e82009dc833c761848}
      \field{labelnamesource}{author}
      \field{labeltitlesource}{title}
      \field{journaltitle}{Journal of the American Mathematical Society}
      \field{number}{2}
      \field{title}{Definable sets, motives and p-adic integrals}
      \field{volume}{14}
      \field{year}{2000}
      \field{pages}{429\bibrangedash 469}
      \range{pages}{41}
      \verb{doi}
      \verb 10.1090/s0894-0347-00-00360-x
      \endverb
    \endentry
    \entry{denef1991functional}{article}{}
      \name{author}{2}{}{%
        {{hash=08db5faaf14779bd520b24f44574c976}{%
           family={Denef},
           familyi={D\bibinitperiod},
           given={Jan},
           giveni={J\bibinitperiod}}}%
        {{hash=e939122b055620757fd1ddd2ead71e00}{%
           family={Meuser},
           familyi={M\bibinitperiod},
           given={Diane},
           giveni={D\bibinitperiod}}}%
      }
      \list{publisher}{1}{%
        {JSTOR}%
      }
      \strng{namehash}{3823357667750d3635e324f28fda00c8}
      \strng{fullhash}{3823357667750d3635e324f28fda00c8}
      \strng{bibnamehash}{3823357667750d3635e324f28fda00c8}
      \strng{authorbibnamehash}{3823357667750d3635e324f28fda00c8}
      \strng{authornamehash}{3823357667750d3635e324f28fda00c8}
      \strng{authorfullhash}{3823357667750d3635e324f28fda00c8}
      \field{sortinit}{D}
      \field{sortinithash}{6f385f66841fb5e82009dc833c761848}
      \field{labelnamesource}{author}
      \field{labeltitlesource}{title}
      \field{journaltitle}{American Journal of Mathematics}
      \field{number}{6}
      \field{title}{A functional equation of Igusa's local zeta function}
      \field{volume}{113}
      \field{year}{1991}
      \field{pages}{1135\bibrangedash 1152}
      \range{pages}{18}
    \endentry
    \entry{encinas1998good}{article}{}
      \name{author}{2}{}{%
        {{hash=1d5ebd9ced1a743322683e73f7cb3af0}{%
           family={Encinas},
           familyi={E\bibinitperiod},
           given={Santiago},
           giveni={S\bibinitperiod}}}%
        {{hash=d5c31f317027c6d8ad035ff933adc2c0}{%
           family={Villamayor},
           familyi={V\bibinitperiod},
           given={Orlando},
           giveni={O\bibinitperiod}}}%
      }
      \list{publisher}{1}{%
        {Springer Netherlands, Dordrecht}%
      }
      \strng{namehash}{97ed7cb70de8e0642a1ac777b35ce88a}
      \strng{fullhash}{97ed7cb70de8e0642a1ac777b35ce88a}
      \strng{bibnamehash}{97ed7cb70de8e0642a1ac777b35ce88a}
      \strng{authorbibnamehash}{97ed7cb70de8e0642a1ac777b35ce88a}
      \strng{authornamehash}{97ed7cb70de8e0642a1ac777b35ce88a}
      \strng{authorfullhash}{97ed7cb70de8e0642a1ac777b35ce88a}
      \field{sortinit}{E}
      \field{sortinithash}{8da8a182d344d5b9047633dfc0cc9131}
      \field{labelnamesource}{author}
      \field{labeltitlesource}{title}
      \field{journaltitle}{Acta mathematica}
      \field{number}{1}
      \field{title}{Good points and constructive resolution of singularities}
      \field{volume}{181}
      \field{year}{1998}
      \field{pages}{109\bibrangedash 158}
      \range{pages}{50}
    \endentry
    \entry{evans2006expected}{article}{}
      \name{author}{1}{}{%
        {{hash=197981035063baa38fb6cb3f792cc7af}{%
           family={Evans},
           familyi={E\bibinitperiod},
           given={Steven},
           giveni={S\bibinitperiod}}}%
      }
      \list{publisher}{2}{%
        {Institute of Mathematical Statistics}%
        {Bernoulli Society}%
      }
      \strng{namehash}{197981035063baa38fb6cb3f792cc7af}
      \strng{fullhash}{197981035063baa38fb6cb3f792cc7af}
      \strng{bibnamehash}{197981035063baa38fb6cb3f792cc7af}
      \strng{authorbibnamehash}{197981035063baa38fb6cb3f792cc7af}
      \strng{authornamehash}{197981035063baa38fb6cb3f792cc7af}
      \strng{authorfullhash}{197981035063baa38fb6cb3f792cc7af}
      \field{sortinit}{E}
      \field{sortinithash}{8da8a182d344d5b9047633dfc0cc9131}
      \field{labelnamesource}{author}
      \field{labeltitlesource}{title}
      \field{journaltitle}{Electronic Communications in Probability}
      \field{title}{The expected number of zeros of a random system of p -adic polynomials}
      \field{volume}{11}
      \field{year}{2006}
      \field{pages}{278\bibrangedash 290}
      \range{pages}{13}
    \endentry
    \entry{groechenig2020geometric}{article}{}
      \name{author}{3}{}{%
        {{hash=cde426fd51033558e2850565feaef2c2}{%
           family={Groechenig},
           familyi={G\bibinitperiod},
           given={Michael},
           giveni={M\bibinitperiod}}}%
        {{hash=3952ccf2b2cc632e9ee489315ee20723}{%
           family={Wyss},
           familyi={W\bibinitperiod},
           given={Dimitri},
           giveni={D\bibinitperiod}}}%
        {{hash=6f07f332d7615fb5609e51b2e42a92bb}{%
           family={Ziegler},
           familyi={Z\bibinitperiod},
           given={Paul},
           giveni={P\bibinitperiod}}}%
      }
      \strng{namehash}{4db2b60ef439945b74d511a9181ecaa0}
      \strng{fullhash}{4db2b60ef439945b74d511a9181ecaa0}
      \strng{bibnamehash}{4db2b60ef439945b74d511a9181ecaa0}
      \strng{authorbibnamehash}{4db2b60ef439945b74d511a9181ecaa0}
      \strng{authornamehash}{4db2b60ef439945b74d511a9181ecaa0}
      \strng{authorfullhash}{4db2b60ef439945b74d511a9181ecaa0}
      \field{extraname}{1}
      \field{sortinit}{G}
      \field{sortinithash}{32d67eca0634bf53703493fb1090a2e8}
      \field{labelnamesource}{author}
      \field{labeltitlesource}{title}
      \field{journaltitle}{Journal of the American Mathematical Society}
      \field{number}{3}
      \field{title}{Geometric stabilisation via p-adic integration}
      \field{volume}{33}
      \field{year}{2020}
      \field{pages}{807\bibrangedash 873}
      \range{pages}{67}
    \endentry
    \entry{groechenig2020mirror}{article}{}
      \name{author}{3}{}{%
        {{hash=cde426fd51033558e2850565feaef2c2}{%
           family={Groechenig},
           familyi={G\bibinitperiod},
           given={Michael},
           giveni={M\bibinitperiod}}}%
        {{hash=3952ccf2b2cc632e9ee489315ee20723}{%
           family={Wyss},
           familyi={W\bibinitperiod},
           given={Dimitri},
           giveni={D\bibinitperiod}}}%
        {{hash=6f07f332d7615fb5609e51b2e42a92bb}{%
           family={Ziegler},
           familyi={Z\bibinitperiod},
           given={Paul},
           giveni={P\bibinitperiod}}}%
      }
      \list{publisher}{1}{%
        {Springer}%
      }
      \strng{namehash}{4db2b60ef439945b74d511a9181ecaa0}
      \strng{fullhash}{4db2b60ef439945b74d511a9181ecaa0}
      \strng{bibnamehash}{4db2b60ef439945b74d511a9181ecaa0}
      \strng{authorbibnamehash}{4db2b60ef439945b74d511a9181ecaa0}
      \strng{authornamehash}{4db2b60ef439945b74d511a9181ecaa0}
      \strng{authorfullhash}{4db2b60ef439945b74d511a9181ecaa0}
      \field{extraname}{2}
      \field{sortinit}{G}
      \field{sortinithash}{32d67eca0634bf53703493fb1090a2e8}
      \field{labelnamesource}{author}
      \field{labeltitlesource}{title}
      \field{journaltitle}{Inventiones mathematicae}
      \field{number}{2}
      \field{title}{Mirror symmetry for moduli spaces of Higgs bundles via p-adic integration}
      \field{volume}{221}
      \field{year}{2020}
      \field{pages}{505\bibrangedash 596}
      \range{pages}{92}
    \endentry
    \entry{hausel2008mixed}{article}{}
      \name{author}{2}{}{%
        {{hash=044c11ee38ec9311f63c208b2db20b4f}{%
           family={Hausel},
           familyi={H\bibinitperiod},
           given={Tamás},
           giveni={T\bibinitperiod}}}%
        {{hash=fce30ac8f3bacc82b9c123d57c7cbdaf}{%
           family={Rodriguez-Villegas},
           familyi={R\bibinithyphendelim V\bibinitperiod},
           given={Fernando},
           giveni={F\bibinitperiod}}}%
      }
      \list{publisher}{1}{%
        {Springer}%
      }
      \strng{namehash}{598e4e65cff9c3fb2d5a45f72b03d521}
      \strng{fullhash}{598e4e65cff9c3fb2d5a45f72b03d521}
      \strng{bibnamehash}{598e4e65cff9c3fb2d5a45f72b03d521}
      \strng{authorbibnamehash}{598e4e65cff9c3fb2d5a45f72b03d521}
      \strng{authornamehash}{598e4e65cff9c3fb2d5a45f72b03d521}
      \strng{authorfullhash}{598e4e65cff9c3fb2d5a45f72b03d521}
      \field{sortinit}{H}
      \field{sortinithash}{23a3aa7c24e56cfa16945d55545109b5}
      \field{labelnamesource}{author}
      \field{labeltitlesource}{title}
      \field{journaltitle}{Inventiones mathematicae}
      \field{number}{3}
      \field{title}{Mixed Hodge polynomials of character varieties}
      \field{volume}{174}
      \field{year}{2008}
      \field{pages}{555\bibrangedash 624}
      \range{pages}{70}
    \endentry
    \entry{igusa2000introduction}{article}{}
      \name{author}{1}{}{%
        {{hash=cc81e61e91386a0b29e6ba5536e0ba5d}{%
           family={Igusa},
           familyi={I\bibinitperiod},
           given={Jun-ichi},
           giveni={J\bibinithyphendelim i\bibinitperiod}}}%
      }
      \strng{namehash}{cc81e61e91386a0b29e6ba5536e0ba5d}
      \strng{fullhash}{cc81e61e91386a0b29e6ba5536e0ba5d}
      \strng{bibnamehash}{cc81e61e91386a0b29e6ba5536e0ba5d}
      \strng{authorbibnamehash}{cc81e61e91386a0b29e6ba5536e0ba5d}
      \strng{authornamehash}{cc81e61e91386a0b29e6ba5536e0ba5d}
      \strng{authorfullhash}{cc81e61e91386a0b29e6ba5536e0ba5d}
      \field{sortinit}{I}
      \field{sortinithash}{8d291c51ee89b6cd86bf5379f0b151d8}
      \field{labelnamesource}{author}
      \field{labeltitlesource}{title}
      \field{journaltitle}{IP Studies in Advanced Mathematics. American Mathematical Society, Providence, RI}
      \field{title}{An introduction to the theory of local zeta functions, volume 14 of AMS}
      \field{volume}{6}
      \field{year}{2000}
    \endentry
    \entry{kac1943average}{article}{}
      \name{author}{1}{}{%
        {{hash=cbe7a5c62ae18623547649ceb3ca1786}{%
           family={Kac},
           familyi={K\bibinitperiod},
           given={Mark},
           giveni={M\bibinitperiod}}}%
      }
      \list{publisher}{1}{%
        {American Mathematical Society}%
      }
      \strng{namehash}{cbe7a5c62ae18623547649ceb3ca1786}
      \strng{fullhash}{cbe7a5c62ae18623547649ceb3ca1786}
      \strng{bibnamehash}{cbe7a5c62ae18623547649ceb3ca1786}
      \strng{authorbibnamehash}{cbe7a5c62ae18623547649ceb3ca1786}
      \strng{authornamehash}{cbe7a5c62ae18623547649ceb3ca1786}
      \strng{authorfullhash}{cbe7a5c62ae18623547649ceb3ca1786}
      \field{sortinit}{K}
      \field{sortinithash}{c02bf6bff1c488450c352b40f5d853ab}
      \field{labelnamesource}{author}
      \field{labeltitlesource}{title}
      \field{journaltitle}{Bulletin of the American Mathematical Society}
      \field{number}{4}
      \field{title}{On the average number of real roots of a random algebraic equation}
      \field{volume}{49}
      \field{year}{1943}
      \field{pages}{314\bibrangedash 320}
      \range{pages}{7}
    \endentry
    \entry{kulkarni2021p}{article}{}
      \name{author}{2}{}{%
        {{hash=f7338e71b69903d39dd78d39c7d25014}{%
           family={Kulkarni},
           familyi={K\bibinitperiod},
           given={Avinash},
           giveni={A\bibinitperiod}}}%
        {{hash=a3125bd69965dde6fa58fba7f4b22a3e}{%
           family={Lerario},
           familyi={L\bibinitperiod},
           given={Antonio},
           giveni={A\bibinitperiod}}}%
      }
      \list{publisher}{1}{%
        {SIAM}%
      }
      \strng{namehash}{34413c920ec6bee522283eb1dc77e63a}
      \strng{fullhash}{34413c920ec6bee522283eb1dc77e63a}
      \strng{bibnamehash}{34413c920ec6bee522283eb1dc77e63a}
      \strng{authorbibnamehash}{34413c920ec6bee522283eb1dc77e63a}
      \strng{authornamehash}{34413c920ec6bee522283eb1dc77e63a}
      \strng{authorfullhash}{34413c920ec6bee522283eb1dc77e63a}
      \field{sortinit}{K}
      \field{sortinithash}{c02bf6bff1c488450c352b40f5d853ab}
      \field{labelnamesource}{author}
      \field{labeltitlesource}{title}
      \field{journaltitle}{SIAM Journal on Applied Algebra and Geometry}
      \field{number}{1}
      \field{title}{p-adic Integral Geometry}
      \field{volume}{5}
      \field{year}{2021}
      \field{pages}{28\bibrangedash 59}
      \range{pages}{32}
    \endentry
    \entry{loeser2009seattle}{article}{}
      \name{author}{1}{}{%
        {{hash=70630960cb67da39c4269310b4475b91}{%
           family={Loeser},
           familyi={L\bibinitperiod},
           given={François},
           giveni={F\bibinitperiod}}}%
      }
      \strng{namehash}{70630960cb67da39c4269310b4475b91}
      \strng{fullhash}{70630960cb67da39c4269310b4475b91}
      \strng{bibnamehash}{70630960cb67da39c4269310b4475b91}
      \strng{authorbibnamehash}{70630960cb67da39c4269310b4475b91}
      \strng{authornamehash}{70630960cb67da39c4269310b4475b91}
      \strng{authorfullhash}{70630960cb67da39c4269310b4475b91}
      \field{sortinit}{L}
      \field{sortinithash}{7c47d417cecb1f4bd38d1825c427a61a}
      \field{labelnamesource}{author}
      \field{labeltitlesource}{title}
      \field{journaltitle}{Algebraic geometry—Seattle 2005}
      \field{title}{Seattle lectures on motivic integration}
      \field{volume}{80}
      \field{year}{2009}
      \field{pages}{745\bibrangedash 784}
      \range{pages}{40}
    \endentry
    \entry{macpherson1998making}{article}{}
      \name{author}{2}{}{%
        {{hash=1a0002b3b9ae04202872e159c0ec9712}{%
           family={MacPherson},
           familyi={M\bibinitperiod},
           given={Robert},
           giveni={R\bibinitperiod}}}%
        {{hash=3271f309c0effca6307597d1a7913f6f}{%
           family={Procesi},
           familyi={P\bibinitperiod},
           given={Claudio},
           giveni={C\bibinitperiod}}}%
      }
      \list{publisher}{1}{%
        {Springer}%
      }
      \strng{namehash}{e37d27dbd76ffef3bb93726c1a9d78c7}
      \strng{fullhash}{e37d27dbd76ffef3bb93726c1a9d78c7}
      \strng{bibnamehash}{e37d27dbd76ffef3bb93726c1a9d78c7}
      \strng{authorbibnamehash}{e37d27dbd76ffef3bb93726c1a9d78c7}
      \strng{authornamehash}{e37d27dbd76ffef3bb93726c1a9d78c7}
      \strng{authorfullhash}{e37d27dbd76ffef3bb93726c1a9d78c7}
      \field{sortinit}{M}
      \field{sortinithash}{4625c616857f13d17ce56f7d4f97d451}
      \field{labelnamesource}{author}
      \field{labeltitlesource}{title}
      \field{journaltitle}{Selecta Mathematica}
      \field{title}{Making conical compactifications wonderful}
      \field{volume}{4}
      \field{year}{1998}
      \field{pages}{125\bibrangedash 139}
      \range{pages}{15}
    \endentry
    \entry{nguyen2021random}{article}{}
      \name{author}{2}{}{%
        {{hash=f6335af274eca8e4ac4f372022111414}{%
           family={Nguyen},
           familyi={N\bibinitperiod},
           given={Oanh},
           giveni={O\bibinitperiod}}}%
        {{hash=8ad506863f61cad43fe4db578d9cdf42}{%
           family={Vu},
           familyi={V\bibinitperiod},
           given={Van},
           giveni={V\bibinitperiod}}}%
      }
      \list{publisher}{1}{%
        {Duke University Press}%
      }
      \strng{namehash}{4359ab8b31ac94d10908a6e5abe29367}
      \strng{fullhash}{4359ab8b31ac94d10908a6e5abe29367}
      \strng{bibnamehash}{4359ab8b31ac94d10908a6e5abe29367}
      \strng{authorbibnamehash}{4359ab8b31ac94d10908a6e5abe29367}
      \strng{authornamehash}{4359ab8b31ac94d10908a6e5abe29367}
      \strng{authorfullhash}{4359ab8b31ac94d10908a6e5abe29367}
      \field{sortinit}{N}
      \field{sortinithash}{22369a73d5f88983a108b63f07f37084}
      \field{labelnamesource}{author}
      \field{labeltitlesource}{title}
      \field{journaltitle}{Duke Mathematical Journal}
      \field{number}{17}
      \field{title}{Random polynomials: central limit theorems for the real roots}
      \field{volume}{170}
      \field{year}{2021}
      \field{pages}{3745\bibrangedash 3813}
      \range{pages}{69}
    \endentry
    \entry{poonen2017rational}{book}{}
      \name{author}{1}{}{%
        {{hash=431e6f90528021b69e26ef43cce177fa}{%
           family={Poonen},
           familyi={P\bibinitperiod},
           given={Bjorn},
           giveni={B\bibinitperiod}}}%
      }
      \list{publisher}{1}{%
        {American Mathematical Soc.}%
      }
      \strng{namehash}{431e6f90528021b69e26ef43cce177fa}
      \strng{fullhash}{431e6f90528021b69e26ef43cce177fa}
      \strng{bibnamehash}{431e6f90528021b69e26ef43cce177fa}
      \strng{authorbibnamehash}{431e6f90528021b69e26ef43cce177fa}
      \strng{authornamehash}{431e6f90528021b69e26ef43cce177fa}
      \strng{authorfullhash}{431e6f90528021b69e26ef43cce177fa}
      \field{sortinit}{P}
      \field{sortinithash}{ff3bcf24f47321b42cb156c2cc8a8422}
      \field{labelnamesource}{author}
      \field{labeltitlesource}{title}
      \field{title}{Rational points on varieties}
      \field{volume}{186}
      \field{year}{2017}
    \endentry
    \entry{reid2002correspondance}{article}{}
      \name{author}{1}{}{%
        {{hash=b36ed2e74acb2c122428f1cb5b896108}{%
           family={Reid},
           familyi={R\bibinitperiod},
           given={Miles},
           giveni={M\bibinitperiod}}}%
      }
      \list{publisher}{1}{%
        {CENTRE NATIONAL DE LA RECHERCHE SCIENTIFIQUE}%
      }
      \strng{namehash}{b36ed2e74acb2c122428f1cb5b896108}
      \strng{fullhash}{b36ed2e74acb2c122428f1cb5b896108}
      \strng{bibnamehash}{b36ed2e74acb2c122428f1cb5b896108}
      \strng{authorbibnamehash}{b36ed2e74acb2c122428f1cb5b896108}
      \strng{authornamehash}{b36ed2e74acb2c122428f1cb5b896108}
      \strng{authorfullhash}{b36ed2e74acb2c122428f1cb5b896108}
      \field{sortinit}{R}
      \field{sortinithash}{5e1c39a9d46ffb6bebd8f801023a9486}
      \field{labelnamesource}{author}
      \field{labeltitlesource}{title}
      \field{journaltitle}{Asterisque-Societe Mathematique de France}
      \field{title}{La correspondance de McKay}
      \field{volume}{276}
      \field{year}{2002}
      \field{pages}{53\bibrangedash 72}
      \range{pages}{20}
    \endentry
    \entry{serre1978formule}{article}{}
      \name{author}{1}{}{%
        {{hash=523a41970d92b80420ca6d5307619da7}{%
           family={Serre},
           familyi={S\bibinitperiod},
           given={Jean-Pierre},
           giveni={J\bibinithyphendelim P\bibinitperiod}}}%
      }
      \strng{namehash}{523a41970d92b80420ca6d5307619da7}
      \strng{fullhash}{523a41970d92b80420ca6d5307619da7}
      \strng{bibnamehash}{523a41970d92b80420ca6d5307619da7}
      \strng{authorbibnamehash}{523a41970d92b80420ca6d5307619da7}
      \strng{authornamehash}{523a41970d92b80420ca6d5307619da7}
      \strng{authorfullhash}{523a41970d92b80420ca6d5307619da7}
      \field{sortinit}{S}
      \field{sortinithash}{b164b07b29984b41daf1e85279fbc5ab}
      \field{labelnamesource}{author}
      \field{labeltitlesource}{title}
      \field{journaltitle}{CR Acad. Sc. Paris A}
      \field{title}{Une formule de masse pour les extensions totalement ramifiees de degres dinne d'un corps local.}
      \field{volume}{286}
      \field{year}{1978}
      \field{pages}{1031\bibrangedash 1036}
      \range{pages}{6}
    \endentry
    \entry{roy}{misc}{}
      \name{author}{1}{}{%
        {{hash=67126962386598937414e6b347f00547}{%
           family={Shmueli},
           familyi={S\bibinitperiod},
           given={Roy},
           giveni={R\bibinitperiod}}}%
      }
      \list{publisher}{1}{%
        {arXiv}%
      }
      \strng{namehash}{67126962386598937414e6b347f00547}
      \strng{fullhash}{67126962386598937414e6b347f00547}
      \strng{bibnamehash}{67126962386598937414e6b347f00547}
      \strng{authorbibnamehash}{67126962386598937414e6b347f00547}
      \strng{authornamehash}{67126962386598937414e6b347f00547}
      \strng{authorfullhash}{67126962386598937414e6b347f00547}
      \field{sortinit}{S}
      \field{sortinithash}{b164b07b29984b41daf1e85279fbc5ab}
      \field{labelnamesource}{author}
      \field{labeltitlesource}{title}
      \field{title}{The probability that a p-adic random étale algebra is an unramified field}
      \field{year}{2022}
      \verb{doi}
      \verb 10.48550/arXiv.2211.12995
      \endverb
      \verb{urlraw}
      \verb https://arxiv.org/abs/2211.12995
      \endverb
      \verb{url}
      \verb https://arxiv.org/abs/2211.12995
      \endverb
      \keyw{Number Theory (math.NT),FOS: Mathematics,FOS: Mathematics}
    \endentry
    \entry{yasuda2017wild}{article}{}
      \name{author}{1}{}{%
        {{hash=bd4c36c7c494307e15372a25d86ef2e7}{%
           family={Yasuda},
           familyi={Y\bibinitperiod},
           given={Takehiko},
           giveni={T\bibinitperiod}}}%
      }
      \strng{namehash}{bd4c36c7c494307e15372a25d86ef2e7}
      \strng{fullhash}{bd4c36c7c494307e15372a25d86ef2e7}
      \strng{bibnamehash}{bd4c36c7c494307e15372a25d86ef2e7}
      \strng{authorbibnamehash}{bd4c36c7c494307e15372a25d86ef2e7}
      \strng{authornamehash}{bd4c36c7c494307e15372a25d86ef2e7}
      \strng{authorfullhash}{bd4c36c7c494307e15372a25d86ef2e7}
      \field{sortinit}{Y}
      \field{sortinithash}{fd67ad5a9ef0f7456bdd9aab10fe1495}
      \field{labelnamesource}{author}
      \field{labeltitlesource}{title}
      \field{journaltitle}{Journal of the European Mathematical Society}
      \field{number}{12}
      \field{title}{The wild McKay correspondence and p-adic measures}
      \field{volume}{19}
      \field{year}{2017}
      \field{pages}{3709\bibrangedash 3743}
      \range{pages}{35}
    \endentry
    \entry{john}{misc}{}
      \name{author}{1}{}{%
        {{hash=a83856f362ed9a405c7a55282ac0ae31}{%
           family={Yin},
           familyi={Y\bibinitperiod},
           given={John},
           giveni={J\bibinitperiod}}}%
      }
      \list{publisher}{1}{%
        {arXiv}%
      }
      \strng{namehash}{a83856f362ed9a405c7a55282ac0ae31}
      \strng{fullhash}{a83856f362ed9a405c7a55282ac0ae31}
      \strng{bibnamehash}{a83856f362ed9a405c7a55282ac0ae31}
      \strng{authorbibnamehash}{a83856f362ed9a405c7a55282ac0ae31}
      \strng{authornamehash}{a83856f362ed9a405c7a55282ac0ae31}
      \strng{authorfullhash}{a83856f362ed9a405c7a55282ac0ae31}
      \field{sortinit}{Y}
      \field{sortinithash}{fd67ad5a9ef0f7456bdd9aab10fe1495}
      \field{labelnamesource}{author}
      \field{labeltitlesource}{title}
      \field{title}{Density of $p$-adic polynomials generating extensions with fixed splitting type}
      \field{year}{2022}
      \verb{doi}
      \verb 10.48550/ARXIV.2211.10425
      \endverb
      \verb{urlraw}
      \verb https://arxiv.org/abs/2211.10425
      \endverb
      \verb{url}
      \verb https://arxiv.org/abs/2211.10425
      \endverb
      \keyw{Number Theory (math.NT),FOS: Mathematics,FOS: Mathematics}
    \endentry
  \enddatalist
\endrefsection
\endinput

